\documentclass[12pt,a4paper,oneside]{report}

% ==============================================================================
% CẤU HÌNH GÓI LỆNH VÀ TƯƠNG THÍCH TRÌNH BIÊN DỊCH (PDFLATEX / XELATEX)
% ==============================================================================
\usepackage{iftex}
\ifXeTeX
  \usepackage{fontspec}
  \setmainfont{Cambria}[
    BoldFont = Cambria Bold,
    ItalicFont = Cambria Italic,
    BoldItalicFont = Cambria Bold Italic
  ]
\else
  \usepackage[utf8]{inputenc}
  \usepackage[vietnamese]{babel}
  \usepackage[T5]{fontenc}
  \usepackage{mathptmx} % Font tương thích chuẩn pdfLaTeX
\fi

% Cấu hình kiểm soát độ co giãn văn bản, triệt tiêu tràn lề phải
\usepackage{microtype}
\emergencystretch=3em
\tolerance=2000
\hfuzz=1pt
\hbadness=2000

% Căn lề trang in chuẩn Luận văn / Báo cáo NCKH Học viện Công nghệ Bưu chính Viễn thông
% Lề trái: 3.5 cm (đóng gáy), Lề phải: 2.0 cm, Lề trên: 2.5 cm, Lề dưới: 2.5 cm
\usepackage[left=3.5cm,right=2.0cm,top=2.5cm,bottom=2.5cm]{geometry}
\setlength{\headheight}{25pt}
\usepackage{setspace}
\onehalfspacing % Giãn dòng 1.5 lines tiêu chuẩn

% Các gói hỗ trợ toán học và biểu thức đại số
\usepackage{amsmath,amssymb,amsfonts,amsthm}
\usepackage{bm}
\usepackage{mathtools}

% Gói hỗ trợ hình ảnh, bảng biểu và đồ thị
\usepackage{graphicx}
\usepackage{pdfpages}
\graphicspath{{paper/fig/}{paper/}{fig/}{./}}
\usepackage{booktabs}
\usepackage{tabularx}
\usepackage{multirow}
\usepackage{multicol}
\usepackage{array}
\usepackage{float}
\usepackage{placeins}
\usepackage{caption}
\usepackage{subcaption}
\usepackage{tikz}
\usetikzlibrary{calc,shapes.geometric,arrows,arrows.meta,positioning,fit}

% Định dạng tiêu đề chương, mục chuẩn học thuật hiện đại
\usepackage{titlesec}
\titleformat{\chapter}[display]
  {\normalfont\bfseries\centering}
  {\fontsize{13}{15}\selectfont\color{black}\MakeUppercase{Chương \thechapter}}
  {0.6ex}
  {\fontsize{16}{19}\selectfont\color{black}\MakeUppercase}
  [\vspace{0.8ex}\color{black}\rule{0.5\textwidth}{0.8pt}]
\titlespacing*{\chapter}{0pt}{10pt}{25pt}

\titleformat{\section}
  {\normalfont\bfseries\fontsize{13}{15}\selectfont\color{black}}
  {\thesection}{1em}{}
\titlespacing*{\section}{0pt}{14pt}{6pt}

\titleformat{\subsection}
  {\normalfont\bfseries\fontsize{12}{14}\selectfont\color{black}}
  {\thesubsection}{1em}{}
\titlespacing*{\subsection}{0pt}{10pt}{4pt}

\titleformat{\subsubsection}
  {\normalfont\itshape\bfseries\fontsize{12}{14}\selectfont\color{black}}
  {\thesubsubsection}{1em}{}
\titlespacing*{\subsubsection}{0pt}{8pt}{3pt}

% Gói hỗ trợ Header và Footer trang trí chuyên nghiệp, sang trọng
% Cấu hình Header và Footer: Bỏ header trên đầu trang, chỉ giữ số trang ở chân trang
\usepackage{fancyhdr}
\pagestyle{fancy}
\fancyhf{}
\fancyfoot[C]{\fontsize{10}{12}\selectfont\thepage}
\renewcommand{\headrulewidth}{0pt}
\renewcommand{\footrulewidth}{0pt}

% Cấu hình liên kết Hyperref
\usepackage{hyperref}
\hypersetup{
    colorlinks=true,
    linkcolor=black,
    citecolor=black,
    urlcolor=black,
    pdftitle={Báo cáo Đề tài NCKH Cấp Học viện - PTIT},
    pdfauthor={Lê Việt Anh},
    pdfsubject={Nghiên cứu dự báo giao thông thông minh}
}

% Định dạng danh sách liệt kê
\usepackage{enumitem}
\setlist[itemize]{noitemsep, topsep=3pt, leftmargin=1.2cm}
\setlist[enumerate]{noitemsep, topsep=3pt, leftmargin=1.2cm}

% Định nghĩa môi trường thuật toán sạch đẹp
\newcounter{algo}
\newenvironment{algorithm_box}[2]{%
  \refstepcounter{algo}%
  \begin{table}[H]
  \centering
  \begin{tabularx}{\textwidth}{|X|}
  \hline
  \textbf{Thuật toán \thealgo: #1} \\
  \hline
  \small #2 \\
  \hline
  \end{tabularx}
  \end{table}
}{}

% ==============================================================================
% BẮT ĐẦU TÀI LIỆU
% ==============================================================================

% Định nghĩa tường minh các tiêu đề danh mục tiếng Việt chuẩn mực
\renewcommand{\contentsname}{MỤC LỤC}
\renewcommand{\listtablename}{DANH MỤC BẢNG BIỂU}
\renewcommand{\listfigurename}{DANH MỤC HÌNH VẼ}
\renewcommand{\bibname}{TÀI LIỆU THAM KHẢO}
\renewcommand{\figurename}{Hình}
\renewcommand{\tablename}{Bảng}
\renewcommand{\chaptername}{Chương}

\begin{document}
% ------------------------------------------------------------------------------
% TRANG BÌA CHÍNH THỨC (XUẤT TỪ BIA.DOCX)
% ------------------------------------------------------------------------------
\includepdf[pages=1]{bia.pdf}


% ------------------------------------------------------------------------------
% LỜI CẢM ƠN
% ------------------------------------------------------------------------------
\pagenumbering{roman}
\setcounter{page}{1}
\chapter*{LỜI CẢM ƠN}
\addcontentsline{toc}{chapter}{Lời cảm ơn}

Trước hết, tôi xin gửi lời cảm ơn tới các thầy cô giảng viên của Học viện Công nghệ Bưu chính Viễn thông, đặc biệt là các thầy cô thuộc Khoa Công nghệ Thông tin I. Trong suốt quá trình học tập, tôi đã được trang bị những kiến thức nền tảng và chuyên môn cần thiết, đồng thời học hỏi được tinh thần làm việc nghiêm túc và trách nhiệm từ thầy cô. Đây là những hành trang quan trọng giúp tôi hoàn thành đề tài nghiên cứu khoa học này.

Tôi xin chân thành cảm ơn giảng viên hướng dẫn, thầy PGS.TS Nguyễn Trọng Khánh, người đã hỗ trợ và góp ý cho tôi trong suốt quá trình thực hiện đề tài. Những nhận xét và hướng dẫn của thầy đã giúp tôi định hướng và hoàn thiện công trình nghiên cứu tốt hơn.

Bên cạnh đó, tôi cũng xin cảm ơn bạn bè và các bạn cùng lớp đã hỗ trợ chia sẻ tài liệu, trao đổi kiến thức và đóng góp ý kiến trong quá trình học tập và thực hiện báo cáo.

Do kiến thức và kinh nghiệm còn hạn chế, bài báo cáo khó tránh khỏi những thiếu sót. Tôi rất mong nhận được những ý kiến đóng góp từ thầy cô và các bạn để có thể tiếp tục hoàn thiện và nâng cao kiến thức của bản thân.

Cuối cùng, tôi xin chân thành cảm ơn tất cả những người đã hỗ trợ và đồng hành cùng tôi trong suốt quá trình thực hiện đề tài nghiên cứu khoa học này.

\vspace{1.5cm}
\begin{flushright}
    \textit{Hà Nội, Tháng 10 năm 2026} \\
    \textbf{SINH VIÊN THỰC HIỆN} \\
    \vspace{1.5cm}
    \textbf{Lê Việt Anh}
\end{flushright}


\chapter*{TÓM TẮT ĐỀ TÀI}
\addcontentsline{toc}{chapter}{Tóm tắt đề tài}

Trong bối cảnh đô thị hóa nhanh chóng tại các quốc gia Đông Nam Á, các mạng lưới camera giám sát công cộng (CCTV) cung cấp một nguồn cấp dữ liệu thị giác phong phú, bao phủ rộng khắp với chi phí vận hành thấp hơn đáng kể so với cảm biến vòng từ truyền thống. Tuy nhiên, hiện nay có rất ít bộ dữ liệu chuẩn mực và công trình nghiên cứu đánh giá đồng thời cả hai khía cạnh: (1) nhận dạng lưu lượng phương tiện ở mức độ camera đơn lẻ và (2) dự báo dòng giao thông không -- thời gian trên quy mô toàn thành phố dưới điều kiện giao thông hỗn hợp với xe máy chiếm ưu thế tuyệt đối ($>80\%$). Đóng góp cốt lõi của đề tài này là xây dựng khung mô hình hai giai đoạn hoàn chỉnh dựa hoàn toàn trên dữ liệu camera và thực hiện phân tích có hệ thống về hiện tượng \textbf{lan truyền sai số nhận dạng thị giác (Error Propagation)} từ camera vào các mô hình mạng nơ-ron đồ thị dự báo đa bước dưới các mức nhiễu cảm biến mô phỏng.

Về mặt kỹ thuật, nghiên cứu đề xuất một khung làm việc hai giai đoạn (Two-Stage Pipeline) và kiểm chứng trên \textbf{IC4SD-Traffic-HCM} -- một bộ dữ liệu chuẩn quy mô thành phố bao gồm \textbf{608 nút camera} tại TP. Hồ Chí Minh được xây dựng và công bố:
\begin{enumerate}
    \item \textbf{Giai đoạn 1 (Vision Perception):} Sử dụng mạng nơ-ron tích chập EfficientNet-B5 kết hợp đầu hồi quy MLP đa tầng thống nhất để dự đoán trực tiếp số lượng phương tiện theo phân lớp (xe ô tô con và xe máy) từ ảnh tĩnh. Cách tiếp cận này loại bỏ hoàn toàn tầng hậu xử lý triệt tiêu cực đại cục bộ (NMS) vốn có độ phức tạp bậc hai $\mathcal{O}(M^2)$ và giải quyết triệt để hiện tượng che khuất tương hỗ nghiêm trọng giữa các cụm xe máy.
    \item \textbf{Giai đoạn 2 (Graph Forecasting):} Đánh giá 7 kiến trúc đồ thị và Transformer tiên tiến nhất (STGCN, GraphWaveNet, ASTGCN, MegaCRN, DSTAGNN, STAEformer, iTransformer) cùng với mô hình đề xuất \textbf{TA-STGCN} (Temporal Attention-Guided Spatio-Temporal Graph Convolutional Network) -- một biến thể cải tiến bổ sung cơ chế Chú ý tự thân theo thời gian đa đầu (Multi-Head Temporal Self-Attention) xen kẽ giữa các khối tích chập đồ thị phổ Chebyshev.
\end{enumerate}

Kết quả thực nghiệm trên 5 seed ngẫu nhiên độc lập khẳng định rằng trên dữ liệu sạch, các mô hình đồ thị đạt độ chính xác gần như tương đương (TA-STGCN đạt sai số MAE $= 3.2401 \pm 0.0108$, chỉ nhỉnh hơn STGCN gốc $\sim 0.9\%$, quy mô ảnh hưởng Cohen's $d = -0.032$). Tuy nhiên, dưới tác động của nhiễu nhận dạng cảm biến mô phỏng thực tế (phân phối Gamma bất đối xứng có phương sai dị biệt tỷ lệ với lưu lượng), khả năng chống chịu nhiễu phân hóa rõ rệt: TA-STGCN vượt trội áp đảo các mô hình không có cơ chế chú ý thời gian (Cohen's $d$ lên tới $+5.84$ so với GraphWaveNet ở mức nhiễu cao nhất), trong khi duy trì sai số tương đương các Transformer khổng lồ nhưng với chi phí tham số chỉ 453K tham số và tiết kiệm $25.2\%$ bộ nhớ GPU đỉnh so với STAEformer. Nghiên cứu mở ra triển vọng triển khai kiến trúc Biên -- Đám mây (Edge--Cloud) truyền vector số đếm bảo mật và tiết kiệm băng thông gấp 62.500 lần cho đô thị thông minh.

\vspace{0.5cm}
\noindent \textbf{Từ khóa:} Giao thông thông minh (ITS), Dự báo giao thông đô thị, IC4SD-Traffic-HCM, Hồi quy đếm xe trực tiếp, TA-STGCN, Mạng nơ-ron đồ thị không -- thời gian, Lan truyền sai số, Độ bền vững với nhiễu.

% ------------------------------------------------------------------------------
% MỤC LỤC, DANH MỤC BẢNG, DANH MỤC HÌNH VẼ, DANH MỤC VIẾT TẮT
% ------------------------------------------------------------------------------
\newpage
\tableofcontents

\newpage
\listoftables
\addcontentsline{toc}{chapter}{Danh mục bảng biểu}

\newpage
\listoffigures
\addcontentsline{toc}{chapter}{Danh mục hình vẽ}

\newpage
\chapter*{DANH MỤC TỪ VIẾT TẮT}
\addcontentsline{toc}{chapter}{Danh mục từ viết tắt}

\begin{table}[H]
\centering
\begin{tabularx}{\textwidth}{l p{4.5cm} X}
\toprule
\textbf{Viết tắt} & \textbf{Thuật ngữ tiếng Anh} & \textbf{Ý nghĩa học thuật} \\
\midrule
ITS & Intelligent Transportation Systems & Hệ thống Giao thông Thông minh \\
CCTV & Closed-Circuit Television & Hệ thống Camera giám sát công cộng \\
CNN & Convolutional Neural Network & Mạng nơ-ron tích chập \\
GCN & Graph Convolutional Network & Mạng tích chập đồ thị \\
STGNN & Spatio-Temporal Graph Neural Network & Mạng nơ-ron đồ thị không -- thời gian \\
STGCN & Spatio-Temporal Graph ConvNet & Mạng tích chập đồ thị không -- thời gian \\
TA-STGCN & Temporal Attention-Guided STGCN & Mạng STGCN điều hướng bởi chú ý thời gian \\
TCN & Temporal Convolutional Network & Mạng tích chập theo thời gian \\
GLU & Gated Linear Unit & Đơn vị tuyến tính có cổng \\
RBF & Radial Basis Function & Hàm cơ sở xuyên tâm \\
OSM & OpenStreetMap & Nền tảng bản đồ đường phố mở \\
MAE & Mean Absolute Error & Sai số tuyệt đối trung bình \\
MAPE & Mean Absolute Percentage Error & Sai số phần trăm tuyệt đối trung bình \\
RMSE & Root Mean Squared Error & Căn bậc hai sai số toàn phương trung bình \\
CI & Confidence Interval & Khoảng tin cậy thống kê \\
CV & Coefficient of Variation & Hệ số biến thiên phương sai \\
FLOPs & Floating Point Operations & Số phép toán dấu phẩy động \\
NMS & Non-Maximum Suppression & Thuật toán triệt tiêu cực đại cục bộ \\
Grad-CAM & Gradient-weighted Class Activation & Trực quan hóa bản đồ kích hoạt đặc trưng \\
SOTA & State-of-the-Art & Trình độ tiên tiến nhất hiện nay \\
PCU & Passenger Car Unit & Hệ số xe con quy đổi \\
VRAM & Video Random Access Memory & Bộ nhớ đồ họa chuyên dụng \\
\bottomrule
\end{tabularx}
\end{table}

% ------------------------------------------------------------------------------
% NỘI DUNG CHÍNH (MAIN CHAPTERS QUA TỪNG TỆP CHUYÊN SÂU)
% ------------------------------------------------------------------------------
\newpage
\pagenumbering{arabic}
\setcounter{page}{1}

% Phần mở đầu
\chapter*{PHẦN MỞ ĐẦU}
\addcontentsline{toc}{chapter}{Phần Mở đầu}

\section*{1. Bối cảnh và Tính cấp thiết của Đề tài}
Trong kỷ nguyên hiện đại, tiến trình đô thị hóa diễn ra với tốc độ vũ bão tại các quốc gia đang phát triển, đặc biệt là khu vực Đông Nam Á, đã kéo theo sự bùng nổ chưa từng có của các phương tiện giao thông cơ giới cá nhân. Tình trạng quá tải hạ tầng đường bộ, ùn tắc giao thông kéo dài và ô nhiễm môi trường đã trở thành bài toán hóc búa, đe dọa trực tiếp đến sự phát triển kinh tế bền vững và chất lượng cuộc sống của hàng chục triệu cư dân đô thị. Tại các siêu đô thị như Thành phố Hồ Chí Minh và Thủ đô Hà Nội, thiệt hại kinh tế hàng năm do ùn tắc giao thông ước tính lên tới hàng tỷ USD, chưa kể đến hàng triệu giờ lao động bị lãng phí và những hệ lụy nặng nề về sức khỏe cộng đồng do khí thải phương tiện gây ra.

Để giải quyết vấn đề nan giải này, việc phát triển và ứng dụng các Hệ thống Giao thông Thông minh (Intelligent Transportation Systems -- ITS) được xem là giải pháp mang tính chiến lược và đột phá. Khác với cách tiếp cận truyền thống vốn tập trung mở rộng hạ tầng cầu đường một cách tốn kém và bị giới hạn bởi quỹ đất đô thị, ITS tận dụng sức mạnh của công nghệ thông tin, trí tuệ nhân tạo (AI) và mạng truyền thông không dây để tối ưu hóa năng lực điều hành và phân phối dòng di chuyển. Một hệ thống ITS hoàn chỉnh đòi hỏi sự phối hợp nhịp nhàng giữa hai năng lực công nghệ cốt lõi:
\begin{enumerate}
    \item \textbf{Năng lực Giám sát và Thu nhận Trạng thái Giao thông Tức thời (Traffic Perception):} Định lượng chính xác mật độ, số lượng và phân loại các phương tiện giao thông đang hiện diện trên từng nút giao cắt hoặc hành lang tuyến đường tại thời điểm thực.
    \item \textbf{Năng lực Dự báo Dòng Giao thông Đa bước trên Toàn Mạng lưới (Network-Wide Multi-Horizon Traffic Forecasting):} Dự đoán chính xác xu hướng biến động của lưu lượng phương tiện trên toàn bộ đồ thị mạng lưới đường bộ thành phố trong khoảng thời gian tương lai từ 15 đến 30 phút ($H = 6$ bước thời gian, với chu kỳ lấy mẫu $\Delta t = 5$ phút).
\end{enumerate}

Tầm nhìn dự báo trong khoảng 15--30 phút (Lookahead Forecasting) giữ vai trò "sống còn" đối với các trung tâm chỉ huy điều hành giao thông đô thị:
\begin{itemize}
    \item \textbf{Điều khiển đèn tín hiệu thích ứng đón đầu (Proactive Adaptive Signal Control):} Trong lý thuyết điều khiển giao thông, chu kỳ đèn tín hiệu tại các ngã tư trọng điểm thường dao động từ 90 đến 180 giây. Việc thay đổi pha đèn cần từ 3 đến 5 chu kỳ liên tiếp để dòng xe tự điều chỉnh và tái lập trạng thái cân bằng. Nếu hệ thống chỉ dựa vào số liệu quan trắc tức thời tại thời điểm hiện tại ($t$), thuật toán điều khiển sẽ rơi vào trạng thái trễ pha (Phase-Lag), khiến hàng chờ xe bị dồn ứ trước khi pha đèn kịp kéo dài. Dự báo chính xác trước 15--30 phút cho phép hệ thống tính toán và kích hoạt kế hoạch phân bổ thời lượng đèn xanh đón đầu trước khi đỉnh tắc nghẽn thực sự ập đến.
    \item \textbf{Kiểm soát sự lan truyền sóng xung kích tắc nghẽn (Shockwave Propagation Control):} Dựa trên Mô hình Truyền sóng Thủy động lực học (Cell Transmission Model -- CTM), khi một nút giao bị tắc nghẽn cục bộ, sóng xung kích tắc nghẽn sẽ lan truyền ngược chiều dòng xe với vận tốc khoảng 15--25 km/h, nhanh chóng làm tê liệt các tuyến đường nhánh thượng lưu. Dự báo trước 30 phút cung cấp "khoảng thời gian vàng" để trung tâm điều hành kịp thời điều chỉnh các biển báo điện tử VMS (Variable Message Signs), cảnh báo người tham gia giao thông trên các ứng dụng bản đồ số và điều động lực lượng cảnh sát giao thông cắm chốt phân luồng từ xa.
\end{itemize}

Tuy nhiên, việc triển khai các công nghệ cảm biến đo đếm truyền thống như vòng từ cảm ứng (Inductive Loop Detectors) hay cảm biến radar dưới lòng đường bộc lộ những rào cản kỹ thuật và kinh tế không thể vượt qua trong bối cảnh giao thông tại Việt Nam:
\begin{enumerate}
    \item \textbf{Đặc thù dòng giao thông hỗn hợp xe máy (Motorcycle-Dominant Mixed Traffic):} Giao thông tại các đô thị Việt Nam mang tính chất hỗn hợp cao độ với tỷ lệ xe máy chiếm từ $75\%$ đến trên $85\%$ tổng lưu lượng. Khác với dòng xe ô tô có tính kỷ luật làn đường cao ở các nước phát triển, xe máy di chuyển tự do, đan xen, liên tục chuyển làn và luồn lách vào mọi khoảng trống trên mặt đường. Cảm biến vòng từ vốn được thiết kế để đếm xe ô tô đi qua một làn cố định sẽ mất hoàn toàn độ chính xác khi nhiều xe máy cùng đi qua một vòng từ hoặc di chuyển chệch giữa hai vòng từ kề nhau.
    \item \textbf{Chi phí hạ tầng và duy tu bảo dưỡng đắt đỏ:} Để phủ kín toàn bộ hàng nghìn giao lộ, việc xẻ đường chôn vòng từ đòi hỏi kinh phí đầu tư ban đầu khổng lồ, làm gián đoạn giao thông trong quá trình thi công và thường xuyên bị hư hỏng khi mặt đường bị sụt lún hoặc đào bới do sửa chữa hệ thống cấp thoát nước, cáp viễn thông.
    \item \textbf{Tiềm năng từ mạng lưới camera giám sát (CCTV):} Ngược lại, hầu hết các đô thị lớn hiện nay đều đã đầu tư hạ tầng camera giám sát giao thông công cộng với độ phủ rất cao tại các nút giao trọng điểm. Nguồn cấp hình ảnh từ camera cung cấp thông tin thị giác trực quan, phong phú, chi phí khai thác thấp và hoàn toàn không xâm lấn vào kết cấu mặt đường.
\end{enumerate}

Mặc dù camera mang lại tiềm năng to lớn, việc ứng dụng luồng dữ liệu hình ảnh vào bài toán dự báo mạng lưới đồ thị đô thị lại đặt ra một thách thức khoa học hoàn toàn mới: \textbf{Hiện tượng Lan truyền Sai số Nhận dạng (Error Propagation).} Số đếm phương tiện thu nhận từ hình ảnh camera thực chất là các giá trị ước lượng mang sai số nhận dạng (\textit{Perception Noise}), bắt nguồn từ hiện tượng che khuất tương hỗ nghiêm trọng giữa các cụm xe máy đông đúc. Sự lan truyền sai số này vào các mô hình mạng nơ-ron đồ thị không -- thời gian đa bước diễn ra như thế nào, và làm thế nào để thiết kế một kiến trúc học sâu có khả năng chống chịu nhiễu bền bỉ là một câu hỏi mở chưa từng được định lượng bài bản trong các nghiên cứu quốc tế.

Xuất phát từ những đòi hỏi cấp bách của thực tiễn quản lý đô thị và khoảng trống nghiên cứu khoa học nêu trên, đề tài \textbf{"Nghiên cứu và phát triển khung mô hình thu nhận hình ảnh và dự báo dòng giao thông đô thị đa bước hai giai đoạn dựa trên mạng đồ thị không -- thời gian từ camera giám sát"} đã được đề xuất và thực hiện.

\section*{2. Mục tiêu Nghiên cứu}
Đề tài hướng tới việc giải quyết trọn vẹn chuỗi bài toán từ thu nhận hình ảnh thị giác máy tính đến dự báo chuỗi thời gian trên đồ thị giao thông thông qua 4 mục tiêu cụ thể:
\begin{enumerate}
    \item \textbf{Xây dựng bộ dữ liệu chuẩn mực IC4SD-Traffic-HCM:} Thu thập dữ liệu thực tế từ hệ thống camera công cộng TP. Hồ Chí Minh, xây dựng bộ dữ liệu ảnh gán nhãn thủ công theo phân lớp (ô tô và xe máy) và chuỗi thời gian giao thông liên tục 86 ngày trên mạng lưới 608 nút camera để làm nền tảng kiểm chứng mở cho cộng đồng nghiên cứu.
    \item \textbf{Thiết kế khung làm việc hai giai đoạn (Two-Stage Framework):} 
    \begin{itemize}
        \item \textit{Giai đoạn 1 (Vision Perception):} Phát triển mô hình hồi quy đếm xe trực tiếp từ ảnh tĩnh độ phân giải cao, giải quyết triệt để vấn đề che khuất nghiêm trọng của xe máy và triệt tiêu điểm nghẽn tính toán NMS.
        \item \textit{Giai đoạn 2 (Graph Forecasting):} Mô hình hóa mạng lưới đường phố thành đồ thị không gian có hướng dựa trên bản đồ OpenStreetMap và dự báo lưu lượng đa bước thời gian ($t+1$ đến $t+6$).
    \end{itemize}
    \item \textbf{Đề xuất kiến trúc cải tiến TA-STGCN:} Tích hợp cơ chế Chú ý tự thân theo thời gian đa đầu (Multi-Head Temporal Self-Attention) vào cấu trúc mạng nơ-ron tích chập đồ thị không -- thời gian, giúp tăng cường khả năng chống chịu nhiễu nhận dạng và kiểm soát tốc độ suy thoái sai số khi dự báo dài hạn.
    \item \textbf{Định lượng có hệ thống hiện tượng Lan truyền Sai số và Khả năng Chống chịu Nhiễu:} Xây dựng mô hình toán học tiêm nhiễu cảm biến thực tế, thực hiện các kiểm định thống kê chuyên sâu (Friedman, Wilcoxon Signed-Rank, Cohen's $d$) để đánh giá sự phân hóa độ bền vững giữa các dòng kiến trúc đồ thị khác nhau.
\end{enumerate}

\section*{3. Đối tượng và Phạm vi Nghiên cứu}
\begin{itemize}
    \item \textbf{Đối tượng nghiên cứu:}
    \begin{itemize}
        \item Các kiến trúc mạng nơ-ron tích chập sâu (CNNs) và Vision Transformers cho bài toán thị giác máy tính.
        \item Lý thuyết đồ thị phổ (Spectral Graph Theory) và mạng nơ-ron đồ thị không -- thời gian (Spatio-Temporal Graph Neural Networks -- STGNNs).
        \item Cơ chế Chú ý tự thân (Self-Attention Mechanism) và lý thuyết ước lượng thống kê bền vững (Robust Statistics).
        \item Các mô hình phân phối sai số nhận dạng thực tế (Heteroscedastic Asymmetric Gamma Distribution).
    \end{itemize}
    \item \textbf{Phạm vi nghiên cứu:}
    \begin{itemize}
        \item \textit{Không gian:} Mạng lưới 608 trạm camera giám sát phủ kín khu vực nội thành và các cửa ngõ giao thông huyết mạch của Thành phố Hồ Chí Minh.
        \item \textit{Thời gian:} Dữ liệu chuỗi thời gian liên tục trong 86 ngày, từ ngày 23/06/2025 đến ngày 16/09/2025 (tổng cộng 24.768 bước thời gian với chu kỳ 5 phút).
        \item \textit{Đối tượng phương tiện:} Tập trung vào hai loại hình phương tiện cơ giới chủ đạo chiếm trên $98\%$ lưu lượng đường bộ đô thị: Xe ô tô con (\textit{Passenger Cars}) và Xe máy (\textit{Motorcycles}).
    \end{itemize}
\end{itemize}

\section*{4. Phương pháp Luận Nghiên cứu Khoa học}
Đề tài tuân thủ phương pháp luận nghiên cứu khoa học thực nghiệm kết hợp với mô hình hóa toán học chặt chẽ:
\begin{enumerate}
    \item \textbf{Phương pháp thu thập và xử lý dữ liệu:} Thu thập tự động qua giao diện lập trình của Cổng thông tin giao thông TP.HCM; áp dụng giao thức gán nhãn 3 thẩm định viên độc lập kết hợp kiểm định độ tin cậy bằng Hệ số tương quan nội lớp (Intraclass Correlation Coefficient -- ICC); áp dụng bộ lọc nhân quả Causal Double-Rolling Filter để làm mượt dữ liệu chuỗi mà không gây rò rỉ dữ liệu tương lai.
    \item \textbf{Phương pháp mô hình hóa toán học:} Biểu diễn mạng lưới đường phố thành đồ thị có hướng $\mathcal{G} = (\mathcal{V}, \mathcal{E}, W)$ dựa trên tô pô đường bộ thực tế trích xuất từ OpenStreetMap; áp dụng phép xấp xỉ bộ lọc phổ ChebNet trên toán tử Laplacian để truyền tin không gian cục bộ; ứng dụng giải tích ma trận để chứng minh hiệu ứng giảm phương sai nhiễu của cơ chế Attention thời gian.
    \item \textbf{Phương pháp thực nghiệm và kiểm định thống kê chuẩn mực:} Huấn luyện tất cả các mô hình qua 5 seed ngẫu nhiên độc lập; đo đạc sai số trên giá trị đếm xe thực tế sau giải chuẩn hóa; sử dụng kiểm định phi tham số Friedman omnibus và kiểm định cặp Wilcoxon Signed-Rank kèm hiệu chỉnh Holm-Bonferroni; đo lường quy mô ảnh hưởng thực chất bằng chỉ số Cohen's $d$ và khoảng tin cậy Bootstrap 95\%.
\end{enumerate}

\section*{5. Những Đóng góp Khoa học và Thực tiễn của Đề tài}
\begin{enumerate}
    \item \textbf{Đóng góp bộ dữ liệu chuẩn mực IC4SD-Traffic-HCM:} Công bố bộ dữ liệu giao thông quy mô đô thị mở đầu tiên cho khu vực Đông Nam Á bảo toàn thông tin phân loại phương tiện ô tô -- xe máy trên mạng lưới 608 nút camera, tạo môi trường thử nghiệm chuẩn mực cho cộng đồng học thuật quốc tế.
    \item \textbf{Đột phá trong kỹ thuật đếm xe hỗn hợp:} Đề xuất và kiểm chứng giải pháp hồi quy đếm xe trực tiếp bằng EfficientNet-B5, vượt trội về tốc độ tính toán và độ bền bỉ trước hiện tượng che khuất so với các phương pháp phát hiện và bám vết truyền thống trong điều kiện xe máy dày đặc.
    \item \textbf{Đề xuất kiến trúc TA-STGCN tối ưu:} Thiết kế thành công mô hình tích hợp cơ chế Chú ý tự thân theo thời gian vào STGCN, đạt được sự cân bằng xuất sắc giữa độ chính xác dự báo, số lượng tham số gọn nhẹ (453K params) và tính kinh tế bộ nhớ GPU (tiết kiệm $25.2\%$ VRAM so với STAEformer).
    \item \textbf{Luận giải khoa học về Lan truyền sai số và Cơ chế lọc nhiễu:} Đưa ra bằng chứng thực nghiệm và lý thuyết đầu tiên chứng minh cơ chế Attention theo trục thời gian hoạt động như một bộ lọc nhiễu tự nhiên, làm phẳng đường cong suy thoái sai số trước các biến động cảm biến camera, vượt trội áp đảo các mô hình không có chú ý (Cohen's $d$ lên tới $+5.84$).
    \item \textbf{Hiện thực hóa kiến trúc Edge--Cloud:} Đề xuất mô hình triển khai phân tán giúp tiết kiệm băng thông mạng truyền thông đô thị tới $62.500\times$ và bảo vệ tuyệt đối quyền riêng tư hình ảnh của người dân.
\end{enumerate}

\section*{6. Bố cục của Cuốn Báo cáo}
Nội dung cuốn báo cáo NCKH được tổ chức thành 6 chương chuyên sâu, kết luận và tài liệu tham khảo:
\begin{itemize}
    \item \textbf{Phần Mở đầu:} Giới thiệu bối cảnh, tính cấp thiết, mục tiêu, đối tượng, phương pháp luận và các đóng góp cốt lõi của đề tài.
    \item \textbf{Chương 1: Tổng quan và Bối cảnh Nghiên cứu:} Khảo sát toàn diện về ITS, bài toán dự báo dòng giao thông đô thị hỗn hợp, hiện trạng các phương pháp đếm xe, các kiến trúc học sâu không -- thời gian và xác định khoảng trống nghiên cứu.
    \item \textbf{Chương 2: Cơ sở Lý thuyết và Nền tảng Toán học:} Trình bày chi tiết, tường minh toán học về CNNs, lý thuyết đồ thị phổ, đa thức Chebyshev, tích chập thời gian GLU, cơ chế Multi-Head Attention, hàm mất mát Huber Loss và lý thuyết giảm phương sai nhiễu.
    \item \textbf{Chương 3: Thu thập Dữ liệu và Xây dựng Đồ thị Giao thông (IC4SD-Traffic-HCM):} Chi tiết quy trình thu thập camera 608 nút, cấu trúc bộ dữ liệu ảnh gán nhãn, chuỗi thời gian, xây dựng đồ thị OSM-RBF và quy trình kiểm soát rò rỉ dữ liệu Zero-Leakage.
    \item \textbf{Chương 4: Phương pháp Nghiên cứu Đề xuất:} Mô tả chi tiết khung làm việc 2 giai đoạn, mô hình hồi quy EfficientNet-B5 và kiến trúc mạng nơ-ron đồ thị TA-STGCN cùng thuật toán huấn luyện.
    \item \textbf{Chương 5: Thực nghiệm, Kết quả và Đánh giá:} Trình bày kết quả đo đạc Stage 1, Stage 2 qua 5 seeds ngẫu nhiên, phân tích lan truyền sai số theo mật độ, nghiên cứu độ bền vững với nhiễu Noise Robustness, phân tích ma trận trọng số Attention và Ablation study.
    \item \textbf{Chương 6: Thảo luận, Mô hình Triển khai và Hướng Phát triển:} Phân tích kiến trúc Edge--Cloud tiết kiệm băng thông $62.500\times$, tính kinh tế tài nguyên GPU, tính công bằng thực nghiệm, giới hạn và hướng mở rộng tương lai.
    \item \textbf{Kết luận và Kiến nghị:} Tổng kết kết quả nghiên cứu và đề xuất ứng dụng thực tiễn.
    \item \textbf{Tài liệu tham khảo:} Danh mục các công trình khoa học quốc tế uy tín được trích dẫn.
\end{itemize}


% Chương 1: Tổng quan và bối cảnh nghiên cứu
\chapter{Tổng quan và bối cảnh nghiên cứu}
\label{chap:tong_quan}

\section{Tổng quan về hệ thống giao thông thông minh (ITS)}
Hệ thống giao thông thông minh (Intelligent Transportation Systems -- ITS) là sự kết hợp giữa công nghệ thông tin, truyền thông và các thuật toán điều khiển hiện đại nhằm nâng cao hiệu quả vận hành mạng lưới giao thông, bảo đảm an toàn, tối ưu hóa năng lượng và giảm thiểu ô nhiễm môi trường đô thị.

Theo các khung kiến trúc chuẩn mực quốc tế (như ISO 14813), một hệ thống ITS đô thị thường gồm ba phân hệ chính:
\begin{enumerate}
    \item \textbf{Hệ thống quản lý giao thông nâng cao (ATMS):} Giám sát trạng thái mạng lưới đường bộ thời gian thực qua hệ thống cảm biến, phát hiện sự cố tự động, điều khiển đèn tín hiệu thích ứng và điều phối lực lượng xử lý tình huống phát sinh.
    \item \textbf{Hệ thống thông tin người tham gia giao thông (ATIS):} Cung cấp thông tin cập nhật liên tục về tình trạng lưu thông, thời gian hành trình dự kiến và hướng dẫn lộ trình tối ưu thông qua biển báo điện tử VMS, ứng dụng di động và đài phát thanh giao thông.
    \item \textbf{Hệ thống vận tải công cộng nâng cao (APTS):} Hỗ trợ quản lý và tối ưu biểu đồ chạy xe buýt, metro, cung cấp thời gian đến bến dự kiến (ETA) và điều phối ưu tiên qua nút giao cho phương tiện công cộng.
\end{enumerate}

Trong đó, bài toán \textbf{dự báo dòng giao thông đô thị (urban traffic forecasting)} đóng vai trò trung tâm trong phân hệ ATMS và ATIS. Việc dự báo chính xác trạng thái mạng lưới trước 15 đến 30 phút cho phép các nhà quản lý chuyển từ trạng thái phản ứng bị động sang điều hành chủ động, kịp thời điều chỉnh chu kỳ đèn tín hiệu và phân luồng từ xa trước khi ùn tắc lan rộng.

\subsection{Tiến trình phát triển của hệ thống giao thông thông minh}
Quá trình phát triển của ITS trải qua ba giai đoạn công nghệ chính:
\begin{enumerate}
    \item \textbf{Thế hệ ITS 1.0 (1970 -- 1990): Kỷ nguyên cảm biến cơ điện cố định.}
    Giai đoạn này gắn liền với các cảm biến vòng từ (inductive loop detectors) chôn dưới mặt đường và hệ thống đèn tín hiệu phối hợp cố định theo thời gian trong ngày (như phương pháp Webster). Dữ liệu thu được chủ yếu phục vụ các tuyến cao tốc, thu thập lưu lượng tổng hợp trên từng làn và chưa có khả năng liên kết mạng lưới không gian cũng như phân loại phương tiện.
    \item \textbf{Thế hệ ITS 2.0 (1990 -- 2015): Kỷ nguyên truyền thông không dây và định vị vệ tinh.}
    Sự phổ biến của hệ thống GPS và mạng di động mở ra nguồn dữ liệu từ phương tiện di động (Floating Car Data -- FCD). Các hệ thống điều khiển đèn thích ứng như SCOOT (Anh) và SCATS (Úc) bắt đầu ứng dụng mô hình hồi quy để điều hòa lưu lượng. Tuy nhiên, chi phí đầu tư và bản quyền của các giải pháp này rất lớn, chủ yếu thiết kế cho dòng xe ô tô chiếm ưu thế tại các quốc gia phát triển.
    \item \textbf{Thế hệ ITS 3.0 (Từ 2015 đến nay): Kỷ nguyên AIoT, thị giác máy tính và học sâu đồ thị.}
    Sự phát triển mạnh mẽ của thị giác máy tính và mạng nơ-ron học sâu giúp tận dụng trực tiếp hạ tầng camera giám sát công cộng (CCTV) sẵn có tại các đô thị. Dữ liệu hình ảnh được chuyển hóa thành các luồng số liệu giao thông liên tục, kết hợp với các mô hình mạng nơ-ron đồ thị không -- thời gian (STGNN) để dự báo toàn mạng lưới theo thời gian thực với chi phí phần cứng hợp lý.
\end{enumerate}

\section{Động học dòng giao thông đô thị và bài toán dự báo đa bước}
\subsection{Đặc tính động học phi tuyến và tính phụ thuộc không -- thời gian}
Dòng giao thông đô thị là một hệ thống động học phi tuyến phức tạp, chịu sự chi phối đồng thời của hai yếu tố tương quan:
\begin{enumerate}
    \item \textbf{Tương quan thời gian (temporal dependencies):} Lưu lượng xe mang tính quy luật rõ rệt theo chu kỳ ngày đêm và chu kỳ tuần (giữa ngày làm việc và ngày nghỉ), đồng thời có quán tính ngắn hạn. Khi dòng phương tiện đông đúc xuất hiện tại một thời điểm, trạng thái này có xu hướng duy trì trong các khoảng thời gian kế tiếp trước khi giải tỏa.
    \item \textbf{Tương quan không gian (spatial dependencies):} Giao thông vận hành trên mạng lưới đường phố liên thông có cấu trúc đồ thị. Trạng thái lưu thông tại một nút giao phụ thuộc trực tiếp vào dòng xe từ các nút thượng lưu đổ về và khả năng thoát xe tại các nút hạ lưu. Mối tương quan này bị ràng buộc bởi hình học mạng lưới, khoảng cách giữa các giao lộ, quy định đường một chiều và số làn xe.
\end{enumerate}

\subsection{Đặc thù dòng giao thông hỗn hợp và động học di chuyển của xe máy}
Tại các đô thị ở Việt Nam, dòng giao thông mang những đặc thù khác biệt đáng kể so với mô hình giao thông tại các nước phát triển:
\begin{itemize}
    \item \textbf{Dòng chảy phi làn (non-lane based flow):} Khác với giả định phương tiện di chuyển bám làn thành hàng dọc trong các lý thuyết giao thông cổ điển, xe máy tại Việt Nam di chuyển linh hoạt theo dòng chảy tự do, liên tục tận dụng các khoảng trống giữa các xe ô tô để vượt lên.
    \item \textbf{Đặc thù di chuyển theo cụm mật độ cao:} Khi dừng chờ đèn đỏ tại các ngã tư, hàng trăm xe máy dồn ứ lại thành cụm dày đặc phía trước ô tô. Khi đèn chuyển xanh, cụm xe này đồng loạt bứt tốc với gia tốc lớn, tạo ra dao động lưu lượng đột biến trong chu kỳ quan sát 5 phút.
    \item \textbf{Hệ số xe con quy đổi biến thiên:} Trong điều kiện lưu thông bình thường, một xe máy tương đương khoảng 0.2 -- 0.25 xe con quy đổi (PCU). Tuy nhiên trong điều kiện ùn ứ tốc độ thấp, xe máy luồn lách dễ dàng hơn trong khi ô tô chiếm dụng diện tích mặt đường lớn. Do đó, việc tách biệt chính xác số lượng ô tô và xe máy là yêu cầu quan trọng để phản ánh đúng hiện trạng tải trọng đường bộ.
\end{itemize}

\subsection{Ý nghĩa của tầm nhìn dự báo 15--30 phút}
Bài toán dự báo dòng giao thông đa bước đặt mục tiêu ước lượng lưu lượng phương tiện qua $H$ bước thời gian liên tiếp trong tương lai:
\begin{equation}
    \hat{\mathbf{Y}}_{t+1}, \hat{\mathbf{Y}}_{t+2}, \dots, \hat{\mathbf{Y}}_{t+H}
\end{equation}
với chu kỳ lấy mẫu tiêu chuẩn $\Delta t = 5$ phút, tầm nhìn dự báo $H = 6$ bước tương ứng với 30 phút tới.

Tầm nhìn dự báo 15--30 phút có ý nghĩa thực tiễn quan trọng:
\begin{itemize}
    \item \textbf{Thời gian đáp ứng của chu kỳ đèn tín hiệu:} Các giải pháp điều khiển đèn phối hợp cần khoảng thời gian tương đương 3 đến 5 chu kỳ đèn để chuyển đổi kế hoạch pha mượt mà mà không gây gián đoạn dòng xe. Dự báo trước 15--30 phút giúp trung tâm điều hành kịp thời điều chỉnh thời lượng đèn xanh đón đầu dòng xe đang tiếp cận.
    \item \textbf{Mô hình lan truyền sóng tắc nghẽn:} Theo mô hình thủy động lực học giao thông (LWR), khi lưu lượng vượt quá dung lượng mặt đường, sóng tắc nghẽn sẽ lan truyền ngược chiều dòng xe với vận tốc khoảng 15--25 km/h. Dự báo sớm 30 phút cho phép hệ thống kịp thời thông báo và điều tiết từ xa, hỗ trợ phân luồng phương tiện trước khi điểm nghẽn hình thành.
\end{itemize}

\section{Khảo cứu các kỹ thuật nhận dạng và đếm phương tiện từ camera}
Để cung cấp số liệu lưu lượng cho các mô hình dự báo mạng lưới, khâu thị giác máy tính đóng vai trò thu thập tín hiệu đầu vào. Hiện nay có ba hướng tiếp cận chính:

\subsection{Tiếp cận bám vết theo phát hiện (tracking-by-detection)}
Phương pháp này kết hợp mô hình phát hiện vật thể (như YOLO, Faster R-CNN) với thuật toán bám vết liên khung hình (như ByteTrack, DeepSORT).
\begin{itemize}
    \item \textbf{Nguyên lý:} Tại mỗi khung hình, mô hình phát hiện xác định các bounding box của từng xe. Thuật toán bám vết liên kết quỹ đạo phương tiện qua các khung hình liên tiếp dựa trên bộ lọc Kalman và chỉ số IoU (Intersection-over-Union). Khi một quỹ đạo cắt qua vạch ảo định sẵn, hệ thống ghi nhận một lượt xe.
    \item \textbf{Hạn chế trong giao thông xe máy mật độ cao:}
    \begin{enumerate}
        \item \textit{Hiện tượng che khuất lẫn nhau:} Vào giờ cao điểm, các xe máy đi sát nhau khiến tỷ lệ che khuất thân xe rất lớn, dẫn đến việc mô hình phát hiện bỏ sót hoặc sinh ra các hộp nhận dạng sai lệch.
        \item \textit{Chi phí tính toán của khâu NMS:} Khi số lượng xe trong khung hình lên tới hàng trăm, mô hình sinh ra hàng nghìn bounding box ứng viên. Thuật toán triệt tiêu cực đại cục bộ (NMS) với độ phức tạp $\mathcal{O}(M^2)$ trở thành điểm nghẽn tính toán trên các thiết bị biên.
        \item \textit{Gián đoạn thời gian trên camera công cộng:} Để tiết kiệm băng thông hạ tầng mạng, nhiều camera giao thông chỉ truyền ảnh tĩnh định kỳ (ví dụ 60 giây/ảnh), làm mất tính liên tục khung hình và khiến các thuật toán bám vết mất dấu chuyển động.
    \end{enumerate}
\end{itemize}

\subsection{Tiếp cận ước lượng bản đồ mật độ (density map estimation)}
Xuất phát từ bài toán đếm đám đông (crowd counting), các mô hình như CSRNet, CAN chuyển đổi bài toán đếm rời rạc thành bài toán hồi quy bản đồ mật độ liên tục:
\begin{itemize}
    \item \textbf{Nguyên lý:} Mỗi phương tiện được biểu diễn bằng một điểm tâm, bản đồ mật độ được tạo ra qua phép tích chập với hàm Gaussian. Tổng số lượng xe được tính bằng tích phân trên toàn bộ bức ảnh.
    \item \textbf{Hạn chế:} Bản đồ mật độ không phân biệt được cụ thể loại phương tiện (ô tô hay xe máy). Trong khi đó, việc mất thông tin phân lớp gây khó khăn cho việc đánh giá tải trọng giao thông. Ngoài ra, việc gán nhãn điểm tâm cho tập dữ liệu lớn đòi hỏi chi phí nhân công đáng kể.
\end{itemize}

\subsection{Tiếp cận hồi quy đếm xe trực tiếp (direct regression)}
Nhằm khắc phục các nhược điểm trên, đề tài lựa chọn và phát triển hướng tiếp cận \textbf{hồi quy đếm xe trực tiếp theo phân lớp}.
\begin{itemize}
    \item \textbf{Nguyên lý:} Mô hình học sâu nhận đầu vào là ảnh tĩnh $\mathbf{I} \in \mathbb{R}^{H \times W \times 3}$ và ánh xạ trực tiếp sang vector số lượng hai chiều:
    \begin{equation}
        f_\theta(I) \longrightarrow \hat{\mathbf{c}} = \begin{bmatrix} \hat{c}_{\text{car}} \\ \hat{c}_{\text{bike}} \end{bmatrix} \in \mathbb{R}^2
    \end{equation}
    \item \textbf{Ưu điểm:}
    \begin{enumerate}
        \item \textit{Không cần khâu NMS:} Mô hình không sinh bounding box, loại bỏ độ phức tạp tính toán bậc hai và đạt tốc độ xử lý nhanh (dưới 35 ms/ảnh trên phần cứng phổ thông).
        \item \textit{Bảo toàn phân lớp phương tiện:} Tách biệt rõ ràng số lượng ô tô và xe máy, cung cấp đầu vào chính xác cho mô hình dự báo không -- thời gian.
        \item \textit{Phù hợp với luồng ảnh tĩnh định kỳ:} Xử lý độc lập trên từng ảnh, không đòi hỏi video liên tục băng thông cao.
        \item \textit{Khắc phục hiện tượng học bám cảnh nền:} Áp dụng các kỹ thuật tăng cường dữ liệu ngẫu nhiên (Color Jittering, Random Erasing) và kiểm chứng sự tập trung đặc trưng qua bản đồ nhiệt Grad-CAM++.
    \end{enumerate}
\end{itemize}

\section{Khảo cứu các mô hình dự báo chuỗi thời gian giao thông}
Bài toán dự báo dòng giao thông không -- thời gian đã trải qua các giai đoạn phát triển chính:

\subsection{Các mô hình thống kê truyền thống}
Trong giai đoạn đầu, các phương pháp thống kê cổ điển giữ vai trò chủ đạo:
\begin{itemize}
    \item \textbf{ARIMA và SARIMA:} Phân tích chuỗi thời gian dựa trên các thành phần tự hồi quy, sai phân và trung bình trượt dưới giả định chuỗi dừng và tuyến tính.
    \item \textbf{VAR (Vector AutoRegression):} Mở rộng ARIMA cho trường hợp đa biến, mô hình hóa tương quan giữa một số nút lân cận.
    \item \textbf{SVR và Gradient Boosting (XGBoost, LightGBM):} Ứng dụng các hàm nhân phi tuyến hoặc cây quyết định để hồi quy lưu lượng từ các đặc trưng trễ trong quá khứ.
\end{itemize}
Tuy nhiên, các phương pháp này gặp khó khăn khi mạng lưới mở rộng lên hàng trăm nút ($N = 608$) do số lượng tham số tăng nhanh, đồng thời thiếu cơ chế biểu diễn mối quan hệ topo không gian của mạng lưới đường bộ.

\subsection{Các mô hình học sâu tuần tự (RNN, LSTM, GRU)}
Với sự phát triển của học sâu, các kiến trúc mạng hồi quy tuần tự được áp dụng:
\begin{itemize}
    \item \textbf{LSTM và GRU:} Khắc phục hiện tượng triệt tiêu gradient của RNN truyền thống qua các cơ chế cổng điều tiết thông tin, ghi nhớ các phụ thuộc thời gian dài hạn hơn.
    \item \textbf{TCN (Temporal Convolutional Networks):} Sử dụng tích chập 1D giãn nở nhân quả, cho phép tính toán song song trên GPU và mở rộng trường tiếp nhận theo cấp số nhân.
\end{itemize}
Mặc dù vậy, các mô hình này chủ yếu xử lý độc lập từng nút hoặc làm phẳng toàn bộ mạng lưới thành một vector không có cấu trúc, bỏ qua thông tin hình học và liên kết vật lý giữa các tuyến đường.

\subsection{Mạng nơ-ron đồ thị không -- thời gian (STGNN)}
Để mô hình hóa đồng thời cả chiều không gian và thời gian, mạng nơ-ron đồ thị không -- thời gian (STGNN) đã ra đời và trở thành hướng nghiên cứu chủ đạo. Các kiến trúc STGNN tiêu biểu gồm:
\begin{enumerate}
    \item \textbf{Tích chập phổ trên đồ thị tĩnh:} Tiêu biểu là STGCN (Yu et al., 2018), sử dụng xấp xỉ ChebNet bậc $K=3$ trên ma trận Laplacian của đồ thị khoảng cách, kết hợp với các khối tích chập thời gian 1D Gated Linear Units (GLU). STGCN có tốc độ tính toán nhanh và khả năng định vị cục bộ tốt, nhưng phụ thuộc vào cấu trúc đồ thị tĩnh ban đầu.
    \item \textbf{Tích chập khuếch tán có hướng:} Tiêu biểu là DCRNN (Li et al., 2018), mô hình hóa dòng giao thông như một quá trình bước ngẫu nhiên có hướng trên đồ thị, kết hợp với tế bào GRU theo kiến trúc Encoder-Decoder. Hạn chế là tốc độ huấn luyện chậm do tính tuần tự của GRU.
    \item \textbf{Học ma trận kề thích ứng:} Tiêu biểu là GraphWaveNet (Wu et al., 2019), bổ sung ma trận kề tự học để nắm bắt tương quan tiềm ẩn giữa các nút xa nhau, kết hợp với tích chập thời gian giãn nở (TCN). Nhược điểm là ma trận kề tự học dạng dày làm tăng chi phí bộ nhớ khi số nút lớn và nhạy cảm với nhiễu.
    \item \textbf{Cơ chế chú ý không gian và thời gian phân tách:} Tiêu biểu là ASTGCN (Guo et al., 2019), tính toán ma trận chú ý động theo cả chiều không gian và thời gian. Tuy nhiên số lượng tham số lớn (1.53 triệu tham số) và chi phí tính toán tăng nhanh theo số lượng nút.
    \item \textbf{Đồ thị động và đồ thị meta:} Tiêu biểu là DSTAGNN (Lan et al., 2022) và MegaCRN (Jiang et al., 2023), xây dựng đồ thị tương quan động dựa trên khoảng cách Wasserstein hoặc phân tách biểu diễn đồ thị meta. Các mô hình này đạt độ chính xác tốt trên dữ liệu chuẩn nhưng dễ bị suy giảm khi đầu vào có nhiễu cảm biến lớn.
    \item \textbf{Mô hình Transformer không -- thời gian:} Các công trình như STAEformer (Liu et al., 2023) và iTransformer (Liu et al., 2024) ứng dụng cơ chế chú ý tự thân trên toàn chuỗi hoặc đảo ngược theo chiều biến số. Mặc dù có năng lực biểu diễn cao, các mô hình này đòi hỏi bộ nhớ VRAM đỉnh lớn (trên 11 GB với đồ thị 608 nút) và dễ bị suy thoái sai số ở các bước dự báo xa khi dữ liệu đầu vào mang nhiễu thực tế.
\end{enumerate}

\section{Hiện tượng lan truyền sai số trong hệ thống học máy nối tầng}
Trong các hệ thống học máy nối tầng, kết quả đầu ra của khâu nhận dạng đóng vai trò là dữ liệu đầu vào trực tiếp cho khâu dự báo:
\begin{equation}
    \mathbf{I} \xrightarrow{\quad f_{\theta_1} \quad} \hat{\mathbf{Y}} = \mathbf{Y}^* + \boldsymbol{\eta} \xrightarrow{\quad g_{\theta_2} \quad} \hat{\mathbf{Z}}
\end{equation}
Hiện tượng \textbf{lan truyền sai số (error propagation)} xuất hiện khi sai số nhận dạng $\boldsymbol{\eta}$ từ khâu $f_{\theta_1}$ không bị triệt tiêu mà lan truyền qua các tầng nơ-ron của khâu $g_{\theta_2}$, làm lệch lạc kết quả dự báo lưu lượng.

\subsection{Cơ chế tích lũy và khuếch đại sai số qua các bước thời gian}
Trong bài toán giao thông, hiện tượng này thể hiện ở hai khía cạnh:
\begin{enumerate}
    \item \textbf{Lan truyền sai số theo chiều không gian:} Phép tích chập đồ thị $g_\theta \star_{\mathcal{G}} \mathbf{x} = \sum_{k=0}^{K-1} \theta_k T_k(\tilde{\mathbf{L}}) \mathbf{x}$ khuếch tán thông tin trạng thái của một nút sang các nút lân cận $(K-1)$-hop. Nếu một camera bị lỗi nhận dạng cục bộ (ví dụ do mưa lớn hoặc lóa đèn pha), sai số này sẽ lan sang các nút xung quanh nếu mô hình không có cơ chế lọc phù hợp.
    \item \textbf{Tích lũy sai số theo chiều thời gian:} Khi dự báo đa bước từ $t+1$ đến $t+6$, sai lệch nhận dạng trong cửa sổ lịch sử 24 bước có thể làm biến dạng xu hướng ngắn hạn, dẫn đến việc ước lượng sai lệch tăng dần ở các bước nhìn xa.
\end{enumerate}

\subsection{Yêu cầu về độ bền vững với nhiễu nhận dạng}
Do đó, việc chỉ đánh giá mô hình trên dữ liệu sạch lý tưởng là chưa đủ đối với các hệ thống camera thực tế. Mô hình dự báo cần có \textbf{độ bền vững với nhiễu (noise robustness)}, tức năng lực tự động làm mượt và giảm thiểu ảnh hưởng của sai số đầu vào ngay trong cấu trúc nơ-ron. Đây là mục tiêu thiết kế then chốt của mô hình TA-STGCN.

\section{Đối sánh các kiến trúc dự báo giao thông tiêu biểu}
Để làm rõ bối cảnh nghiên cứu, Bảng \ref{tab:so_sanh_sota_toan_dien} tổng hợp các đặc tính kỹ thuật chính giữa các kiến trúc tiêu biểu hiện nay và mô hình đề xuất trong đề tài.

\begin{table}[htbp]
\centering
\footnotesize
\caption{Đối sánh đặc điểm các kiến trúc dự báo giao thông không -- thời gian tiêu biểu.}
\label{tab:so_sanh_sota_toan_dien}
\renewcommand{\arraystretch}{1.2}
\setlength{\tabcolsep}{3.5pt}
\begin{tabularx}{\textwidth}{l >{\raggedright\arraybackslash}p{2.6cm} >{\raggedright\arraybackslash}p{2.3cm} c >{\raggedright\arraybackslash}X c}
\toprule
\textbf{Mô hình} & \textbf{Cấu trúc không gian} & \textbf{Cơ chế thời gian} & \textbf{Cảm biến} & \textbf{Bối cảnh thực nghiệm} & \textbf{Phân loại} \\
\midrule
DCRNN (2018) & Khuếch tán đồ thị & GRU tuần tự & Vòng từ & Cao tốc ô tô (PeMS) & Không \\
STGCN (2018) & ChebNet ($K=3$) & Khối 1D GLU & Vòng từ & Cao tốc ô tô (PeMS) & Không \\
Graph WaveNet (2019) & Đồ thị thích ứng & Dilated TCN & Vòng từ & Cao tốc ô tô (PeMS) & Không \\
ASTGCN (2019) & ChebNet + Chú ý & Chú ý thời gian & Vòng từ & Cao tốc ô tô (PeMS) & Không \\
DSTAGNN (2022) & Đồ thị động đa kênh & Transformer thời gian & Vòng từ & Cao tốc ô tô (PeMS) & Không \\
MegaCRN (2023) & Meta-Graph thích ứng & GRU Encoder--Decoder & Vòng từ & Cao tốc ô tô (PeMS) & Không \\
STAEformer (2023) & Nhúng không gian & ST-Transformer & Vòng từ & Cao tốc ô tô (PeMS) & Không \\
iTransformer (2024) & Đảo chiều biến số & Inverted MLP & Chuỗi & Chuỗi đa biến chung & Không \\
\midrule
\textbf{TA-STGCN (Đề xuất)} & \textbf{ChebNet ($K=3$) thực} & \textbf{MHTA + GLU} & \textbf{Camera (608 nút)} & \textbf{Đô thị ($>80\%$ xe máy)} & \textbf{Có (2 lớp)} \\
\bottomrule
\end{tabularx}
\end{table}

\section{Khoảng trống nghiên cứu}
Từ khảo cứu trên, đề tài xác định ba khoảng trống nghiên cứu chính:
\begin{enumerate}
    \item \textbf{Thiếu tập dữ liệu chuẩn mực về giao thông hỗn hợp quy mô lớn:} Các tập dữ liệu phổ biến (PeMSD4, PeMSD8, METR-LA) đều thu thập từ cảm biến vòng từ trên cao tốc tại Mỹ, đặc trưng bởi dòng xe ô tô đơn nhất đi đúng làn. Chưa có tập dữ liệu chuẩn mực nào phục vụ bài toán dự báo toàn diện mạng lưới đô thị hỗn hợp với xe máy chiếm ưu thế trên quy mô toàn thành phố (608 camera).
    \item \textbf{Mất thông tin phân lớp phương tiện:} Hầu hết các mô hình hiện nay chỉ dự báo một kênh lưu lượng gộp chung hoặc tốc độ trung bình, làm mất đi đặc trưng động học riêng biệt của ô tô và xe máy.
    \item \textbf{Thiếu phân tích thực nghiệm về lan truyền sai số nhận dạng:} Các nghiên cứu trước đây thường giả định dữ liệu đầu vào là chính xác tuyệt đối. Vấn đề sai số nhận dạng từ camera ảnh hưởng như thế nào đến độ chính xác dự báo đa bước và cơ chế kiến trúc nào giúp nâng cao độ bền vững với nhiễu vẫn chưa được khảo sát đầy đủ.
\end{enumerate}


% Chương 2: Cơ sở lý thuyết và nền tảng toán học
\chapter{Cơ sở lý thuyết và nền tảng toán học}
\label{chap:co_so_ly_thuyet}

Chương này trình bày hệ thống nền tảng toán học, giải tích đồ thị phổ và lý thuyết học sâu làm cơ sở khoa học cho kiến trúc Two-Stage Pipeline của đề tài. Cấu trúc chương được tổ chức tập trung vào hai trụ cột công nghệ cốt lõi:
\begin{enumerate}
    \item \textbf{Mạng xương sống EfficientNet và Mô hình hồi quy đếm xe trực tiếp (Giai đoạn 1):} Trình bày cơ sở toán học của bài toán đếm xe qua ánh xạ trực tiếp từ không gian ảnh tĩnh sang vector mật độ phương tiện theo phân lớp; phân tích nguyên lý mở rộng đồng thời (Compound Scaling), cấu trúc khối tích chập MBConv tích hợp cơ chế chú ý kênh Squeeze-and-Excitation, tầng nén đặc trưng Bottleneck MLP và hàm mất mát bền vững Huber Loss.
    \item \textbf{Mạng nơ-ron đồ thị không -- thời gian nhận biết thời gian TA-STGCN (Giai đoạn 2):} Trình bày mô hình hóa mạng lưới giao thông trên đồ thị phi Euclid; phân rã giải tích đồ thị phổ với toán tử xấp xỉ Chebyshev, cơ chế tích chập thời gian có cổng GLU; và kiến trúc đề xuất TA-STGCN tích hợp cơ chế Self-Attention theo thời gian đa đầu (Multi-Head Temporal Self-Attention) giải quyết hiện tượng lan truyền sai số cảm biến.
\end{enumerate}

% ==============================================================================
\section{Mạng xương sống EfficientNet và Mô hình hồi quy đếm xe trực tiếp}
\label{sec:kien_truc_efficientnet}

\subsection{Tổng quan nguyên lý và bài toán đếm xe qua hồi quy trực tiếp}
\label{sec:nguyen_ly_hoi_quy_dem_xe}

Trong giám sát giao thông đô thị bằng thị giác máy tính, khung làm việc truyền thống thường dựa trên mô hình \textit{Tracking-by-Detection} (kết hợp các mạng phát hiện hộp bao như Faster R-CNN, YOLO với thuật toán bám vết liên khung hình như SORT, DeepSORT). Tuy nhiên, tại các đô thị đặc thù ở Việt Nam với tỷ trọng phương tiện cá nhân hai bánh chiếm trên $80\%$, mô hình này bộc lộ những giới hạn căn bản:
\begin{enumerate}
    \item \textbf{Che khuất tương hỗ dày đặc (Severe Mutual Occlusion):} Dòng xe máy di chuyển thành cụm với cự ly rất hẹp, dẫn đến hiện tượng phần lớn thân xe bị che khuất. Các thuật toán sinh vùng bao thường xuyên bỏ sót phương tiện hoặc gộp nhiều xe máy liền kề vào cùng một hộp định vị.
    \item \textbf{Nút thắt tính toán của thuật toán triệt tiêu cực đại cục bộ (NMS):} Khi mật độ xe trong khung hình vượt quá ngưỡng cao (trên 100 phương tiện), số lượng hộp bao dự đoán sinh ra tăng đột biến. Thuật toán NMS với độ phức tạp tính toán tiệm cận $\mathcal{O}(M^2)$ (với $M$ là số hộp bao ứng viên) tạo thành điểm nghẽn trễ lớn, không đáp ứng yêu cầu xử lý thời gian thực trên các thiết bị biên.
\end{enumerate}

Để khắc phục triệt để các hạn chế trên, đề tài xây dựng bài toán đếm xe theo hướng \textbf{Hồi quy trực tiếp từ ảnh tĩnh} (Direct Density Regression). Về mặt toán học, bài toán được phát biểu dưới dạng tìm ánh xạ phi tuyến tham số hóa $f_\theta$:
\begin{equation}
    f_\theta: \mathbf{I} \in \mathbb{R}^{3 \times H \times W} \longrightarrow \hat{\mathbf{c}} = \begin{bmatrix} \hat{c}_{\text{car}} \\ \hat{c}_{\text{bike}} \end{bmatrix} \in \mathbb{R}^2
\end{equation}
trong đó $\mathbf{I}$ là tensor biểu diễn ảnh chụp tĩnh từ camera quan sát có kích thước $H \times W$ trên 3 kênh màu RGB; $\hat{\mathbf{c}}$ là vector số thực hai chiều ước lượng số lượng ô tô và xe máy hiện diện trong khung cảnh.

Mô hình học trực tiếp mối quan hệ giữa các mẫu hoa văn mật độ điểm ảnh toàn cục và tổng lượng phương tiện mà không cần định vị tọa độ từng xe riêng lẻ. Cách tiếp cận này triệt tiêu hoàn toàn khâu hậu xử lý NMS, đưa độ trễ suy luận về mức cố định (dưới 35 ms trên GPU biên), đồng thời duy trì độ chính xác cao ngay cả khi các phương tiện bị che khuất nghiêm trọng.

\subsection{Nguyên lý mở rộng đồng thời (Compound Scaling) của EfficientNet}

Để đảm nhận vai trò mạng xương sống (backbone) trích xuất đặc trưng hình ảnh cho bài toán hồi quy đếm xe, mô hình đòi hỏi khả năng biểu diễn phong phú song phải tối ưu hóa tài nguyên tính toán. Mạng nơ-ron tích chập truyền thống thường cải thiện độ chính xác bằng cách mở rộng độc lập một trong ba chiều: chiều sâu mạng ($d$), chiều rộng số kênh ($w$), hoặc độ phân giải ảnh đầu vào ($r$). Tuy nhiên, Tan và Le \cite{efficientnet} đã chỉ ra rằng việc tăng đơn lẻ từng chiều nhanh chóng dẫn đến hiện tượng suy giảm hiệu suất biên (diminishing returns).

Nhằm khắc phục nhược điểm trên, nguyên lý \textbf{Compound Scaling} điều chỉnh đồng bộ cả ba chiều thông qua một hệ số mở rộng duy nhất $\phi$:
\begin{equation}
    d = \alpha^\phi, \quad w = \beta^\phi, \quad r = \gamma^\phi
\end{equation}
với ràng buộc cân bằng mức tăng trưởng phép tính:
\begin{equation}
    \alpha \cdot \beta^2 \cdot \gamma^2 \approx 2, \quad \alpha \ge 1, \ \beta \ge 1, \ \gamma \ge 1
\end{equation}
trong đó $\alpha, \beta, \gamma$ là các hằng số phân bổ xác định qua tìm kiếm kiến trúc trên mạng cơ sở EfficientNet-B0. Bằng cách mở rộng cân đối theo $\phi$, kiến trúc duy trì tỷ lệ tối ưu giữa kích thước trường thụ cảm và dung lượng kênh, đạt độ chính xác trích xuất cao trên cùng mức chi phí tính toán (FLOPs).

\subsection{Cấu trúc khối MBConv và cơ chế chú ý kênh Squeeze-and-Excitation}

Đơn vị cấu trúc cốt lõi của EfficientNet là khối \textbf{MBConv (Mobile Inverted Bottleneck Convolution)}, tích hợp ba kỹ thuật tối ưu:

\textbf{1. Tích chập tách rời theo chiều sâu (Depthwise Separable Convolution):}
Phép toán phân tách tích chập tiêu chuẩn thành hai bước độc lập: \textit{Depthwise Convolution} (áp dụng bộ lọc $k \times k$ riêng cho từng kênh) và \textit{Pointwise Convolution} (tích chập $1 \times 1$ để kết hợp các kênh). Tỷ số chi phí tính toán đạt được là:
\begin{equation}
    \frac{\text{FLOPs}_{\text{separable}}}{\text{FLOPs}_{\text{standard}}} = \frac{1}{C_{\text{out}}} + \frac{1}{k^2}
\end{equation}
Với kernel kích thước $k = 3$, phương pháp này cắt giảm khoảng $85\% - 89\%$ lượng phép tính so với tích chập 2D thông thường mà vẫn bảo toàn năng lực biểu diễn không gian.

\textbf{2. Khối phần dư nghịch đảo (Inverted Residuals):}
Khác với khối Bottleneck truyền thống (nén kênh ở giữa), MBConv mở rộng số kênh ở tầng giữa lên gấp $t$ lần ($t=6$) trước khi áp dụng Depthwise Convolution, sau đó chiếu ngược về số kênh thấp kết hợp đường dẫn tắt (shortcut connection), giúp bảo toàn tối đa thông tin ngữ nghĩa.

\textbf{3. Cơ chế chú ý kênh Squeeze-and-Excitation (SE):}
Khối SE tự động tái hiệu chuẩn trọng số các kênh đặc trưng qua hai pha:
\begin{align}
    \mathbf{z} &= \text{GAP}(\mathbf{X}) \in \mathbb{R}^C \quad &\text{(Pha Squeeze)} \\
    \mathbf{s} &= \sigma\left( \mathbf{W}_2 \cdot \text{SiLU}(\mathbf{W}_1 \mathbf{z}) \right) \in \mathbb{R}^C \quad &\text{(Pha Excitation)}
\end{align}
với $\mathbf{W}_1 \in \mathbb{R}^{\frac{C}{r} \times C}$ (tỷ lệ co thắt $r=4$), $\mathbf{W}_2 \in \mathbb{R}^{C \times \frac{C}{r}}$, và $\sigma(\cdot)$ là hàm Sigmoid. Bản đồ đặc trưng đầu ra $\tilde{\mathbf{X}}_c = s_c \cdot \mathbf{X}_c$ giúp mô hình khuếch đại các kênh nhạy cảm với mật độ phương tiện và kìm hãm phản hồi từ cảnh nền tĩnh.

\subsection{Cấu trúc đầu hồi quy nén đặc trưng (Bottleneck MLP)}
\label{sec:bottleneck_mlp}

\textbf{Tầng gộp trung bình toàn cục (Global Average Pooling -- GAP):}
Tại tầng cuối cùng của mạng xương sống, tensor đặc trưng $\mathbf{F} \in \mathbb{R}^{C \times H_f \times W_f}$ được tổng hợp thành vector biểu diễn $\mathbf{f} \in \mathbb{R}^C$ thông qua phép lấy trung bình không gian:
\begin{equation}
    f_c = \frac{1}{H_f \cdot W_f} \sum_{i=1}^{H_f} \sum_{j=1}^{W_f} \mathbf{F}_c(i, j) \quad (\forall c = 1, \dots, C)
\end{equation}
Phép biến đổi này không sinh thêm tham số học được, loại bỏ nguy cơ quá khớp của các tầng kết nối dày đặc cổ điển, đồng thời bảo đảm tính bất biến đối với vị trí xuất hiện của xe trên ảnh.

\textbf{Cấu trúc đầu hồi quy Bottleneck MLP:}
Từ vector ngữ nghĩa $\mathbf{f} \in \mathbb{R}^{2048}$ xuất ra từ EfficientNet-B5, đầu hồi quy được thiết kế theo cấu trúc nén thắt cổ chai:
\begin{align}
    \mathbf{h}_1 &= \text{LayerNorm}(\mathbf{W}_1 \mathbf{f} + \mathbf{b}_1), \quad &\mathbf{W}_1 \in \mathbb{R}^{128 \times 2048}, \ \mathbf{b}_1 \in \mathbb{R}^{128} \\
    \mathbf{h}_2 &= \text{ReLU}(\mathbf{h}_1) \in \mathbb{R}^{128} \\
    \hat{\mathbf{c}} &= \mathbf{W}_2 \mathbf{h}_2 + \mathbf{b}_2, \quad &\mathbf{W}_2 \in \mathbb{R}^{2 \times 128}, \ \mathbf{b}_2 \in \mathbb{R}^2
\end{align}
Tầng thắt cổ chai 128 chiều đóng vai trò như một bộ lọc thông tin (Information Bottleneck), ép mạng phải tổng hợp các đặc trưng trừu tượng phản ánh mật độ và loại bỏ các đặc trưng cục bộ thừa thãi. Tầng Layer Normalization giúp ổn định phân phối kích hoạt trước các biến thiên điều kiện ánh sáng mạnh giữa ban ngày và ban đêm.

\subsection{Hàm mất mát bền vững Huber Loss}
\label{sec:huber_loss}
Trong điều kiện quan trắc camera thực địa, dữ liệu hình ảnh chịu ảnh hưởng từ các yếu tố môi trường cực đoan (lóa đèn xe rọi thẳng vào ống kính ban đêm, giọt nước mưa làm nhòe kính hoặc camera bị rung lắc do gió lớn), dẫn đến sự xuất hiện của các điểm ngoại lai (outliers) có sai số đếm bất thường.

Nếu sử dụng hàm mất mát trung bình bình phương sai số ($L_2$ Loss):
\begin{equation}
    \mathcal{L}_2(e) = \frac{1}{2} e^2, \quad \frac{\partial \mathcal{L}_2}{\partial e} = e
\end{equation}
gradient tăng tuyến tính theo sai số $e = c - \hat{c}$. Khi gặp một điểm ngoại lai lớn ($|e| \gg 1$), gradient bùng nổ sẽ làm biến dạng các trọng số học được, gây mất ổn định mô hình. Ngược lại, hàm mất mát sai số tuyệt đối ($L_1$ Loss) có đạo hàm không liên tục tại $e=0$, gây dao động số học quanh điểm tối ưu.

Mô hình áp dụng hàm mất mát \textbf{Huber Loss} với ngưỡng chuyển tiếp $\delta = 1.0$:
\begin{equation}
    \mathcal{L}_{\mathrm{Huber}}(e) =
    \begin{cases}
    \frac{1}{2} e^2, & \text{khi } |e| < \delta, \\[4pt]
    \delta \left( |e| - \frac{1}{2} \delta \right), & \text{khi } |e| \ge \delta
    \end{cases}
\end{equation}
Hàm ảnh hưởng (Influence Function -- đạo hàm bậc nhất của hàm mất mát) có dạng:
\begin{equation}
    \psi(e) = \frac{\partial \mathcal{L}_{\mathrm{Huber}}}{\partial e} =
    \begin{cases}
    e, & \text{khi } |e| < \delta, \\[4pt]
    \delta \cdot \text{sign}(e), & \text{khi } |e| \ge \delta
    \end{cases}
\end{equation}
Ở lân cận sai số nhỏ ($|e| < \delta$), hàm cư xử như $L_2$ Loss cho phép gradient triệt tiêu mượt mà về 0, bảo đảm tính hội tụ chính xác. Ở vùng sai số lớn ($|e| \ge \delta$), gradient bị chặn trên bởi hằng số $\pm \delta$, cô lập tác động phá hủy của các điểm dị biệt.

Hàm mất mát tổng thể để tối ưu hóa Giai đoạn 1 là:
\begin{equation}
    \mathcal{L}_{\text{Stage 1}} = \frac{1}{2} \left( \mathcal{L}_{\mathrm{Huber}}(c_{\text{car}} - \hat{c}_{\text{car}}) + \mathcal{L}_{\mathrm{Huber}}(c_{\text{bike}} - \hat{c}_{\text{bike}}) \right)
\end{equation}

% ==============================================================================
\section{Mạng nơ-ron đồ thị không -- thời gian nhận biết thời gian TA-STGCN}
\label{sec:kien_truc_ta_stgcn}

\subsection{Tổng quan bài toán dự báo dòng giao thông không -- thời gian trên đồ thị}
\label{sec:bieu_dien_khong_gian_thoi_gian}

Dòng giao thông trên mạng lưới đường bộ đô thị vận hành trong không gian phi Euclid có cấu trúc liên kết hình học phức tạp. Sự tương tác giữa các luồng xe không chỉ phụ thuộc vào khoảng cách địa lý Euclid thuần túy mà bị chi phối mạnh mẽ bởi cấu trúc hình học của đường (chiều lưu thông, dải phân cách, số làn và các nút giao cắt). 

Mạng lưới giao thông gồm $N$ vị trí quan sát được biểu diễn dưới dạng một đồ thị có hướng, có trọng số $\mathcal{G} = (\mathcal{V}, \mathcal{E}, \mathbf{W})$, trong đó:
\begin{itemize}
    \item $\mathcal{V} = \{v_1, v_2, \dots, v_N\}$ là tập hợp $N$ đỉnh đồ thị, tương ứng với $N$ vị trí cảm biến / camera thu thập dữ liệu trên mạng lưới.
    \item $\mathcal{E} \subseteq \mathcal{V} \times \mathcal{V}$ là tập hợp các cạnh có hướng, thể hiện hành lang lưu thông trực tiếp từ nút $v_i$ sang nút $v_j$.
    \item $\mathbf{W} \in \mathbb{R}^{N \times N}$ là ma trận trọng số không gian. Phần tử $W_{ij} > 0$ biểu thị mức độ liên kết không gian giữa hai nút $v_i$ và $v_j$ dựa trên khoảng cách mạng lưới thực tế; $W_{ij} = 0$ nếu không tồn tại kết nối trực tiếp.
\end{itemize}

Tại mỗi bước thời gian $t$, trạng thái giao thông trên toàn mạng lưới được mô tả bởi ma trận $\mathbf{X}_t \in \mathbb{R}^{N \times C}$, trong đó $C$ là số lượng thuộc tính quan trắc (chẳng hạn lưu lượng đếm được của từng loại xe). Xét cửa sổ quan sát lịch sử gồm $T$ bước thời gian liên tiếp, dữ liệu đầu vào của bài toán là tensor ba chiều:
\begin{equation}
    \mathcal{X} = [\mathbf{X}_{t-T+1}, \mathbf{X}_{t-T+2}, \dots, \mathbf{X}_t] \in \mathbb{R}^{T \times N \times C}
\end{equation}

Mục tiêu của bài toán dự báo không -- thời gian là tìm ánh xạ hàm học được $\mathcal{F}_\Theta$ tham số hóa bởi $\Theta$:
\begin{equation}
    \hat{\mathcal{Y}} = [\hat{\mathbf{X}}_{t+1}, \hat{\mathbf{X}}_{t+2}, \dots, \hat{\mathbf{X}}_{t+H}] = \mathcal{F}_\Theta(\mathcal{X}; \mathcal{G}) \in \mathbb{R}^{H \times N \times C}
\end{equation}
nhằm dự báo chính xác trạng thái giao thông tại toàn bộ $N$ nút trên mạng lưới qua $H$ bước thời gian tiếp theo trong tương lai.

\subsection{Nền tảng giải tích đồ thị phổ và toán tử tích chập Chebyshev}
\label{sec:ly_thuyet_do_thi_pho}

Để trích xuất các tương quan không gian phi Euclid trên đồ thị $\mathcal{G}$, lý thuyết đồ thị phổ (Spectral Graph Theory) cung cấp cơ sở giải tích cho phép định nghĩa phép biến đổi Fourier và tích chập trên miền phổ của toán tử vi phân đồ thị.

\textbf{1. Toán tử Laplacian đối xứng chuẩn hóa:}
Cho ma trận trọng số $\mathbf{W} \in \mathbb{R}^{N \times N}$ và ma trận bậc đường chéo $\mathbf{D} \in \mathbb{R}^{N \times N}$ với $D_{ii} = \sum_{j=1}^N W_{ij}$. Toán tử Laplacian đối xứng chuẩn hóa $\mathbf{L} \in \mathbb{R}^{N \times N}$ được định nghĩa:
\begin{equation}
    \mathbf{L} = \mathbf{I}_N - \mathbf{D}^{-1/2} \mathbf{W} \mathbf{D}^{-1/2}
\end{equation}
trong đó $\mathbf{I}_N$ là ma trận đơn vị cấp $N$.

\textbf{Các đặc tính đại số quan trọng của $\mathbf{L}$:}
\begin{enumerate}
    \item $\mathbf{L}$ là ma trận thực, đối xứng và nửa xác định dương. Với mọi vector tín hiệu đồ thị $\mathbf{x} \in \mathbb{R}^N$:
    \begin{equation}
        \mathbf{x}^T \mathbf{L} \mathbf{x} = \frac{1}{2} \sum_{i=1}^N \sum_{j=1}^N W_{ij} \left( \frac{x_i}{\sqrt{D_{ii}}} - \frac{x_j}{\sqrt{D_{jj}}} \right)^2 \ge 0
    \end{equation}
    Do đó, toàn bộ các giá trị riêng của $\mathbf{L}$ đều là số thực không âm.
    \item Phổ giá trị riêng của $\mathbf{L}$ bị chặn trong đoạn hữu hạn:
    \begin{equation}
        0 = \lambda_1 \le \lambda_2 \le \dots \le \lambda_N \le 2
    \end{equation}
\end{enumerate}

Nhờ tính đối xứng thực, ma trận $\mathbf{L}$ có thể phân rã phổ trực chuẩn:
\begin{equation}
    \mathbf{L} = \mathbf{U} \boldsymbol{\Lambda} \mathbf{U}^T
\end{equation}
với $\mathbf{U} = [\mathbf{u}_1, \mathbf{u}_2, \dots, \mathbf{u}_N] \in \mathbb{R}^{N \times N}$ là ma trận trực giao chứa các vector riêng thỏa mãn $\mathbf{U} \mathbf{U}^T = \mathbf{I}_N$, và $\boldsymbol{\Lambda} = \text{diag}(\lambda_1, \dots, \lambda_N)$ là ma trận đường chéo chứa các giá trị riêng tương ứng.

\textbf{2. Biến đổi Fourier và định lý tích chập trên đồ thị:}
Các vector riêng $\mathbf{u}_k$ đóng vai trò như hệ hàm sóng cơ sở trên đồ thị. Phép biến đổi Fourier trên đồ thị của tín hiệu $\mathbf{x} \in \mathbb{R}^N$ và biến đổi nghịch đảo được xác định bởi:
\begin{equation}
    \hat{\mathbf{x}} = \mathbf{U}^T \mathbf{x}, \quad \mathbf{x} = \mathbf{U} \hat{\mathbf{x}}
\end{equation}
Theo định lý tích chập, phép tích chập giữa tín hiệu $\mathbf{x}$ và bộ lọc tham số $g_\theta$ trên đồ thị là phép nhân từng phần tử trong miền tần số:
\begin{equation}
    \mathbf{x} \ast_{\mathcal{G}} g_\theta = \mathbf{U} \left( (\mathbf{U}^T \mathbf{x}) \odot (\mathbf{U}^T g_\theta) \right) = \mathbf{U} g_\theta(\boldsymbol{\Lambda}) \mathbf{U}^T \mathbf{x}
\end{equation}
trong đó $g_\theta(\boldsymbol{\Lambda}) = \text{diag}(g_\theta(\lambda_1), \dots, g_\theta(\lambda_N))$ là ma trận đáp ứng phổ của bộ lọc.

\textit{Hạn chế tính toán:} Việc tính toán trực tiếp công thức trên đòi hỏi phép phân rã phổ của $\mathbf{L}$ với độ phức tạp $\mathcal{O}(N^3)$ và phép nhân ma trận dày đặc $\mathbf{U}$ với độ phức tạp $\mathcal{O}(N^2)$, hoàn toàn không khả thi trên các mạng lưới đô thị có số nút lớn.

\textbf{3. Xấp xỉ đa thức Chebyshev cắt cụt (ChebNet):}
Để triệt tiêu nút thắt tính toán trên, Defferrard et al. \cite{chebnet} đề xuất xấp xỉ bộ lọc phổ $g_\theta(\boldsymbol{\Lambda})$ bằng chuỗi đa thức trực giao Chebyshev loại I cắt cụt bậc $K-1$.

Đa thức Chebyshev $T_k(x)$ được định nghĩa trên đoạn $[-1, 1]$ thông qua công thức truy hồi:
\begin{equation}
    T_0(x) = 1, \quad T_1(x) = x, \quad T_k(x) = 2x T_{k-1}(x) - T_{k-2}(x) \quad (\forall k \ge 2)
\end{equation}
Thực hiện phép biến đổi affine đưa phổ trị riêng $[0, \lambda_{\max}]$ về $[-1, 1]$:
\begin{equation}
    \tilde{\mathbf{L}} = \frac{2}{\lambda_{\max}} \mathbf{L} - \mathbf{I}_N
\end{equation}
Bộ lọc phổ được xấp xỉ qua $K$ số hạng đầu tiên:
\begin{equation}
    g_\theta \ast_{\mathcal{G}} \mathbf{x} \approx \sum_{k=0}^{K-1} \theta_k T_k(\tilde{\mathbf{L}}) \mathbf{x}
\end{equation}
với $\theta \in \mathbb{R}^K$ là vector các tham số học được của bộ lọc.

\textbf{Định lý 2.1 (Tính định vị cục bộ $(K-1)$-hop):} \textit{Toán tử tích chập ChebNet bậc $K-1$ đảm bảo tính định vị không gian trong phạm vi đúng $K-1$ bước nhảy lân cận trên đồ thị.}
\begin{proof}
Ma trận $\tilde{\mathbf{L}}$ là hàm bậc nhất của ma trận kề $\mathbf{W}$. Với hai nút $i$ và $j$ bất kỳ, nếu khoảng cách ngắn nhất trên đồ thị $d_{\mathcal{G}}(i, j) > k$, thì phần tử $[\tilde{\mathbf{L}}^k]_{ij} = 0$. Vì $T_k(\tilde{\mathbf{L}})$ là đa thức đại số bậc $k$ của $\tilde{\mathbf{L}}$, nó được biểu diễn thành tổ hợp tuyến tính của các lũy thừa $\tilde{\mathbf{L}}^m$ ($m \le k$). Do đó, $[T_k(\tilde{\mathbf{L}})]_{ij} = 0$ khi $d_{\mathcal{G}}(i, j) > k$. Với bậc xấp xỉ $K=3$, thông tin giao thông chỉ lan truyền trong phạm vi 2-hop lân cận, hoàn toàn tương thích với động học lan truyền sóng giao thông thực tế.
\end{proof}

\textbf{Định lý 2.2 (Độ phức tạp thời gian tuyến tính):} \textit{Phép tích chập ChebNet có độ phức tạp thời gian $\mathcal{O}(K \cdot |\mathcal{E}|)$, tuyến tính theo số cạnh của đồ thị.}
\begin{proof}
Đại lượng $\mathbf{z}_k = T_k(\tilde{\mathbf{L}}) \mathbf{x}$ được tính toán đệ quy: $\mathbf{z}_0 = \mathbf{x}$, $\mathbf{z}_1 = \tilde{\mathbf{L}}\mathbf{x}$, $\mathbf{z}_k = 2\tilde{\mathbf{L}}\mathbf{z}_{k-1} - \mathbf{z}_{k-2}$. Mỗi bước đệ quy chỉ gồm phép nhân ma trận thưa $\tilde{\mathbf{L}}$ với vector $\mathbf{z}_{k-1}$. Với đồ thị mạng lưới đường bộ có cấu trúc thưa ($|\mathcal{E}| \ll N^2$), toàn bộ quá trình tính toán $K$ bước chỉ tiêu tốn chi phí $\mathcal{O}(K \cdot |\mathcal{E}|)$, loại bỏ hoàn toàn chi phí $\mathcal{O}(N^3)$ của phép phân rã phổ.
\end{proof}

\subsection{Tích chập thời gian có cổng (Temporal Convolution with GLU)}
\label{sec:tich_chap_thoi_gian_co_cong}

Để mô hình hóa chiều thời gian trong dữ liệu chuỗi giao thông, các mô hình hồi quy tuần hoàn (RNN, LSTM, GRU) thường gặp hai hạn chế kỹ thuật:
\begin{enumerate}
    \item \textbf{Nút thắt tính toán tuần tự (Sequential Bottleneck):} Việc cập nhật trạng thái ẩn $\mathbf{h}_t = f(\mathbf{h}_{t-1}, \mathbf{x}_t)$ buộc phải thực hiện tuần tự qua từng bước thời gian, không thể khai thác kiến trúc xử lý song song khối lượng lớn của GPU.
    \item \textbf{Suy thoái gradient (Gradient Degradation):} Việc nhân chuỗi liên tiếp các ma trận trọng số qua nhiều bước thời gian dài khiến gradient dễ bị triệt tiêu hoặc bùng nổ trong quá trình lan truyền ngược qua thời gian (BPTT).
\end{enumerate}

Để khắc phục nhược điểm này, mô hình áp dụng phép tích chập thời gian 1D kết hợp cơ chế \textbf{Cổng tuyến tính (Gated Linear Units -- GLU)}.

Xét tensor tín hiệu của một nút dọc theo trục thời gian $\mathbf{X} \in \mathbb{R}^{T \times C_{\text{in}}}$. Tầng tích chập thời gian áp dụng hai bộ lọc tích chập 1D độc lập $\boldsymbol{\Theta}_1, \boldsymbol{\Theta}_2 \in \mathbb{R}^{K_t \times C_{\text{in}} \times C_{\text{out}}}$ với độ dài kernel thời gian $K_t$:
\begin{equation}
    \Gamma \ast_{\mathcal{T}} \mathbf{X} = (\mathbf{X} \star \boldsymbol{\Theta}_1 + \mathbf{b}_1) \odot \sigma(\mathbf{X} \star \boldsymbol{\Theta}_2 + \mathbf{b}_2)
\end{equation}
trong đó $\star$ là phép tích chập 1D dọc trục thời gian; $\odot$ là phép nhân Hadamard từng phần tử; và $\sigma(z) = (1 + e^{-z})^{-1}$ là hàm Sigmoid đóng vai trò van cổng đóng mở dòng thông tin.

\textbf{Phân tích giải tích luồng gradient của GLU:}
Đặt $\mathbf{P} = \mathbf{X} \star \boldsymbol{\Theta}_1 + \mathbf{b}_1$ là nhánh tuyến tính và $\mathbf{Q} = \mathbf{X} \star \boldsymbol{\Theta}_2 + \mathbf{b}_2$ là nhánh điều khiển cổng. Đầu ra của cổng là $\mathbf{Y} = \mathbf{P} \odot \sigma(\mathbf{Q})$. Đạo hàm của $\mathbf{Y}$ theo nhánh tuyến tính $\mathbf{P}$ là:
\begin{equation}
    \frac{\partial \mathbf{Y}}{\partial \mathbf{P}} = \sigma(\mathbf{Q})
\end{equation}
Khi van cổng mở ($\sigma(\mathbf{Q}) \to 1$), gradient $\frac{\partial \mathcal{L}}{\partial \mathbf{Y}}$ được truyền thẳng trực tiếp về $\mathbf{P}$ như một đường truyền tuyến tính, hoàn toàn không bị suy giảm bởi đạo hàm của các hàm phi tuyến tính như trong RNN thông thường. Cơ chế này cho phép huấn luyện các mạng tích chập sâu một cách ổn định, đồng thời toàn bộ chuỗi thời gian được xử lý song song với tốc độ cao.

\subsection{Cơ chế Self-Attention theo thời gian đa đầu (Multi-Head Temporal Self-Attention)}

Mặc dù tích chập thời gian GLU tính toán nhanh và ổn định, kiến trúc STGCN tiêu chuẩn chỉ sử dụng kernel thời gian cục bộ nhỏ ($K_t = 3$). Thiết kế này bộc lộ một \textbf{điểm yếu chí tử trong hệ thống Two-Stage Pipeline}:
\begin{itemize}
    \item Sai số ước lượng từ Giai đoạn 1 (do rung lắc camera, chói sáng hoặc xe buýt che khuất tức thời) tạo ra các điểm đột biến dị thường (outliers) trong chuỗi đếm.
    \item Do trường tiếp nhận thời gian cục bộ hẹp, kernel tích chập GLU không thể nhận biết được xung nhiễu này là dị biệt so với xu thế toàn cục. Xung lỗi lập tức truyền thẳng vào tầng tích chập không gian ChebNet và lan truyền sang toàn bộ các nút lân cận trên đồ thị mạng lưới đường bộ (hiện tượng \textit{Error Propagation}).
\end{itemize}

Để giải quyết vấn đề này, đề tài đề xuất tích hợp \textbf{Cơ chế Self-Attention theo thời gian đa đầu (Multi-Head Temporal Self-Attention -- MHTA)} vào khối xử lý không -- thời gian.

\textbf{1. Thiết kế chú ý thời gian tách rời (Decoupled Temporal Attention):}
Thay vì tính toán cơ chế chú ý liên hợp trên toàn bộ không gian và thời gian $\mathbb{R}^{N \times T}$ (gây bùng nổ bộ nhớ với đồ thị quy mô lớn $N = 608$), cơ chế chú ý thời gian được thực thi độc lập cho từng đỉnh $i \in \{1, \dots, N\}$ dọc theo chuỗi quan sát $T$ bước:
\begin{equation}
    \mathbf{Z}_i = \mathbf{X}_{i, :, :} \in \mathbb{R}^{T \times C}
\end{equation}
Để mô hình nhận biết được thứ tự thời gian của các bước quan sát, vector mã hóa vị trí (Positional Encoding) được cộng trực tiếp vào chuỗi đặc trưng:
\begin{equation}
    \tilde{\mathbf{Z}}_i = \mathbf{Z}_i + \mathbf{PE}
\end{equation}
với các thành phần hàm điều hòa:
\begin{equation}
    \mathbf{PE}_{(pos, 2j)} = \sin\left(\frac{pos}{10000^{2j/C}}\right), \quad \mathbf{PE}_{(pos, 2j+1)} = \cos\left(\frac{pos}{10000^{2j/C}}\right)
\end{equation}

Chuỗi đặc trưng $\tilde{\mathbf{Z}}_i$ được ánh xạ tuyến tính thành ba ma trận Truy vấn (Query), Khóa (Key) và Giá trị (Value):
\begin{equation}
    \mathbf{Q}_i = \tilde{\mathbf{Z}}_i \mathbf{W}^Q, \quad \mathbf{K}_i = \tilde{\mathbf{Z}}_i \mathbf{W}^K, \quad \mathbf{V}_i = \tilde{\mathbf{Z}}_i \mathbf{W}^V
\end{equation}
trong đó $\mathbf{W}^Q, \mathbf{W}^K \in \mathbb{R}^{C \times d_k}$ và $\mathbf{W}^V \in \mathbb{R}^{C \times d_v}$ là các ma trận trọng số học được ($d_k = d_v = C / h$).

Ma trận trọng số chú ý thời gian $\mathbf{A}_i \in \mathbb{R}^{T \times T}$ được tính bằng tích vô hướng chia tỷ lệ:
\begin{equation}
    \mathbf{A}_i = \text{Softmax}\left( \frac{\mathbf{Q}_i \mathbf{K}_i^T}{\sqrt{d_k}} \right)
\end{equation}
Đầu ra của cơ chế chú ý tại đỉnh $i$ là:
\begin{equation}
    \text{Attention}(\mathbf{Q}_i, \mathbf{K}_i, \mathbf{V}_i) = \mathbf{A}_i \mathbf{V}_i
\end{equation}

\textbf{2. Cơ chế đa đầu (Multi-Head Attention):}
Để đồng thời trích xuất nhiều dạng quan hệ phụ thuộc thời gian (chu kỳ tín hiệu đèn, xu thế lưu lượng tăng/giảm và mức nền ổn định), cơ chế chú ý được triển khai song song trên $h = 4$ đầu chú ý độc lập:
\begin{align}
    \text{head}_m^{(i)} &= \text{Attention}(\tilde{\mathbf{Z}}_i \mathbf{W}_m^Q, \tilde{\mathbf{Z}}_i \mathbf{W}_m^K, \tilde{\mathbf{Z}}_i \mathbf{W}_m^V) \quad (m = 1, \dots, h) \\
    \text{MHTA}(\mathbf{Z}_i) &= \text{concat}\left( \text{head}_1^{(i)}, \text{head}_2^{(i)}, \dots, \text{head}_h^{(i)} \right) \mathbf{W}^O
\end{align}
với $\mathbf{W}^O \in \mathbb{R}^{C \times C}$ là ma trận chiếu đầu ra.

\textbf{Ý nghĩa giải tích của MHTA:} Cơ chế Self-Attention theo thời gian đóng vai trò như một \textbf{bộ lọc thông thấp thích ứng dữ liệu (Data-driven Adaptive Low-pass Filter)}. Khi một bước thời gian xuất hiện xung lỗi đột biến (outlier từ Giai đoạn 1), vector truy vấn và khóa của bước thời gian đó sẽ có độ tương đồng thấp với toàn bộ bối cảnh lịch sử còn lại trong cửa sổ $T$. Do đó, hàm Softmax sẽ tự động phân bổ trọng số chú ý cực nhỏ cho bước lỗi này, làm mượt xung nhiễu trước khi tín hiệu được truyền vào đồ thị không gian.

\subsection{Kiến trúc tích hợp TA-ST-Conv và phân tích độ phức tạp tiệm cận}
\label{sec:kien_truc_stgcn_tieu_chuan}

\textbf{1. Cấu trúc khối TA-ST-Conv hoàn chỉnh:}
Mỗi khối xử lý không -- thời gian trong TA-STGCN được cấu hình theo quy trình tích hợp ba tầng chặt chẽ:
\begin{align}
    \mathbf{Z}^{(l)} &= \text{LayerNorm}\left( \mathbf{X}^{(l)} + \text{MHTA}(\mathbf{X}^{(l)}) \right) \\
    \mathbf{S}^{(l)} &= \text{SpatialChebConv}\left( \mathbf{Z}^{(l)}; \tilde{\mathbf{L}}, K=3 \right) \\
    \mathbf{X}^{(l+1)} &= \text{LayerNorm}\left( \mathbf{S}^{(l)} + \text{TemporalGLU}(\mathbf{S}^{(l)}) \right)
\end{align}
Bằng cách bố trí khối MHTA ngay tại cửa ngõ đầu vào của khối trước toán tử ChebNet, mọi xung nhiễu cục bộ tại từng nút đều được làm sạch trước khi diễn ra quá trình khuếch tán thông tin trên đồ thị không gian. Tầng tích chập thời gian GLU đặt ở cuối đảm nhiệm vai trò cập nhật động học và điều tiết luồng gradient.

\textbf{2. Khối xuất kết quả (Output Block):}
Sau khi tín hiệu truyền qua hai khối TA-ST-Conv liên tiếp, kích thước chiều thời gian được nén về 1 thông qua một tầng tích chập thời gian 1D có kích thước kernel bằng đúng chiều dài chuỗi còn lại. Tầng tích chập $1 \times 1$ cuối cùng ánh xạ số kênh về đúng tầm nhìn dự báo $H = 6$ bước:
\begin{equation}
    \hat{\mathcal{Y}} = \text{Conv}_{1 \times 1}\left( \text{ReLU}\left( \text{Conv}_{\text{time}}(\mathbf{X}^{(2)}) \right) \right) \in \mathbb{R}^{H \times N \times 1}
\end{equation}

\textbf{3. Đối sánh kiến trúc và phân tích độ phức tạp tiệm cận:}
Bảng \ref{tab:so_sanh_kien_truc_stgcn} tổng hợp đặc tính cấu trúc và độ phức tạp tính toán giữa TA-STGCN và các mô hình STGNN kinh điển trong y văn.

\begin{table}[htbp]
\centering
\caption{So sánh đặc tính kiến trúc và độ phức tạp tính toán giữa các mô hình STGNN.}
\label{tab:so_sanh_kien_truc_stgcn}
\footnotesize
\setlength{\tabcolsep}{4pt}
\renewcommand{\arraystretch}{1.2}
\begin{tabularx}{\textwidth}{l >{\raggedright\arraybackslash}p{2.3cm} >{\raggedright\arraybackslash}p{2.2cm} >{\raggedright\arraybackslash}p{2.8cm} >{\raggedright\arraybackslash}X}
\toprule
\textbf{Mô hình} & \textbf{Không gian} & \textbf{Thời gian} & \textbf{Độ phức tạp} & \textbf{Khả năng lọc nhiễu cảm biến} \\
\midrule
DCRNN \cite{dcrnn} & Lan truyền ngẫu nhiên & GRU tuần tự & $\mathcal{O}(K |\mathcal{E}| T + N^2)$ & Thấp (nhiễu tích lũy qua các bước hồi quy BPTT). \\
STGCN \cite{stgcn} & ChebNet ($K=3$) & Gated CNN 1D & $\mathcal{O}(K |\mathcal{E}| T)$ & Trung bình (kernel cục bộ $K_t=3$, dễ bị sốc ngoại lai). \\
ASTGCN \cite{astgcn} & ChebNet + Spatial Attn & Gated CNN + Temp Attn & $\mathcal{O}(T N^2 C + K |\mathcal{E}| T)$ & Trung bình (chú ý không gian bị nhiễu cục bộ đánh lừa). \\
Graph WaveNet \cite{wavenet} & GCN tự học & Dilated Causal Conv & $\mathcal{O}(N^2 + K |\mathcal{E}| T)$ & Trung bình (không có cơ chế tự làm mượt trên trục thời gian). \\
\textbf{TA-STGCN (Đề xuất)} & \textbf{ChebNet ($K=3$)} & \textbf{MHTA + GLU} & $\mathcal{O}(N T^2 C + K |\mathcal{E}| T)$ & \textbf{Rất cao (cơ chế chú ý hoạt động như bộ lọc thông thấp thích ứng dữ liệu).} \\
\bottomrule
\end{tabularx}
\end{table}

\textbf{Ưu thế vượt trội về độ phức tạp của cơ chế chú ý tách rời:}
\begin{itemize}
    \item \textbf{Cơ chế chú ý liên hợp không -- thời gian (Joint Attention):} Nếu áp dụng cơ chế chú ý trên toàn bộ không gian tensor $(N \times T)$, ma trận chú ý có kích thước $(N \cdot T) \times (N \cdot T)$ với độ phức tạp $\mathcal{O}((N \cdot T)^2 \cdot C)$. Với $N = 608, T = 24$, số phần tử ma trận lên tới $(14.592)^2 \approx 2.13 \times 10^8$, gây tràn bộ nhớ GPU ngay từ bước lan truyền xuôi đầu tiên.
    \item \textbf{Cơ chế chú ý không gian (Spatial Attention):} Việc tính toán tương quan toàn thể giữa mọi cặp nút tại mỗi bước thời gian khiến độ phức tạp tăng bậc hai theo số lượng đỉnh $\mathcal{O}(T \cdot N^2 \cdot C)$, tạo thành điểm nghẽn nghiêm trọng khi mở rộng quy mô mạng lưới.
    \item \textbf{Cơ chế chú ý thời gian tách rời trong TA-STGCN:} Do phép chú ý được thực thi độc lập và song song trên từng đỉnh dọc theo trục thời gian, độ phức tạp của thành phần chú ý chỉ là $\mathcal{O}(N \cdot T^2 \cdot C)$, tăng tuyến tính theo số đỉnh đồ thị $N$. Trong các mạng lưới giao thông thực tế, số lượng nút quan trắc luôn lớn hơn rất nhiều so với độ dài cửa sổ thời gian lịch sử ($N \gg T$, cụ thể $608 \gg 24$). Do đó, thiết kế này bảo đảm khả năng mở rộng quy mô vượt trội và tiết kiệm bộ nhớ phần cứng tối đa.
\end{itemize}


% Chương 3: Thu thập dữ liệu và xây dựng đồ thị giao thông
\chapter{THU THẬP DỮ LIỆU VÀ XÂY DỰNG ĐỒ THỊ GIAO THÔNG (IC4SD-TRAFFIC-HCM)}
\label{chap:du_lieu}

Để giải quyết bài toán giao thông hỗn hợp tại các siêu đô thị Đông Nam Á và đánh giá một cách khoa học hiện tượng lan truyền sai số nhận dạng, nghiên cứu đã tiến hành thu thập, chuẩn hóa và công bố bộ dữ liệu chuẩn mực quy mô toàn thành phố mang tên \textbf{IC4SD-Traffic-HCM}. Bộ dữ liệu tích hợp hai thành phần: (1) Bộ dữ liệu thị giác hình ảnh gán nhãn chi tiết phục vụ huấn luyện và kiểm định mô hình đếm xe ở Giai đoạn 1; (2) Bộ dữ liệu chuỗi thời gian liên tục kết hợp đồ thị đường bộ thực tế phục vụ bài toán dự báo ở Giai đoạn 2.

\section{Hạ tầng Camera Giám sát và Hệ thống Thu thập Tự động}
\subsection{Hạ tầng Camera Giám sát TP. Hồ Chí Minh}
Thành phố Hồ Chí Minh là trung tâm kinh tế, tài chính lớn nhất của cả nước với diện tích hơn 2.061 $\text{km}^2$ và dân số thực tế hơn 10 triệu người. Mạng lưới đường bộ thành phố chịu áp lực cực lớn từ hơn 8.5 triệu xe máy cá nhân và gần 1 triệu xe ô tô. Để phục vụ công tác điều hành giao thông, Sở Giao thông Vận tải TP.HCM đã lắp đặt hệ thống camera giám sát giao thông (CCTV) phủ kín các tuyến đường huyết mạch, cầu vượt, hầm ngầm và các nút giao thông trọng điểm.

Nguồn cấp hình ảnh từ các camera này được công khai thông qua Cổng thông tin giao thông thành phố Hồ Chí Minh (\url{http://giaothong.hochiminhcity.gov.vn/}) dưới dạng các ảnh tĩnh chụp định kỳ, phục vụ giám sát dân sinh và nghiên cứu khoa học phi thương mại.

\subsection{Hệ thống Thu thập Dữ liệu Đa luồng Tự động}
Nhóm nghiên cứu đã phát triển một hệ thống thu thập dữ liệu tự động (Automated Ingestion Client) được tối ưu hóa cao độ, vận hành liên tục 24/7 trên máy chủ nghiên cứu chuyên dụng đặt tại Lab IC4SD:
\begin{enumerate}
    \item \textbf{Cơ chế Quản lý Phiên và Đa luồng (HTTP Session Pooling \& Multi-Threading):} Để thu thập đồng thời từ hàng trăm camera mà không làm nghẽn đường truyền hoặc bị máy chủ nguồn từ chối dịch vụ (Rate-limiting), chương trình sử dụng thư viện kết nối đa luồng bất đồng bộ kết hợp hàng đợi công việc (Thread Pool Worker). Các phiên kết nối HTTP được tái sử dụng để giảm thiểu độ trễ bắt tay TCP/TLS.
    \item \textbf{Đồng bộ hóa Chu kỳ Lấy mẫu Tuyệt đối (Strict 60-Second Cycle):} Cứ sau đúng 60 giây, chương trình kích hoạt một đợt quét đồng thời trên toàn bộ các endpoint camera đã đăng ký. Toàn bộ hình ảnh phản hồi từ hơn 600 camera được tải về trong cửa sổ thời gian từ 2 đến 4 giây, đảm bảo tính đồng bộ thời gian (Temporal Alignment) giữa tất cả các nút giao trên toàn thành phố.
    \item \textbf{Lọc trạm ngoại lệ và Xử lý Gián đoạn Kết nối:} Danh mục ban đầu ghi nhận 657 camera có tọa độ hợp lệ. Trong quá trình vận hành thử nghiệm 14 ngày, hệ thống tự động ghi nhận các mã lỗi phản hồi: mã lỗi 404 (mất tín hiệu), mã lỗi 503 (quá tải máy chủ), các khung hình trả về hình ảnh đen hoàn toàn, hình ảnh thông báo "Camera mất kết nối" hoặc camera bị lệch góc nhìn vào tường nhà dân. Các camera có tỷ lệ mất kết nối tích lũy vượt quá $5\%$ trong thời gian quan sát đã bị loại bỏ. Kết quả sàng lọc giữ lại đúng \textbf{608 trạm camera hoạt động hoàn hảo và liên tục}, tạo thành tập đỉnh không gian chính thức $\mathcal{V}$ ($N = |\mathcal{V}| = 608$).
\end{enumerate}

Hình \ref{fig:anh_samples_camera_chapter3} minh họa một số mẫu hình ảnh thu thập thực tế từ hệ thống, cho thấy sự đa dạng cao về góc chụp (từ trên cao nhìn xuống, góc chéo ngang mép đường), điều kiện ánh sáng (ban ngày rực rỡ, chạng vạng hoàng hôn, ban đêm dưới ánh đèn cao áp) và mức độ tập trung dày đặc của các luồng xe máy.

\begin{figure}[htbp]
\centering
\includegraphics[width=0.95\linewidth]{fig/samples_pictures_get_from_api.png}
\caption{Hình ảnh thực tế thu thập từ hệ thống camera giao thông TP.HCM thể hiện góc nhìn đa dạng và sự tập trung dày đặc của xe máy trong giờ cao điểm.}
\label{fig:anh_samples_camera_chapter3}
\end{figure}

\section{Bộ Dữ liệu Hình ảnh IC4SD-Traffic-HCM-Vision}
Từ kho dữ liệu hàng triệu khung hình thu thập được, nhóm nghiên cứu đã trích xuất một tập con đại diện gồm \textbf{5.012 ảnh chụp độ phân giải cao} ($1920 \times 1080$ hoặc $1280 \times 720$) để xây dựng bộ dữ liệu thị giác IC4SD-Traffic-HCM-Vision.

\subsection{Giao thức Gán nhãn Độc lập Ba Thẩm định viên}
Để loại bỏ thiên kiến cá nhân và sai số ngẫu nhiên của con người khi đếm các cảnh xe máy chen chúc, quy trình gán nhãn được tổ chức theo quy chuẩn nghiêm ngặt:
\begin{itemize}
    \item Ba thẩm định viên nghiên cứu độc lập (được đào tạo thống nhất về định nghĩa phương tiện và quy tắc xác định xe khuất) tiến hành gán nhãn toàn bộ 5.012 ảnh trên phần mềm chuyên dụng. Mỗi ảnh được gắn nhãn độc lập bởi cả ba chuyên viên.
    \item Định nghĩa hai phân lớp mục tiêu bảo toàn:
    \begin{enumerate}
        \item \textbf{Xe ô tô con (Passenger Cars):} Xe ô tô dưới 9 chỗ, xe taxi, xe SUV và xe bán tải đang lưu thông hoặc dừng chờ đèn tín hiệu trên lòng đường quan sát được.
        \item \textbf{Xe máy (Motorcycles):} Tất cả các loại xe mô tô hai bánh, xe tay ga, xe số và xe máy điện -- phương tiện chủ đạo trong dòng giao thông đô thị Việt Nam.
    \end{enumerate}
    \item Các ảnh có độ lệch số đếm giữa ba người vượt quá $5\%$ đối với ô tô hoặc vượt quá $10\%$ đối với xe máy được đưa vào phiên thẩm tra chung (Joint Reconciliation Session). Trong phiên này, cả ba chuyên viên cùng phóng to hình ảnh, rà soát từng cụm xe nghi vấn và thống nhất số đếm mặt đất (Ground Truth) cuối cùng.
\end{itemize}

\subsection{Đánh giá Độ tin cậy Gán nhãn bằng Hệ số Tương quan Nội lớp (ICC)}
Để định lượng mức độ tin cậy và sự thống nhất tuyệt đối của dữ liệu nhãn trước khi tiến hành phiên thẩm tra chung, nhóm nghiên cứu đã tính toán Hệ số tương quan nội lớp (Intraclass Correlation Coefficient -- ICC) trên một tập con ngẫu nhiên gồm 500 ảnh (chiếm xấp xỉ $10\%$ dữ liệu). 

Mô hình được áp dụng là $\text{ICC}(3, k)$ -- mô hình hiệu ứng hỗn hợp hai chiều (Two-way mixed-effects model) đánh giá sự đồng thuận tuyệt đối (Absolute Agreement) dựa trên giá trị trung bình của $k$ đánh giá ($k=3$):
\begin{equation}
    \text{ICC}(3, k) = \frac{\text{MS}_R - \text{MS}_E}{\text{MS}_R}
\end{equation}
trong đó $\text{MS}_R$ là phương sai trung bình giữa các ảnh (Mean Square for Rows) và $\text{MS}_E$ là phương sai trung bình của sai số dư thừa (Mean Square for Error).

\textbf{Kết quả đo đạc thống kê:}
\begin{itemize}
    \item Đối với phân lớp xe ô tô con: $\text{ICC} = 0.984$ (khoảng tin cậy 95\%: $0.981 - 0.987, p < 0.001$).
    \item Đối với phân lớp xe máy: $\text{ICC} = 0.962$ (khoảng tin cậy 95\%: $0.956 - 0.968, p < 0.001$).
    \item Tổng thể kết hợp cả hai lớp: $\text{ICC} = 0.973$ ($p < 0.001$).
\end{itemize}
Theo các tiêu chuẩn đo lường thống kê quốc tế (Koo \& Li, 2016), chỉ số ICC vượt trên $0.90$ phản ánh độ tin cậy "xuất sắc" (Excellent Reliability). Phân tích sâu hơn theo các điều kiện lưu thông chỉ ra rằng ở điều kiện đường vắng vào ban đêm, chỉ số ICC đạt tới $0.995$. Trong điều kiện ùn tắc cực đoan vào giờ cao điểm, chỉ số ICC của xe máy giảm nhẹ xuống $0.941$, phản ánh sự mập mờ thị giác khách quan khi các xe máy bị che khuất sâu trong đám đông chỉ để lộ một phần gương hoặc mũ bảo hiểm.

\subsection{Đặc tính Thống kê và Phân bố Khung giờ Diurnal}
Dữ liệu nhãn thể hiện tính chất đa dạng cao độ:
\begin{itemize}
    \item Trung bình số lượng phương tiện trên một ảnh là $30.28$ xe ($\text{Std} = 25.07$), giá trị trung vị là 24.00 xe, dải phân bố rộng từ 0 xe (đường vắng lúc rạng sáng) lên tới 194 xe (ngã tư giờ cao điểm chiều).
    \item Phân lớp xe máy áp đảo với trung bình $23.02$ xe/ảnh (cực đại 178 xe máy trong một khung hình), trong khi ô tô con đạt trung bình $7.26$ xe/ảnh (cực đại 50 xe).
    \item Dữ liệu có phân bố lệch phải mạnh (Right-skewed), phản ánh chính xác quy luật biến động của giao thông đô thị thực tế.
\end{itemize}

Bảng \ref{tab:thong_ke_vision_chapter3} mô tả chi tiết sự phân bổ của 5.012 ảnh gán nhãn qua 5 khung thời gian sinh hoạt trong ngày (Diurnal time windows), đảm bảo mô hình thị giác được tiếp xúc đồng đều với mọi điều kiện chiếu sáng và bối cảnh vận hành.

\begin{table}[htbp]
\centering
\caption{Phân bố các mẫu ảnh gán nhãn IC4SD-Traffic-HCM-Vision theo khung giờ hoạt động.}
\label{tab:thong_ke_vision_chapter3}
\small
\resizebox{\textwidth}{!}{%
\begin{tabular}{llcc}
\toprule
\textbf{Khung thời gian} & \textbf{Bối cảnh hoạt động thực tế} & \textbf{Số lượng ảnh} & \textbf{Tỷ lệ (\%)} \\
\midrule
00:00 -- 05:59 & Đêm khuya (Lưu lượng cực thấp, đường vắng, đèn vàng) & 1.311 & $26.16\%$ \\
07:00 -- 10:59 & Giờ cao điểm sáng (Dòng người đi làm, học sinh tới trường) & 964 & $19.23\%$ \\
14:00 -- 15:59 & Đầu giờ chiều (Dòng lưu thông ổn định, ánh nắng gắt) & 645 & $12.87\%$ \\
17:00 -- 18:59 & Giờ cao điểm chiều tối (Tắc nghẽn cục bộ, chạng vạng) & 656 & $13.09\%$ \\
21:00 -- 22:59 & Đêm muộn (Lưu lượng giảm dần, lóa đèn pha phương tiện) & 1.436 & $28.65\%$ \\
\midrule
\textbf{Tổng cộng} & \textbf{Tập dữ liệu ảnh chuẩn hóa gán nhãn đầy đủ} & \textbf{5.012} & $\mathbf{100.0\%}$ \\
\bottomrule
\end{tabular}%
}
\end{table}

\subsection{Kỹ thuật Tăng cường Dữ liệu Bảo toàn Nhãn (Label-Preserving Augmentation)}
Vì mục tiêu là hồi quy chính xác số lượng phương tiện có trong ảnh, các phương pháp cắt ngẫu nhiên (Random Crop) truyền thống sẽ làm mất một phần phương tiện ở rìa ảnh mà không thể cập nhật nhãn chính xác. Nhóm nghiên cứu thiết kế một đường ống tăng cường dữ liệu bảo toàn nhãn (Label-Preserving Pipeline):
\begin{enumerate}
    \item \textbf{Lật ngang ngẫu nhiên (Random Horizontal Flip):} Áp dụng với xác suất $p = 0.5$. Hướng di chuyển của xe bị đảo chiều (tương tự như nhìn qua gương) nhưng tổng số lượng ô tô và xe máy được bảo toàn nguyên vẹn $100\%$.
    \item \textbf{Biến đổi màu sắc trắc quang (Photometric Color Jittering):} Điều chỉnh ngẫu nhiên 4 thông số: độ sáng (Brightness $\pm 15\%$), độ tương phản (Contrast $\pm 15\%$), độ bão hòa màu (Saturation $\pm 15\%$) và sắc độ (Hue $\pm 15\%$). Phép biến đổi này mô phỏng sự biến thiên của điều kiện thời tiết tự nhiên (trời nắng chang chang, trời mưa râm mát, bóng râm của các tòa nhà cao tầng) mà không làm suy chuyển cấu trúc hình học của xe.
    \item \textbf{Xóa ngẫu nhiên diện tích nhỏ (Small-scale Random Erasing):} Áp dụng với xác suất $p = 0.2$, tỷ lệ diện tích bị xóa chỉ từ $2\%$ đến $20\%$ bề mặt ảnh. Kỹ thuật này giúp mô hình không bị quá phụ thuộc vào một chi tiết đơn lẻ, rèn luyện khả năng nhận dạng xe ngay cả khi bị cành cây ven đường hoặc biển báo giao thông che khuất một phần.
\end{enumerate}

\section{Bộ Dữ liệu Chuỗi Thời gian IC4SD-Traffic-HCM-Graph}
Để phục vụ bài toán dự báo đồ thị ở Giai đoạn 2, mô hình thị giác tối ưu ở Giai đoạn 1 (EfficientNet-B5) được triển khai để suy luận liên tục trên toàn bộ kho ảnh của 608 trạm camera trong suốt \textbf{86 ngày liên tục} (từ ngày 23/06/2025 đến ngày 16/09/2025). 

\subsection{Quy trình Lọc Nhân quả Causal Double-Rolling Filter}
Dữ liệu suy luận ở tần suất 1 phút phản ánh lưu lượng tức thời tại thời điểm bấm máy. Do góc chụp camera và chuyển động đèn tín hiệu, chuỗi 1 phút xuất hiện các dao động tần số cao (High-frequency jitter). Nhóm nghiên cứu áp dụng quy trình lọc nhân quả hai bước (Causal Double-Rolling Filter):
\begin{enumerate}
    \item \textbf{Lọc trung bình trượt hồi quy (Backward Rolling Mean):} Sử dụng cửa sổ trượt 3 phút nhìn ngược về quá khứ ($window_1 = 3$, $center = \text{False}$):
    \begin{equation}
        s_{1}(t) = \frac{1}{3} \sum_{k=0}^{2} y(t - k)
    \end{equation}
    \item \textbf{Làm mượt hàm mũ xuôi chiều (Forward-in-time EMA):} Áp dụng làm mượt hàm mũ với tham số $span = 5$ (tương đương hệ số làm mượt $\alpha = 2/(5+1) \approx 0.333$):
    \begin{equation}
        s_{2}(t) = \alpha s_{1}(t) + (1 - \alpha) s_{2}(t - 1)
    \end{equation}
\end{enumerate}
Chuỗi sau khi làm mượt được lấy mẫu lại định kỳ đúng 5 phút một lần ($\Delta t = 5$ phút). Quy trình này đảm bảo tính nhân quả tuyệt đối: giá trị làm mượt tại thời điểm $t$ hoàn toàn không phụ thuộc vào bất kỳ điểm dữ liệu nào trong tương lai $t+1, t+2, \dots$

Kết quả thu được là một tensor dữ liệu chuỗi thời gian hoàn chỉnh có kích thước:
\begin{equation}
    \mathbf{Y}_{\text{total}} \in \mathbb{R}^{24768 \times 608 \times 2}
\end{equation}
bao gồm đúng 24.768 bước thời gian liên tục, phủ kín 608 nút giao thông với 2 kênh lưu lượng phương tiện độc lập.

\section{Xây dựng Đồ thị Giao thông và Ma trận Kề RBF dựa trên OpenStreetMap}
\subsection{Trích xuất Cạnh Đồ thị theo Tuyến đường Thực tế}
Để đồ thị phản ánh đúng các ràng buộc vật lý của mạng lưới đường bộ thành phố (như các tuyến đường một chiều, dải phân cách cứng, cầu vượt hoặc rào cản tự nhiên như sông Sài Gòn và kênh rạch), nhóm nghiên cứu trích xuất dữ liệu bản đồ OpenStreetMap (OSM) của TP. Hồ Chí Minh.

Một cạnh có hướng $(i, j) \in \mathcal{E}$ chỉ được thiết lập nếu tồn tại một hành lang giao thông đường bộ hợp lệ cho phép phương tiện di chuyển trực tiếp từ vị trí camera $i$ đến vị trí camera $j$ theo đúng chiều quy định của luật giao thông mà không đi qua bất kỳ camera quan trắc trung gian nào khác. Tổng cộng có \textbf{2.452 cạnh có hướng hợp lệ} được thiết lập trên 608 đỉnh.

\subsection{Xác định Trọng số Ma trận Kề RBF}
Với mỗi cặp nút có kết nối hợp lệ $(i, j) \in \mathcal{E}$, khoảng cách định tuyến đường bộ thực tế ngắn nhất $d_{ij}$ (Network Routing Distance, tính bằng kilomet) được tính toán thông qua thuật toán Dijkstra trên đồ thị đường phố OSM. Trọng số liên kết $W_{ij} \in [0, 1]$ được tính toán bằng hàm cơ sở xuyên tâm (Radial Basis Function -- RBF):
\begin{equation}
W_{ij} =
\begin{cases}
\exp\left(-\frac{d_{ij}^2}{\sigma_{\mathrm{RBF}}^2}\right), & \text{nếu } (i,j) \in \mathcal{E} \text{ và } d_{ij} \le \tau, \\
0, & \text{ngược lại.}
\end{cases}
\end{equation}
trong đó:
\begin{itemize}
    \item $\sigma_{\mathrm{RBF}} = 1.0887\text{ km}$ là độ lệch chuẩn thực nghiệm của khoảng cách các cạnh kết nối hợp lệ trong mạng lưới đường phố thành phố.
    \item $\tau = 3.3945\text{ km}$ là ngưỡng bán kính cắt lọc không gian, đảm bảo tính thưa (sparsity) của ma trận kề và loại bỏ các liên kết quá xa vốn có sự tương quan giao thông mờ nhạt.
\end{itemize}

\subsection{Chuẩn hóa Toán tử Laplacian Chebyshev Có Hướng}
Ma trận trọng số $W$ được chuẩn hóa bậc hai phía nhằm bảo toàn cấu trúc luồng xe một chiều:
\begin{equation}
    L = I_N - D^{-1/2} W D^{-1/2}, \quad \tilde{L} = \frac{2}{\lambda_{\max}} L - I_N
\end{equation}
với $D_{ii} = \sum_j W_{ij}$ là ma trận bậc ra và $\lambda_{\max} \approx 2.0$ là giá trị riêng cực đại dùng để co giãn toán tử phục vụ khai triển trực giao Chebyshev.

\section{Chiến lược Phân chia Dữ liệu và Kiểm soát Rò rỉ Thông tin (Zero-Leakage Protocol)}
Để đảm bảo các kết quả đánh giá thực nghiệm có giá trị khoa học tin cậy và không bị thổi phồng do hiện tượng rò rỉ dữ liệu (Data Leakage), đề tài thiết lập các ranh giới cô lập nghiêm ngặt:

\begin{enumerate}
    \item \textbf{Phân lập Không gian cho Giai đoạn 1 (Spatial Isolation):} 5.012 ảnh được phân chia theo vị trí camera: 4.010 ảnh huấn luyện được lấy từ một tập các vị trí camera xác định; trong khi 501 ảnh Validation và 501 ảnh Test được lấy từ các camera đặt tại các ngã tư hoàn toàn khác biệt mà mô hình thị giác chưa từng nhìn thấy trong tập Train. Điều này đảm bảo đánh giá chính xác năng lực tổng quát hóa của backbone trước các góc quay mới lạ.
    \item \textbf{Phân lập Thời gian cho Giai đoạn 2 (Temporal Isolation):} Chuỗi 24.768 bước thời gian được chia tuần tự theo đúng dòng thời gian thực tế: $80\%$ đầu tiên làm tập Train (19.814 bước), $10\%$ tiếp theo làm tập Validation (2.477 bước) và $10\%$ cuối cùng làm tập Test (2.477 bước).
    \item \textbf{Đóng băng tham số chuẩn hóa Z-score:} Tuyệt đối không tính toán giá trị trung bình $\mu$ hay độ lệch chuẩn $\sigma$ trên toàn bộ tập dữ liệu. Các thông số $\mu_i$ và $\sigma_i$ của từng nút $i$ được tính toán duy nhất trên tập Train (80\% đầu) và đóng băng hoàn toàn khi áp dụng để chuẩn hóa tập Validation và Test.
\end{enumerate}

\section{Phân tích Định lượng Cấu trúc Tô pô Đồ thị Giao thông}
Để hiểu rõ các thuộc tính hình học và cấu trúc liên kết của đồ thị IC4SD-Traffic-HCM, Bảng \ref{tab:topo_dothi_chap3} tổng hợp các chỉ số tô pô đồ thị chuẩn mực được tính toán trên 608 đỉnh và 2.452 cạnh có hướng.

\begin{table}[htbp]
\centering
\caption{Các chỉ số định lượng đặc trưng tô pô đồ thị IC4SD-Traffic-HCM.}
\label{tab:topo_dothi_chap3}
\small
\resizebox{\textwidth}{!}{%
\begin{tabular}{llc}
\toprule
\textbf{Chỉ số tô pô} & \textbf{Công thức / Ký hiệu} & \textbf{Giá trị Định lượng} \\
\midrule
Số lượng đỉnh quan trắc & $N = |\mathcal{V}|$ & 608 nút camera \\
Số lượng cạnh liên kết có hướng & $|\mathcal{E}|$ & 2.452 cạnh \\
Bậc nút trung bình & $\langle k \rangle = |\mathcal{E}| / N$ & 4.033 cạnh/nút \\
Bậc vào cực đại / Bậc ra cực đại & $k_{\text{in}}^{\max} / k_{\text{out}}^{\max}$ & 8 / 9 \\
Mật độ đồ thị (Graph Density) & $\rho = \frac{|\mathcal{E}|}{N(N-1)}$ & 0.00664 (Độ thưa $99.34\%$) \\
Hệ số phân cụm trung bình & $C_{\text{avg}} = \frac{1}{N} \sum_i C_i$ & 0.2842 \\
Khoảng cách địa lý cạnh trung bình & $\bar{d}_{\text{edge}}$ & $1.412 \pm 0.684\text{ km}$ \\
Đường kính mạng lưới lớn nhất & $D_{\text{network}}$ & $38.64\text{ km}$ \\
Tham số RBF Bandwidth & $\sigma_{\mathrm{RBF}}$ & $1.0887\text{ km}$ \\
Ngưỡng bán kính cắt lọc RBF & $\tau$ & $3.3945\text{ km}$ \\
\bottomrule
\end{tabular}%
}
\end{table}

\textbf{Phân tích đặc tính cấu trúc mạng lưới:}
\begin{enumerate}
    \item \textbf{Cấu trúc Bậc nút tự nhiên:} Bậc nút trung bình $\langle k \rangle \approx 4.03$ phản ánh chính xác cấu trúc hình học của mạng lưới đường phố đô thị thực tế, nơi đa phần các giao lộ trọng điểm là các ngã tư (4 hướng lưu thông) và ngã ba (3 hướng). Bậc cực đại đạt 9 tại các vòng xoay đa nhánh lớn (như Ngã sáu Cộng Hòa, Vòng xoay Cây Gõ và Vòng xoay An Lạc).
    \item \textbf{Đặc trưng Mạng lưới Thế giới Nhỏ (Small-World Property):} Mật độ cạnh rất thưa ($\rho \approx 0.0066$) kết hợp với hệ số phân cụm tương đối cao ($C_{\text{avg}} = 0.284$) là minh chứng điển hình của mạng lưới Small-World. Điều này khẳng định rằng dù đồ thị có quy mô lớn, thông tin tắc nghẽn giao thông có thể lan truyền nhanh chóng giữa các cụm nút giao thông lân cận thông qua một số ít các trục đại lộ xương sống.
\end{enumerate}

\section{Đặc tính Chuỗi Thời gian và Biến động Ngoại sinh trong 86 Ngày}
Chuỗi thời gian liên tục 86 ngày ghi nhận đầy đủ các sắc thái vận động của giao thông TP. Hồ Chí Minh:
\begin{itemize}
    \item \textbf{Chu kỳ Ngày đêm (Diurnal Cycle):} Hình thái lưu lượng thể hiện cấu trúc "hai đỉnh nhọn" (Bimodal distribution) rất điển hình:
    \begin{itemize}
        \item \textit{Đỉnh cao điểm sáng (07:00 -- 08:30):} Lưu lượng xe máy tăng vọt từ $300\%$ đến $500\%$ so với mức trung bình ban ngày, tập trung đổ dồn từ các khu đô thị vệ tinh (TP. Thủ Đức, Quận 12, Bình Tân) về khu trung tâm Quận 1 và Quận 3.
        \item \textit{Đỉnh cao điểm chiều (17:00 -- 18:30):} Luồng phương tiện đảo chiều thoát ra các cửa ngõ thành phố với thời gian giải tỏa kéo dài hơn do người dân kết hợp mua sắm và sinh hoạt.
        \item \textit{Đáy ban đêm (00:00 -- 05:00):} Lưu lượng xe máy giảm xuống dưới $5\%$ công suất đường, trong khi tỷ lệ xe tải nặng và xe container tăng nhẹ trên các trục vành đai.
    \end{itemize}
    \item \textbf{Chu kỳ Tuần (Weekly Periodicity):} Lưu lượng các ngày trong tuần (Thứ Hai đến Thứ Sáu) có tính ổn định rất cao với hệ số tương quan chéo ngày-với-ngày đạt $\rho > 0.94$. Vào Thứ Bảy và Chủ Nhật, đỉnh cao điểm sáng biến mất, thay vào đó là sự gia tăng lưu lượng đồng đều từ 10:00 đến 21:00 tại các trục đường thương mại và giải trí.
    \item \textbf{Biến động Thời tiết Cực đoan (Triều cường và Mưa giông nhiệt đới):} Khoảng thời gian thu thập từ tháng 6 đến tháng 9 trùng với mùa mưa Nam Bộ. Bộ dữ liệu ghi nhận 18 đợt mưa lớn trên diện rộng và 4 đợt triều cường làm ngập sâu nhiều tuyến đường ven sông (đường Nguyễn Hữu Cảnh, Huỳnh Tấn Phát, Trần Xuân Soạn). Trong các thời điểm này, vận tốc dòng xe giảm mạnh, dòng xe máy dồn ứ bất thường và lưu lượng đo được tại các camera lân cận biến động dị thường, tạo thành tập mẫu thử thách khắc nghiệt cho khả năng tự thích ứng của mô hình TA-STGCN.
\end{itemize}



% Chương 4: Phương pháp nghiên cứu đề xuất
\chapter{PHƯƠNG PHÁP NGHIÊN CỨU ĐỀ XUẤT}
\label{chap:phuong_phap}

Chương này trình bày chi tiết kiến trúc kỹ thuật và thuật toán của khung làm việc hai giai đoạn (Two-Stage Perception-to-Forecasting Pipeline) kết hợp giữa thị giác máy tính và học sâu trên đồ thị không -- thời gian. Trọng tâm của chương tập trung vào thiết kế chi tiết của mô hình hồi quy đếm xe Giai đoạn 1 và mô hình đồ thị cải tiến \textbf{TA-STGCN} ở Giai đoạn 2.

\section{Kiến trúc Tổng thể Khung làm việc Hai Giai đoạn}
Mô hình tổng thể được tổ chức thành hai giai đoạn phân tách rõ ràng về mặt chức năng và không gian tính toán:
\begin{itemize}
    \item \textbf{Giai đoạn 1 -- Thu nhận Thị giác tại Camera (Vision Perception Stage):} Tiếp nhận các ảnh tĩnh độc lập từ camera giám sát tại từng ngã tư giao thông, áp dụng mạng nơ-ron tích chập sâu với đầu hồi quy MLP để ước lượng trực tiếp số lượng xe ô tô con và số lượng xe máy tức thời.
    \item \textbf{Giai đoạn 2 -- Dự báo Đồ thị Không -- Thời gian Toàn mạng lưới (Network-Wide Forecasting Stage):} Tổng hợp các vector số đếm từ toàn bộ 608 camera thành chuỗi thời gian phân tán trên đồ thị không gian có hướng, áp dụng mô hình TA-STGCN để dự báo lưu lượng đa bước trong khoảng thời gian nhìn trước 30 phút ($H=6$ bước, mỗi bước 5 phút).
\end{itemize}

\begin{figure}[htbp]
\centering
\includegraphics[width=0.98\linewidth]{fig/ta_stgcn_architecture.pdf}
\caption{Kiến trúc tổng thể của mô hình TA-STGCN. Luồng dữ liệu đi qua hai khối ST-Conv sandwich, ở giữa là tầng Multi-Head Temporal Self-Attention điều phối trọng số thời gian.}
\label{fig:kien_truc_tastgcn_chap4}
\end{figure}

\section{Giai đoạn 1: Hồi quy Đếm xe Trực tiếp theo Phân lớp}
\subsection{Mạng Xương sống (Vision Backbone) và Trích xuất Đặc trưng}
Mô hình tiếp nhận một ảnh tĩnh đầu vào $\mathbf{I} \in \mathbb{R}^{3 \times 224 \times 224}$ đã được chuẩn hóa theo kỳ vọng và độ lệch chuẩn của tập dữ liệu ImageNet. 

Nghiên cứu lựa chọn kiến trúc \textbf{EfficientNet-B5} làm mạng xương sống chính thức (dựa trên các kết quả tối ưu hóa đa mục tiêu về độ chính xác và FLOPs được chứng minh tại Chương 5). EfficientNet-B5 bao gồm 7 cụm khối tích chập MBConv (Mobile Inverted Bottleneck Convolution) kết hợp cơ chế chú ý kênh Squeeze-and-Excitation (SE). Đầu ra của tầng tích chập cuối cùng là tensor đặc trưng bậc 3:
\begin{equation}
    \mathbf{F}_{\text{conv}} \in \mathbb{R}^{2048 \times 7 \times 7}
\end{equation}
Áp dụng tầng gộp trung bình toàn cục (Global Average Pooling -- GAP) trên hai chiều không gian:
\begin{equation}
    \mathbf{f} = \text{GAP}(\mathbf{F}_{\text{conv}}) = \frac{1}{7 \times 7} \sum_{i=1}^{7} \sum_{j=1}^{7} \mathbf{F}_{\text{conv}}(:, i, j) \in \mathbb{R}^{2048}
\end{equation}

\subsection{Đầu Hồi quy MLP Đa tầng Thống nhất (Unified MLP Regression Head)}
Để đảm bảo tính so sánh công bằng giữa các kiến trúc backbone khác nhau, tầng phân loại 1.000 lớp ImageNet gốc được loại bỏ và thay thế bằng cùng một đầu hồi quy MLP đa tầng thống nhất:
\begin{align}
    \mathbf{h}_1 &= \text{LayerNorm}(\mathbf{W}_1 \mathbf{f} + \mathbf{b}_1), \quad \mathbf{W}_1 \in \mathbb{R}^{128 \times 2048}, \ \mathbf{b}_1 \in \mathbb{R}^{128} \\
    \mathbf{h}_2 &= \text{ReLU}(\mathbf{h}_1) \in \mathbb{R}^{128} \\
    \hat{\mathbf{c}} &= \mathbf{W}_2 \mathbf{h}_2 + \mathbf{b}_2, \quad \mathbf{W}_2 \in \mathbb{R}^{2 \times 128}, \ \mathbf{b}_2 \in \mathbb{R}^{2}
\end{align}
trong đó $\hat{\mathbf{c}} = [\hat{c}_{\text{car}}, \hat{c}_{\text{bike}}]^T \in \mathbb{R}^2$ là vector số lượng xe ô tô con và xe máy dự đoán.

\textbf{Lý do thiết kế kiến trúc đầu hồi quy:}
\begin{itemize}
    \item Việc sử dụng tầng ẩn 128 chiều đóng vai trò như một phễu nén đặc trưng (Bottleneck Compression), buộc mô hình phải chắt lọc các biểu diễn tinh túy nhất liên quan đến mật độ xe thay vì ghi nhớ chi tiết cảnh nền.
    \item Tầng LayerNorm được đặt trước hàm kích hoạt phi tuyến ReLU giúp ổn định phân phối giá trị kích hoạt, ngăn ngừa hiện tượng "chết nơ-ron" (Dead ReLU) khi huấn luyện trên dữ liệu ảnh có độ biến động ánh sáng lớn.
\end{itemize}

\subsection{Hàm Mất mát Smooth L1 (Huber Loss)}
Để huấn luyện mô hình Giai đoạn 1 có khả năng chống chịu bền bỉ trước các ảnh bị lóa sáng hoặc che khuất đột ngột, hàm mất mát Huber Loss ($\delta = 1.0$) được áp dụng cho từng phân lớp:
\begin{align}
    e_{\text{car}} &= c_{\text{car}} - \hat{c}_{\text{car}} \\
    e_{\text{bike}} &= c_{\text{bike}} - \hat{c}_{\text{bike}}
\end{align}
Hàm mục tiêu tối ưu hóa toàn cục của Giai đoạn 1 là:
\begin{equation}
    \mathcal{L}_{\text{Stage 1}} = \frac{1}{2} \left( \mathcal{L}_{\mathrm{Huber}}(e_{\text{car}}) + \mathcal{L}_{\mathrm{Huber}}(e_{\text{bike}}) \right)
\end{equation}
với:
\begin{equation}
\mathcal{L}_{\mathrm{Huber}}(e) =
\begin{cases}
\frac{1}{2} e^2, & \text{khi } |e| < 1.0, \\
|e| - 0.5, & \text{khi } |e| \ge 1.0.
\end{cases}
\end{equation}

\section{Giai đoạn 2: Mô hình Dự báo Đồ thị Không -- Thời gian TA-STGCN}
\subsection{Cấu trúc Tensor Đặc trưng Đa chiều}
Tại mỗi thời điểm dự báo $t$, đầu vào của mô hình là một chuỗi quan sát lịch sử gồm $T_{\text{in}} = 24$ bước thời gian (tương đương 120 phút quan sát). Tại mỗi bước thời gian $\tau \in \{t - T_{\text{in}} + 1, \dots, t\}$ và tại mỗi nút camera $i \in \{1, \dots, N\}$, một vector đặc trưng $F = 5$ chiều được cấu tạo:
\begin{enumerate}
    \item Số lượng xe ô tô con đã chuẩn hóa Z-score per-node: $\tilde{c}_{\text{car}}^{(i, \tau)}$.
    \item Số lượng xe máy đã chuẩn hóa Z-score per-node: $\tilde{c}_{\text{bike}}^{(i, \tau)}$.
    \item Đặc trưng sin của chu kỳ ngày đêm: $\sin\left(\frac{2\pi \cdot \text{tod}(\tau)}{1440}\right)$.
    \item Đặc trưng cos của chu kỳ ngày đêm: $\cos\left(\frac{2\pi \cdot \text{tod}(\tau)}{1440}\right)$.
    \item Chỉ số giờ chuẩn hóa trong ngày: $\frac{\text{hour}(\tau)}{24}$.
\end{enumerate}
với $\text{tod}(\tau)$ là số phút tính từ đầu ngày (từ 0 đến 1439). 

Tensor đầu vào cho một mini-batch gồm $B$ mẫu có kích thước:
\begin{equation}
    \mathbf{X} \in \mathbb{R}^{B \times T_{\text{in}} \times N \times 5}
\end{equation}
Mục tiêu là dự báo ma trận lưu lượng trong $H = 6$ bước thời gian kế tiếp (30 phút tới):
\begin{equation}
    \hat{\mathbf{Y}} \in \mathbb{R}^{B \times H \times N \times 2}
\end{equation}

\subsection{Chi tiết Thiết kế Khối ST-Conv (Spatio-Temporal Convolutional Block)}
Khối ST-Conv được thiết kế theo cấu trúc "Sandwich" kinh điển gồm ba tầng tuần tự:
\begin{enumerate}
    \item \textbf{Tầng Tích chập Thời gian GLU thứ nhất (Temporal GLU 1):} Tiếp nhận tensor đặc trưng, áp dụng tích chập 1D dọc theo trục thời gian với kernel $K_t = 3$ để mở rộng số kênh ẩn lên $C = 80$:
    \begin{equation}
        \mathbf{H}_1 = (\mathbf{X} \star \mathbf{\Theta}_{1,1} + \mathbf{b}_{1,1}) \odot \sigma(\mathbf{X} \star \mathbf{\Theta}_{1,2} + \mathbf{b}_{1,2}) \in \mathbb{R}^{B \times C \times N \times T_{\text{in}}}
    \end{equation}
    \item \textbf{Tầng Tích chập Không gian ChebNet (Spatial Chebyshev Graph Conv):} Áp dụng đa thức Chebyshev bậc $K = 3$ trên ma trận Laplacian $\tilde{\mathbf{L}}$ của đồ thị OpenStreetMap:
    \begin{equation}
        \mathbf{H}_2 = \text{ReLU}\left( \sum_{k=0}^{K-1} T_k(\tilde{\mathbf{L}}) \mathbf{H}_1 \mathbf{\Theta}_k^{\text{spatial}} \right) \in \mathbb{R}^{B \times C \times N \times T_{\text{in}}}
    \end{equation}
    với $\mathbf{\Theta}_k^{\text{spatial}} \in \mathbb{R}^{C \times C}$ là ma trận trọng số biến đổi kênh học được.
    \item \textbf{Tầng Tích chập Thời gian GLU thứ hai (Temporal GLU 2):} Tiếp tục tinh lọc đặc trưng thời gian sau khi đã được tổng hợp không gian:
    \begin{equation}
        \mathbf{H}_3 = (\mathbf{H}_2 \star \mathbf{\Theta}_{2,1} + \mathbf{b}_{2,1}) \odot \sigma(\mathbf{H}_2 \star \mathbf{\Theta}_{2,2} + \mathbf{b}_{2,2}) \in \mathbb{R}^{B \times C \times N \times T_{\text{in}}}
    \end{equation}
    Đầu ra của khối được kết hợp với kết nối tắt Residual, Dropout ($p = 0.25$) và chuẩn hóa LayerNorm:
    \begin{equation}
        \mathbf{H}_{\text{ST-Conv}} = \text{LayerNorm}\left( \mathbf{H}_{\text{residual}} + \text{Dropout}(\mathbf{H}_3, p=0.25) \right)
    \end{equation}
\end{enumerate}

\subsection{Mô-đun Multi-Head Temporal Self-Attention Đặt ở Giữa}
Điểm cải tiến kiến trúc cốt lõi của \textbf{TA-STGCN} nằm ở việc chèn mô-đun Chú ý tự thân theo thời gian đa đầu vào vị trí giữa hai khối ST-Conv.

\subsubsection{Cơ chế Định hình lại Tensor và Tính toán Attention}
Gọi đầu ra sau khối ST-Conv thứ nhất là $\mathbf{H}^{(1)} \in \mathbb{R}^{B \times C \times N \times T_{\text{in}}}$ với $C = 80$ channels và $T_{\text{in}} = 24$ steps. Để triệt tiêu hoàn toàn chi phí tương tác bậc hai theo số nút, tensor được hoán vị và làm phẳng:
\begin{equation}
    \mathbf{Z} = \text{Reshape}(\text{Permute}(\mathbf{H}^{(1)}, [0, 2, 3, 1])) \in \mathbb{R}^{(B \cdot N) \times T_{\text{in}} \times C}
\end{equation}
Mỗi nút trong batch giờ đây được xem như một chuỗi độc lập có độ dài $T_{\text{in}} = 24$. Áp dụng cơ chế Multi-Head Attention với $h = 4$ đầu (chiều mỗi đầu $d_k = C/h = 20$):
\begin{align}
    \mathbf{Q}_i &= \mathbf{Z} \mathbf{W}_i^Q \in \mathbb{R}^{(B \cdot N) \times T_{\text{in}} \times d_k} \\
    \mathbf{K}_i &= \mathbf{Z} \mathbf{W}_i^K \in \mathbb{R}^{(B \cdot N) \times T_{\text{in}} \times d_k} \\
    \mathbf{V}_i &= \mathbf{Z} \mathbf{W}_i^V \in \mathbb{R}^{(B \cdot N) \times T_{\text{in}} \times d_k} \quad (\forall i = 1, \dots, h) \\
    \text{head}_i &= \text{Softmax}\left(\frac{\mathbf{Q}_i \mathbf{K}_i^T}{\sqrt{d_k}}\right) \mathbf{V}_i \in \mathbb{R}^{(B \cdot N) \times T_{\text{in}} \times d_k} \\
    \text{MultiHead}(\mathbf{Z}) &= \left[ \text{head}_1 \,\|\, \text{head}_2 \,\|\, \dots \,\|\, \text{head}_h \right] \mathbf{W}^O \in \mathbb{R}^{(B \cdot N) \times T_{\text{in}} \times C}
\end{align}

\subsubsection{Khối Mạng Truyền thẳng (FFN) và Chuẩn hóa Residual}
Đầu ra sau Attention được đưa qua mạng truyền thẳng hai tầng tuyến tính với kích hoạt GELU và hệ số mở rộng kênh gấp 4 lần ($4C = 320$ chiều):
\begin{align}
    \mathbf{Z}_1 &= \text{LayerNorm}(\mathbf{Z} + \text{Dropout}(\text{MultiHead}(\mathbf{Z}), p=0.1)) \\
    \text{FFN}(\mathbf{Z}_1) &= \text{GELU}(\mathbf{Z}_1 \mathbf{W}_{\text{ffn}, 1} + \mathbf{b}_{\text{ffn}, 1}) \mathbf{W}_{\text{ffn}, 2} + \mathbf{b}_{\text{ffn}, 2} \\
    \mathbf{Z}_{\text{out}} &= \text{LayerNorm}(\mathbf{Z}_1 + \text{Dropout}(\text{FFN}(\mathbf{Z}_1), p=0.1))
\end{align}
với $\mathbf{W}_{\text{ffn}, 1} \in \mathbb{R}^{C \times 4C}$ và $\mathbf{W}_{\text{ffn}, 2} \in \mathbb{R}^{4C \times C}$.

Tensor $\mathbf{Z}_{\text{out}}$ sau đó được phục hồi kích thước $\mathbb{R}^{B \times C \times N \times T_{\text{in}}}$ và truyền vào khối ST-Conv thứ hai để tiếp tục tinh chỉnh đặc trưng trước khi xuất ra tầng dự báo.

\subsection{Tầng Chiếu Thời gian Ngõ ra và Hàm Mục tiêu Huấn luyện}
Sau khi đi qua khối ST-Conv thứ hai, tensor có kích thước $\mathbf{H}^{(2)} \in \mathbb{R}^{(B \cdot N) \times C \times T_{\text{in}}}$. Mô hình áp dụng một tầng tích chập 1D chuẩn theo chiều thời gian với kích thước kernel đúng bằng $T_{\text{in}} = 24$ để thu gọn chiều thời gian từ 24 về 1:
\begin{equation}
    \mathbf{Y}_{\text{proj}} = \text{Conv1D}(\mathbf{H}^{(2)}, \text{kernel\_size}=T_{\text{in}}) \in \mathbb{R}^{(B \cdot N) \times (H \cdot 2) \times 1}
\end{equation}
Tensor được định hình lại thành $\hat{\mathbf{Y}} \in \mathbb{R}^{B \times H \times N \times 2}$ tương ứng với dự báo 6 bước tương lai cho 2 loại xe trên toàn bộ 608 nút.

Mô hình được tối ưu hóa bằng hàm mất mát Huber Loss trên toàn bộ các nút và các bước dự báo:
\begin{equation}
    \mathcal{L}_{\text{Stage 2}} = \frac{1}{B \cdot H \cdot N \cdot 2} \sum_{b=1}^{B} \sum_{h=1}^{H} \sum_{i=1}^{N} \sum_{m=1}^{2} \mathcal{L}_{\mathrm{Huber}}\left( Y_{b,h,i,m} - \hat{Y}_{b,h,i,m} \right)
\end{equation}
với ngưỡng chuyển tiếp chuẩn $\delta = 1.0$.

\section{Thuật toán Huấn luyện và Suy luận Toàn diện}
Thuật toán \ref{algo:training_pipeline} mô tả tóm tắt toàn bộ quy trình huấn luyện hai giai đoạn từ ảnh thô camera đến mô hình dự báo đồ thị hoàn chỉnh.

\begin{algorithm_box}{Quy trình Huấn luyện Hai Giai đoạn (Two-Stage Training Pipeline)}{
\textbf{Đầu vào:} Tập ảnh camera $\mathcal{D}_{\text{vision}} = \{(I_k, \mathbf{c}_k)\}_{k=1}^{5012}$, Đồ thị đường phố OpenStreetMap $\mathcal{G}_{\text{OSM}}$, Ma trận trọng số $W$. \\
\textbf{Đầu ra:} Mô hình thị giác $f_{\theta^*}$, Mô hình đồ thị dự báo $\mathcal{M}_{\Phi^*}$. \\
\textbf{Pha 1: Huấn luyện Mô hình Thị giác (Stage 1):}
\begin{enumerate}
    \item Khởi tạo backbone EfficientNet-B5 với trọng số tiền huấn luyện ImageNet.
    \item Khởi tạo đầu Unified MLP Regression Head ngẫu nhiên.
    \item \textbf{For} epoch = 1 \textbf{to} 120 \textbf{do}:
    \begin{itemize}
        \item Lấy batch ảnh $B=32$, áp dụng Label-preserving Augmentation.
        \item Tính dự đoán $\hat{\mathbf{c}} = f_\theta(I)$ và tính hàm mất mát Huber Loss $\mathcal{L}_{\text{Stage 1}}$.
        \item Cập nhật trọng số bằng thuật toán tối ưu AdamW kết hợp Cosine Annealing.
    \end{itemize}
    \item Lưu trọng số tối ưu $\theta^*$ theo sai số MAE nhỏ nhất trên tập Validation.
\end{enumerate}
\textbf{Pha 2: Sinh Dữ liệu Chuỗi Thời gian và Tiền xử lý:}
\begin{enumerate}
    \item Sử dụng $f_{\theta^*}$ suy luận liên tục trên 608 camera trong 86 ngày để tạo chuỗi 1 phút.
    \item Áp dụng bộ lọc Causal Double-Rolling Filter và lấy mẫu đều đặn $\Delta t = 5$ phút để tạo tensor $\mathbf{Y}_{\text{total}} \in \mathbb{R}^{24768 \times 608 \times 2}$.
    \item Phân chia chuỗi thời gian: $80\%$ Train, $10\%$ Val, $10\%$ Test.
    \item Chuẩn hóa Z-score per-node đóng băng thông số $\mu_i, \sigma_i$ từ tập Train.
\end{enumerate}
\textbf{Pha 3: Huấn luyện Mô hình Đồ thị TA-STGCN (Stage 2):}
\begin{enumerate}
    \item Tính toán ma trận Laplacian chuẩn hóa $\tilde{\mathbf{L}}$ từ ma trận khoảng cách RBF.
    \item Khởi tạo mô hình TA-STGCN với 2 khối ST-Conv và Temporal Multi-Head Attention ($h=4$).
    \item \textbf{For} epoch = 1 \textbf{to} 120 \textbf{do}:
    \begin{itemize}
        \item Lấy batch chuỗi thời gian $B=64$, cửa sổ $T_{\text{in}}=24$, đích $H=6$.
        \item Lan truyền xuôi qua ST-Conv 1 $\to$ Temporal Attention $\to$ ST-Conv 2 $\to$ Conv1D Output.
        \item Tính hàm mất mát Huber Loss $\mathcal{L}_{\text{Stage 2}}$ và cập nhật trọng số bằng Adam.
        \item Điều chỉnh Learning Rate theo cơ chế ReduceLROnPlateau.
    \end{itemize}
    \item Dừng sớm sau 20 epoch nếu Validation MAE không cải thiện. Lưu trọng số tối ưu $\Phi^*$.
\end{enumerate}
}
\label{algo:training_pipeline}
\end{algorithm_box}

\section{Đặc tả Kiến trúc Chi tiết Từng Tầng của TA-STGCN}
Bảng \ref{tab:layer_spec_tastgcn_chap4} mô tả chi tiết kích thước tensor đầu vào, đầu ra, kích thước bộ lọc và số lượng tham số học được của từng tầng trong mô hình TA-STGCN với kích thước batch $B$, số lượng nút $N=608$, cửa sổ lịch sử $T_{\text{in}}=24$ và tầm nhìn dự báo $H=6$.

\begin{table}[htbp]
\centering
\caption{Đặc tả chi tiết cấu trúc từng tầng của mô hình đề xuất TA-STGCN.}
\label{tab:layer_spec_tastgcn_chap4}
\small
\resizebox{\textwidth}{!}{%
\begin{tabular}{lllcc}
\toprule
\textbf{Tên khối chức năng} & \textbf{Tầng nơ-ron thành phần} & \textbf{Kích thước Tensor đầu ra} & \textbf{Cấu hình Bộ lọc} & \textbf{Số tham số} \\
\midrule
Đầu vào & Tensor dữ liệu lịch sử & $(B, 2, 24, 608)$ & -- & 0 \\
\midrule
\multirow{4}{*}{\textbf{ST-Conv Block 1}} & 1D Causal Conv GLU 1 & $(B, 64, 22, 608)$ & $K_t = 3, C: 2 \to 64 \times 2$ & 832 \\
 & Spectral ChebNet GCN 1 & $(B, 64, 22, 608)$ & $K = 3, C: 64 \to 64$ & 12.352 \\
 & 1D Causal Conv GLU 2 & $(B, 64, 20, 608)$ & $K_t = 3, C: 64 \to 64 \times 2$ & 24.704 \\
 & LayerNorm \& Residual 1 & $(B, 64, 20, 608)$ & Affine $\gamma, \beta \in \mathbb{R}^{64}$ & 128 \\
\midrule
\multirow{4}{*}{\textbf{Temporal Attention}} & Linear Projections ($Q, K, V$) & $(B \cdot 608, 20, 64)$ & $W^Q, W^K, W^V \in \mathbb{R}^{64 \times 64}$ & 12.480 \\
 & Multi-Head Attention ($h=4$) & $(B \cdot 608, 20, 64)$ & $d_k = 16, \text{Scaling } 1/\sqrt{16}$ & 0 \\
 & Output Projection $W^O$ & $(B \cdot 608, 20, 64)$ & $W^O \in \mathbb{R}^{64 \times 64}$ & 4.160 \\
 & LayerNorm \& Residual 2 & $(B, 64, 20, 608)$ & Affine $\gamma, \beta \in \mathbb{R}^{64}$ & 128 \\
\midrule
\multirow{4}{*}{\textbf{ST-Conv Block 2}} & 1D Causal Conv GLU 3 & $(B, 64, 18, 608)$ & $K_t = 3, C: 64 \to 64 \times 2$ & 24.704 \\
 & Spectral ChebNet GCN 2 & $(B, 64, 18, 608)$ & $K = 3, C: 64 \to 64$ & 12.352 \\
 & 1D Causal Conv GLU 4 & $(B, 64, 16, 608)$ & $K_t = 3, C: 64 \to 64 \times 2$ & 24.704 \\
 & LayerNorm \& Residual 3 & $(B, 64, 16, 608)$ & Affine $\gamma, \beta \in \mathbb{R}^{64}$ & 128 \\
\midrule
\multirow{3}{*}{\textbf{Output Block}} & Temporal Conv Projection & $(B, 64, 1, 608)$ & $K_t = 16, C: 64 \to 64$ & 65.600 \\
 & Channel Conv Projection & $(B, 12, 1, 608)$ & $K_t = 1, C: 64 \to 12$ & 780 \\
 & Reshape \& De-normalization & $(B, 6, 608, 2)$ & Ánh xạ $H=6$ bước, 2 lớp & 0 \\
\midrule
\multicolumn{4}{l}{\textbf{Tổng số tham số học được (Trainable Parameters):}} & \textbf{453.250} (453K) \\
\bottomrule
\end{tabular}%
}
\end{table}

\textbf{Lý do thiết kế vị trí kẹp ở giữa (Inter-Block Placement) của Temporal Attention:}
Việc đặt tầng Multi-Head Temporal Self-Attention ở giữa hai khối ST-Conv mang lại lợi thế giải tích vượt bậc so với việc đặt ở đầu hoặc cuối mạng:
\begin{enumerate}
    \item \textbf{Khối ST-Conv 1 đóng vai trò bộ tiền trích xuất đặc trưng cục bộ (Local Spatio-Temporal Feature Extractor):} Tầng ChebNet thứ nhất đã kết hợp thông tin topo 2-hop từ các nút lân cận, giúp các vector đặc trưng đưa vào tầng Attention không còn là các giá trị điểm ảnh thô rời rạc mà đã chứa đựng ngữ cảnh cục bộ của ngã tư.
    \item \textbf{Tầng Attention thực hiện tái cân chỉnh trọng số thời gian toàn cục (Global Temporal Reweighting):} Tại vị trí này, tầng Attention tính toán ma trận tương quan giữa tất cả các bước thời gian còn lại ($T=20$), cho phép triệt tiêu các dao động nhiễu đột ngột từ cảm biến trước khi bước vào khối ST-Conv thứ hai.
    \item \textbf{Khối ST-Conv 2 thực hiện tinh chỉnh và khuếch tán không gian thứ cấp (Secondary Spatial Refinement):} Tín hiệu thời gian đã được làm mượt qua Attention tiếp tục được khuếch tán không gian một lần nữa qua đồ thị OSM, đảm bảo tính liên tục không gian -- thời gian trước khi chiếu sang các bước dự báo tương lai.
\end{enumerate}

\section{Quy trình Vận hành Suy luận Thời gian Thực (Real-Time Inference Protocol)}
Để ứng dụng mô hình vào trung tâm điều hành giao thông đô thị thực tế, quy trình suy luận được chuẩn hóa thành 5 bước liên hoàn:
\begin{enumerate}
    \item \textbf{Bước 1 -- Thu nhận và Tiền xử lý Ảnh tại Biên (Edge Ingestion):} Cứ sau mỗi chu kỳ $\Delta t = 5$ phút, các camera tại 608 ngã tư chụp ảnh hiện trường. Mạng EfficientNet-B5 lượng tử hóa INT8 thực thi suy luận trong vòng $35$ ms, trích xuất vector số lượng xe tức thời $[\hat{c}_{\text{car}}, \hat{c}_{\text{bike}}]^T$.
    \item \textbf{Bước 2 -- Cập nhật Bộ đệm Chuỗi Thời gian (FIFO Rolling Buffer):} Máy chủ trung tâm nhận 608 vector số đếm (tổng cộng $4.86\text{ KB}$), đẩy vào hàng đệm trượt First-In-First-Out (FIFO) có độ dài cố định $T_{\text{in}} = 24$ bước (tương đương 120 phút gần nhất). Điểm dữ liệu cũ nhất tại bước $t-24$ bị loại bỏ.
    \item \textbf{Bước 3 -- Chuẩn hóa Z-Score Cục bộ (Local Z-Score Scaling):} Dữ liệu trong hàng đệm được chuẩn hóa tức thì bằng các tham số $\mu_i, \sigma_i$ tương ứng của từng nút: $x_{i, t} = (y_{i, t} - \mu_i) / \sigma_i$.
    \item \textbf{Bước 4 -- Dự báo Toàn mạng lưới qua TA-STGCN (Graph Inference):} Tensor kích thước $(1, 2, 24, 608)$ được đưa qua mô hình TA-STGCN trên GPU máy chủ trung tâm. Thời gian thực thi suy luận toàn mạng lưới chỉ mất đúng \textbf{$4.76$ mili-giây}.
    \item \textbf{Bước 5 -- Giải chuẩn hóa và Phát tín hiệu Điều khiển (Actuation \& Alerting):} Đầu ra được giải chuẩn hóa về số xe thực tế: $\hat{y}_{i, t+h} = \hat{z}_{i, t+h} \cdot \sigma_i + \mu_i$. Các vector dự báo 6 bước tương lai ($t+5$ đến $t+30$ phút) được truyền qua API thời gian thực đến bộ điều khiển chu kỳ đèn tín hiệu thích ứng tại từng ngã tư và đẩy lên ứng dụng bản đồ số cảnh báo người dân.
\end{enumerate}



% Chương 5: Thực nghiệm, kết quả và đánh giá
\chapter{Thực nghiệm, kết quả và đánh giá}
\label{chap:thuc_nghiem}

Chương này trình bày toàn diện các kết quả thực nghiệm, đối sánh định lượng và kiểm định thống kê chuẩn mực cho cả hai giai đoạn của hệ thống Two-Stage Pipeline trên bộ dữ liệu thực tế IC4SD-Traffic-HCM: Giai đoạn 1 (ước lượng số lượng phương tiện từ ảnh thị giác) và Giai đoạn 2 (dự báo lưu lượng không -- thời gian trên mạng lưới đồ thị 608 nút camera). Trọng tâm phân tích được đặt vào cơ chế \textbf{lan truyền sai số nhận dạng (Error Propagation)} và \textbf{năng lực kháng nhiễu cảm biến (Noise Robustness)} giữa các dòng kiến trúc học sâu không -- thời gian.

% ==============================================================================
\section{Thiết lập môi trường và tiêu chí đánh giá thực nghiệm}
\label{sec:thiet_lap_moi_truong}

\subsection{Môi trường phần cứng và phần mềm}
Các mô hình trong nghiên cứu được cài đặt trên nền tảng PyTorch 2.x và PyTorch Geometric. Quá trình huấn luyện và đánh giá được thực thi trên hệ thống máy chủ tính toán hiệu năng cao trang bị \textbf{2 GPU NVIDIA RTX 4090 (24GB VRAM mỗi GPU)}, bộ vi xử lý AMD EPYC 7543 32 lõi và 128GB RAM DDR4. Quá trình huấn luyện mô hình đồ thị áp dụng cơ chế song song hóa dữ liệu (DataParallel) phân bổ đều 32 mẫu trên mỗi GPU (tổng batch size $B=64$).

\subsection{Đặc tả cấu trúc tham số và bộ nhớ kích hoạt GPU của các mô hình Stage 2}
Bảng \ref{tab:sieu_tham_so_chap5_models} đối sánh cấu trúc tham số, cơ chế xử lý không gian -- thời gian và mức tiêu thụ bộ nhớ kích hoạt GPU đỉnh (Peak GPU Activation Memory) của 8 mô hình dự báo không -- thời gian được đánh giá trong Giai đoạn 2.

\begin{table}[htbp]
\centering
\caption{Đối sánh kiến trúc, tham số và bộ nhớ GPU đỉnh trong Giai đoạn 2 ($N=608$ nút, $T_{\text{in}}=24, H=6$).}
\label{tab:sieu_tham_so_chap5_models}
\small
\resizebox{\textwidth}{!}{%
\begin{tabular}{lcccccc}
\toprule
\textbf{Kiến trúc mô hình} & \textbf{Số tham số} & \textbf{RAM đỉnh (MB)} & \textbf{Toán tử Không gian} & \textbf{Toán tử Thời gian} & \textbf{Số đầu ($h$)} & \textbf{Số khối} \\
\midrule
STGCN Baseline \cite{stgcn} & 375.212 & 4.837,6 & ChebNet ($K=3$) & 1D Gated GLU & 0 & 2 Khối \\
STGCN (Large) & 487.420 & 5.793,9 & ChebNet ($K=3$) & 1D Gated GLU & 0 & 2 Khối \\
Graph WaveNet \cite{wavenet} & 548.620 & 7.731,4 & Adaptive GCN & Dilated Causal Conv & 0 & 4 Khối \\
ASTGCN \cite{astgcn} & 1.532.566 & 9.260,7 & ChebNet + Spatial Attn & 1D Conv + Temporal Attn & Không/Thời gian & 2 Khối \\
STAEformer \cite{staeformer} & 568.492 & 11.509,1 & Adaptive Spatial Emb & Spatio-Temporal Trans & $h=4$ & 1 Khối \\
MegaCRN \cite{megacrn} & 527.402 & 4.850,4 & Meta-Graph Emb & GRU Enc-Dec & 0 & GRU Layers \\
DSTAGNN \cite{dstagnn} & 910.796 & 9.801,5 & Dyn. Softmax Graph & Temporal Transformer & $h=4$ & 1 Khối \\
iTransformer \cite{itransformer} & \textbf{149.516} & \textbf{265,2} & Variate Attention & Inverted MLP & $h=8$ & 2 Tầng \\
\midrule
\textbf{TA-STGCN (Đề xuất)} & \textbf{453.052} & \textbf{8.606,9} & ChebNet ($K=3$) & 1D GLU + Temporal Self-Attn & $h=4$ & 2 Khối \\
\bottomrule
\end{tabular}%
}
\end{table}

\subsection{Cấu hình siêu tham số huấn luyện}
Bảng \ref{tab:sieu_tham_so_chap5} tổng hợp cấu hình siêu tham số huấn luyện chuẩn mực được áp dụng đồng nhất cho Giai đoạn 1 và Giai đoạn 2.

\begin{table}[htbp]
\centering
\caption{Cấu hình siêu tham số huấn luyện chi tiết cho Giai đoạn 1 và Giai đoạn 2.}
\label{tab:sieu_tham_so_chap5}
\small
\begin{tabularx}{\textwidth}{l X X}
\toprule
\textbf{Siêu tham số} & \textbf{Giai đoạn 1 (Nhận dạng thị giác)} & \textbf{Giai đoạn 2 (Dự báo đồ thị)} \\
\midrule
Thuật toán tối ưu hóa & AdamW & Adam \\
Tốc độ học khởi tạo ($\eta$) & $3 \times 10^{-4}$ & $1 \times 10^{-3}$ \\
Hệ số suy giảm trọng số (Weight Decay) & $1 \times 10^{-2}$ & $1 \times 10^{-4}$ \\
Chiến lược điều chỉnh tốc độ học & Cosine Annealing ($\eta_{\min} = 10^{-6}$) & ReduceLROnPlateau ($\gamma=0.5$, patience=10) \\
Số epoch tối đa & 120 epochs & 120 epochs \\
Ngưỡng dừng sớm (Early Stopping) & 20 epochs & 20 epochs \\
Kích thước mini-batch & 32 & 64 (32 mẫu / GPU) \\
Hàm mất mát & Smooth $L_1$ (Huber, $\delta=1.0$) & Huber Loss ($\delta=1.0$, dùng chung) \\
Tăng cường dữ liệu & Lật ngang, Color Jitter, Random Erase & Không (Bảo toàn trật tự chuỗi thời gian) \\
\bottomrule
\end{tabularx}
\end{table}

\subsection{Tiêu chí đánh giá và quy chuẩn báo cáo}
Để phản ánh trung thực năng lực vận hành thực tế, toàn bộ các chỉ số sai số được tính toán trên \textbf{giá trị số lượng phương tiện thực tế sau khi đã giải chuẩn hóa} thông qua biến đổi nghịch đảo: $y_{\text{real}} = y_{\text{z}} \cdot \sigma_i + \mu_i$, trong đó $\mu_i$ và $\sigma_i$ là kỳ vọng và độ lệch chuẩn của từng nút tính toán cố định trên tập huấn luyện (Train).

Ba thước đo sai số chuẩn mực được sử dụng gồm:
\begin{align}
    \text{MAE} &= \frac{1}{N \cdot H \cdot M} \sum_{i=1}^{N} \sum_{h=1}^{H} \sum_{k=1}^{M} \left| y_{i,h}^{(k)} - \hat{y}_{i,h}^{(k)} \right| \\
    \text{RMSE} &= \sqrt{\frac{1}{N \cdot H \cdot M} \sum_{i=1}^{N} \sum_{h=1}^{H} \sum_{k=1}^{M} \left( y_{i,h}^{(k)} - \hat{y}_{i,h}^{(k)} \right)^2} \\
    \text{MAPE} &= \frac{1}{|\Omega|} \sum_{(i,h,k) \in \Omega} \frac{\left| y_{i,h}^{(k)} - \hat{y}_{i,h}^{(k)} \right|}{y_{i,h}^{(k)} + \epsilon_{\mathrm{MAPE}}} \times 100\%
\end{align}
trong đó $N = 608$ nút mạng, $H = 6$ bước thời gian dự báo, $M$ là số mẫu chuỗi kiểm thử, $\Omega = \{(i,h,k) \mid y_{i,h}^{(k)} > 0.5\}$ là tập các điểm quan trắc có dòng xe thực tế lưu thông nhằm tránh sai số vô nghĩa khi lưu lượng bằng 0, và $\epsilon_{\mathrm{MAPE}} = 10^{-5}$ là hằng số chống chia cho 0.

Toàn bộ các thí nghiệm được lặp lại độc lập trên \textbf{5 seed ngẫu nhiên} $\mathcal{S} = \{42, 100, 2024, 22, 99\}$. Kết quả báo cáo theo định dạng thống kê chuẩn mực $\text{Mean} \pm \text{Std}$ cùng với khoảng tin cậy Bootstrap 95\% ($B = 1.000$ lần lấy mẫu lại có hoàn lại).

% ==============================================================================
\section{Kết quả thực nghiệm Giai đoạn 1 (Vision Perception Benchmark)}
\label{sec:thuc_nghiem_giai_doan_1}

\subsection{Đối sánh định lượng các kiến trúc mạng xương sống}
Bảng \ref{tab:stage1_benchmark_full_chap5} trình bày kết quả đánh giá toàn diện 5 kiến trúc thị giác máy tính tiêu biểu trên 501 ảnh kiểm thử độc lập từ camera giao thông TP. Hồ Chí Minh qua 5 seed ngẫu nhiên.

\begin{table}[htbp]
\centering
\caption{Đánh giá toàn diện các mô hình đếm xe Stage 1 qua 5 seed ngẫu nhiên ($\text{Mean} \pm \text{Std}$).}
\label{tab:stage1_benchmark_full_chap5}
\small
\resizebox{\textwidth}{!}{%
\begin{tabular}{lccccc}
\toprule
\textbf{Chỉ số đánh giá} & \textbf{EfficientNet-B5} & \textbf{ResNet-50} & \textbf{MobileNetV3-Large} & \textbf{ViT-Base/16} & \textbf{ConvNeXt-Tiny} \\
\midrule
Số tham số (Params) & 28.3M & 23.5M & \textbf{4.2M} & 85.8M & 27.8M \\
Độ phức tạp FLOPs & 75.3 GFLOPs & 132.2 GFLOPs & \textbf{6.9 GFLOPs} & 539.1 GFLOPs & 142.6 GFLOPs \\
Độ trễ suy luận (ms/batch 32) & $35.2 \pm 1.0$ & $34.6 \pm 0.7$ & $\mathbf{34.3 \pm 0.5}$ & $45.0 \pm 0.8$ & $34.9 \pm 0.9$ \\
\midrule
\textbf{MAE Tổng thể} & $\mathbf{3.6363 \pm 0.0619}$ & $3.6840 \pm 0.0534$ & $4.0126 \pm 0.0578$ & $5.5783 \pm 0.3120$ & $6.2903 \pm 3.6079$ \\
95\% Bootstrap CI MAE & $[3.5595 - 3.7131]$ & $[3.6177 - 3.7503]$ & $[3.9408 - 4.0844]$ & $[5.1908 - 5.9657]$ & $[1.8105 - 10.7701]$ \\
MAPE Tổng thể (\%) & $\mathbf{43.84\% \pm 1.20\%}$ & $45.83\% \pm 0.26\%$ & $52.05\% \pm 2.21\%$ & $68.79\% \pm 6.09\%$ & $80.11\% \pm 48.03\%$ \\
RMSE Tổng thể & $\mathbf{5.5345 \pm 0.0781}$ & $5.5745 \pm 0.0705$ & $6.0903 \pm 0.0756$ & $8.3289 \pm 0.5077$ & $9.6716 \pm 5.6654$ \\
\midrule
MAE Xe ô tô con & $\mathbf{2.1723 \pm 0.0415}$ & $2.2240 \pm 0.0642$ & $2.4679 \pm 0.0433$ & $3.3948 \pm 0.2252$ & $3.3374 \pm 1.5211$ \\
MAE Xe máy & $\mathbf{5.1003 \pm 0.0969}$ & $5.1440 \pm 0.0448$ & $5.5572 \pm 0.1267$ & $7.7617 \pm 0.4021$ & $9.2432 \pm 5.6963$ \\
\midrule
Hệ số biến thiên phương sai (CV) & $1.70\%$ & $1.45\%$ & $\mathbf{1.44\%}$ & $5.59\%$ & $57.36\%$ \\
\bottomrule
\end{tabular}%
}
\end{table}

\begin{figure}[htbp]
\centering
\includegraphics[width=0.88\linewidth]{fig/stage1_density_mae.pdf}
\caption{So sánh sai số MAE của các kiến trúc thị giác Stage 1 theo từng phân tầng mật độ phương tiện.}
\label{fig:stage1_density_mae_chap5}
\end{figure}

\subsection{Phân tích phương sai khởi tạo và độ ổn định tối ưu hóa}
Bảng \ref{tab:stage1_seed_variance} chi tiết hóa phương sai khởi tạo ($\sigma^2$), độ lệch chuẩn ($\sigma$), khoảng tin cậy Bootstrap 95\% và hệ số biến thiên ($\text{CV} = \sigma / \mu \times 100\%$) của chỉ số MAE tổng thể trên 5 seed ngẫu nhiên.

\begin{table}[htbp]
\centering
\caption{Thống kê phương sai và độ ổn định khởi tạo của các mô hình Stage 1 qua 5 seed ngẫu nhiên.}
\label{tab:stage1_seed_variance}
\small
\begin{tabular}{lcccc}
\toprule
\textbf{Kiến trúc Backbone} & \textbf{Mean MAE} & \textbf{95\% Bootstrap CI} & \textbf{Phương sai ($\sigma^2$)} & \textbf{CV (\%)} \\
\midrule
MobileNetV3-Large \cite{mobilenet} & $4.0126$ & $3.9408 - 4.0844$ & $3.34 \times 10^{-3}$ & $\mathbf{1.44\%}$ \\
ResNet-50 \cite{resnet} & $3.6840$ & $3.6177 - 3.7503$ & $2.85 \times 10^{-3}$ & $1.45\%$ \\
EfficientNet-B5 \cite{efficientnet} & $\mathbf{3.6363}$ & $\mathbf{3.5595 - 3.7131}$ & $3.83 \times 10^{-3}$ & $1.70\%$ \\
ViT-Base/16 \cite{vit} & $5.5783$ & $5.1908 - 5.9657$ & $9.74 \times 10^{-2}$ & $5.59\%$ \\
ConvNeXt-Tiny \cite{convnext} & $6.2903$ & $1.8105 - 10.7701$ & $1.30 \times 10^{1}$ & $57.36\%$ \\
\bottomrule
\end{tabular}
\end{table}

\textbf{Luận giải khoa học từ kết quả Giai đoạn 1:}
\begin{enumerate}
    \item \textbf{Hiệu năng vượt trội và tính kinh tế của EfficientNet-B5:} EfficientNet-B5 đạt độ chính xác cao nhất trên mọi tiêu chí (MAE tổng thể đạt $3.6363 \pm 0.0619$, MAE ô tô đạt $2.1723$, MAE xe máy đạt $5.1003$). Đáng chú ý, EfficientNet-B5 chỉ tiêu tốn 75.3 GFLOPs, tiết kiệm $43.0\%$ lượng phép tính so với ResNet-50 (132.2 GFLOPs) trong khi đạt độ trễ suy luận tương đương ($35.2$ ms so với $34.6$ ms trên batch 32). Kết quả này khẳng định tính ưu việt của nguyên lý mở rộng đồng thời (Compound Scaling), xác lập EfficientNet-B5 là mạng xương sống phù hợp nhất để sinh nhãn dữ liệu cho Giai đoạn 2.
    \item \textbf{Giải pháp siêu nhẹ cho phần cứng biên (MobileNetV3-Large):} MobileNetV3-Large thể hiện tính kinh tế ấn tượng: chỉ với 4.2M tham số và 6.9 GFLOPs (nhẹ hơn 6.7 lần về tham số và gần 11 lần về FLOPs so với EfficientNet-B5), mô hình vẫn đạt sai số chấp nhận được là $\text{MAE} = 4.0126 \pm 0.0578$ với độ ổn định khởi tạo cao nhất ($\text{CV} = 1.44\%$). Đây là cấu hình thay thế khả thi cho các cụm camera biên có kinh phí phần cứng hạn hẹp.
    \item \textbf{Độ nhạy cảm của các kiến trúc thuần Attention khi tinh chỉnh trên dữ liệu vừa:} Các kiến trúc CNN truyền thống tích hợp Batch Normalization (MobileNetV3, ResNet-50, EfficientNet-B5) duy trì tính hội tụ cực kỳ ổn định với phương sai rất nhỏ ($\sigma^2 < 3.84 \times 10^{-3}$, $\text{CV} \le 1.70\%$). Ngược lại, ViT-Base/16 có phương sai lớn hơn ($\text{CV} = 5.59\%$), trong khi ConvNeXt-Tiny thể hiện sự dao động mạnh giữa các lần chạy ngẫu nhiên ($\text{CV} = 57.36\%, \sigma^2 = 13.017$). Điều này phản ánh rằng các kiến trúc sử dụng Layer Normalization và kernel tích chập lớn $7 \times 7$ nhạy cảm hơn với trọng số khởi tạo khi thích nghi vào bài toán hồi quy liên tục trên tập dữ liệu kích thước vừa.
    \item \textbf{Kiểm định ý nghĩa thống kê:} Kiểm định phi tham số Friedman omnibus cho thấy sự phân hóa hiệu năng giữa 5 mạng xương sống có ý nghĩa thống kê cực kỳ rõ rệt ($\chi^2 = 863.55, p < 0.001$). Kiểm định Wilcoxon Signed-Rank hậu nghiệm kèm hiệu chỉnh Holm-Bonferroni xác nhận rằng EfficientNet-B5 vượt trội có ý nghĩa thống kê so với ViT-Base/16 ($p < 0.001$, Cohen's $d = 0.522$), ConvNeXt-Tiny ($p < 0.001$, Cohen's $d = 0.447$) và MobileNetV3-Large ($p < 0.001$, Cohen's $d = 0.164$). Khoảng cách sai số giữa EfficientNet-B5 và ResNet-50 không có ý nghĩa thống kê ($p = 0.253$, Cohen's $d = -0.022$), khẳng định cả hai cấu trúc CNN đều đạt độ chính xác hàng đầu trong bài toán nhận dạng.
\end{enumerate}

% ==============================================================================
\section{Phân tích sai số nhận dạng theo mật độ giao thông và kiểm chứng Grad-CAM++}
\label{sec:phan_tich_sai_so_theo_mat_do}

\subsection{Độ chệch nhận dạng phân tầng theo mật độ}
Để đánh giá tác động của mật độ phương tiện tới sai số thị giác, tập kiểm thử gồm 501 ảnh được phân tầng thành ba mức độ hoạt động: Mật độ Thấp ($<10$ xe/ảnh, 103 ảnh), Mật độ Trung bình ($10-25$ xe/ảnh, 151 ảnh) và Mật độ Cao ($>25$ xe/ảnh, 247 ảnh). Bảng \ref{tab:stage1_density_bias_chap5} chi tiết hóa độ chệch dự đoán ($\text{Bias} = \hat{y} - y$) theo phân lớp phương tiện trên từng phân tầng.

\begin{table}[htbp]
\centering
\caption{Độ chệch nhận dạng Stage 1 phân tầng theo mật độ giao thông ($\text{Mean} \pm \text{Std}$).}
\label{tab:stage1_density_bias_chap5}
\small
\resizebox{\textwidth}{!}{%
\begin{tabular}{llccc}
\toprule
\textbf{Kiến trúc Backbone} & \textbf{Phân loại} & \textbf{Mật độ Thấp ($<10$)} & \textbf{Mật độ TB ($10-25$)} & \textbf{Mật độ Cao ($>25$)} \\
\midrule
\multirow{2}{*}{\textbf{EfficientNet-B5}} & Ô tô con & $-0.03 \pm 0.24$ & $-0.30 \pm 0.30$ & $\mathbf{-0.24 \pm 0.50}$ \\
 & Xe máy & $+0.51 \pm 0.74$ & $+0.31 \pm 0.56$ & $\mathbf{-0.63 \pm 0.87}$ \\
\midrule
\multirow{2}{*}{ResNet-50} & Ô tô con & $+0.14 \pm 0.17$ & $-0.05 \pm 0.11$ & $-0.13 \pm 0.42$ \\
 & Xe máy & $+0.43 \pm 0.53$ & $+0.12 \pm 0.54$ & $-0.70 \pm 1.20$ \\
\midrule
\multirow{2}{*}{MobileNetV3-Large} & Ô tô con & $-0.03 \pm 0.09$ & $-0.33 \pm 0.23$ & $-0.37 \pm 0.74$ \\
 & Xe máy & $+0.56 \pm 0.31$ & $+0.14 \pm 0.56$ & $-0.85 \pm 1.66$ \\
\midrule
\multirow{2}{*}{ViT-Base/16} & Ô tô con & $+0.51 \pm 0.32$ & $-0.30 \pm 0.25$ & $-1.20 \pm 0.48$ \\
 & Xe máy & $+1.97 \pm 1.00$ & $+0.99 \pm 0.93$ & $-4.63 \pm 0.80$ \\
\midrule
\multirow{2}{*}{ConvNeXt-Tiny} & Ô tô con & $+1.31 \pm 1.54$ & $-0.26 \pm 0.32$ & $-2.63 \pm 3.22$ \\
 & Xe máy & $+5.57 \pm 6.91$ & $+1.80 \pm 2.36$ & $-10.22 \pm 13.11$ \\
\bottomrule
\end{tabular}%
}
\end{table}

\textbf{Phân tích hành vi độ chệch:}
\begin{itemize}
    \item Tại mật độ thấp ($<10$ xe), tất cả các mô hình xuất hiện độ chệch dương nhẹ đối với xe máy ($+0.43$ đến $+0.56$), do mạng thỉnh thoảng phản ứng nhầm với các chi tiết biển quảng cáo hoặc vệt sơn trên mặt đường.
    \item Khi mật độ tăng lên mức cao ($>25$ xe), \textbf{tất cả các kiến trúc đều chuyển sang trạng thái chệch âm (đếm thiếu phương tiện)} đối với xe máy do sự che khuất tương hỗ nghiêm trọng trong các dòng xe máy chen chúc.
    \item EfficientNet-B5 kiểm soát độ chệch âm tốt nhất (chỉ chệch $-0.63 \pm 0.87$ xe máy), trong khi ViT-Base/16 đếm thiếu trung bình $-4.63$ xe và ConvNeXt-Tiny đếm thiếu tới $-10.22$ xe. Đây chính là bằng chứng thực nghiệm giải thích nguồn gốc của sai số cảm biến mang tính hệ thống truyền vào Giai đoạn 2.
\end{itemize}

\subsection{Kiểm chứng trực quan hóa vùng chú ý bằng Grad-CAM++}
\label{sec:truc_quan_hoa_gradcam}
Để xác minh mô hình hồi quy đếm xe thực sự dựa trên các thuộc tính của phương tiện thay vì bám vào các tương quan giả tạo của cảnh nền, Hình \ref{fig:gradcam_chap5} minh họa bản đồ phân bổ đặc trưng Grad-CAM++ trích xuất từ tầng tích chập cuối của EfficientNet-B5.

\begin{figure}[htbp]
\centering
\includegraphics[width=0.92\linewidth]{fig/gradcam++.png}
\caption{Trực quan hóa bản đồ nhiệt Grad-CAM++ trên các khung hình camera giao thông, xác nhận vùng kích hoạt tập trung chuẩn xác vào thân xe ô tô và các luồng xe máy.}
\label{fig:gradcam_chap5}
\end{figure}

Bản đồ nhiệt Grad-CAM++ cho thấy các vùng kích hoạt cường độ cao (vùng màu đỏ và vàng) tập trung chính xác vào thân xe ô tô, mũ bảo hiểm và cụm đầu xe máy đang lưu thông trên làn đường. Các khu vực nền tĩnh như vỉa hè, biển báo, tòa nhà và mặt đường không có xe hoàn toàn không phát sinh kích hoạt. Điều này khẳng định mạng nơ-ron đã học được biểu diễn ngữ nghĩa thực chất của mật độ phương tiện.

% ==============================================================================
\section{Kết quả thực nghiệm Giai đoạn 2 (Graph Forecasting Benchmark)}
\label{sec:thuc_nghiem_giai_doan_2}

\subsection{Đối sánh đa chiều trên đồ thị 608 nút giao thông}
Bảng \ref{tab:stage2_benchmark_chap5} trình bày hiệu năng dự báo đa bước thời gian trung bình qua 5 seed ngẫu nhiên độc lập ($\mathcal{S} = \{42, 100, 2024, 22, 99\}$) trên toàn bộ 608 nút camera của mạng lưới đường bộ đô thị. Các mô hình được so sánh tại các mốc thời gian đại diện: $t+1$ (5 phút), $t+3$ (15 phút) và $t+6$ (30 phút phía trước).

\begin{table}[htbp]
\centering
\caption{Hiệu năng dự báo dòng giao thông Stage 2 qua 5 seed ngẫu nhiên ($\text{Mean} \pm \text{Std}$).}
\label{tab:stage2_benchmark_chap5}
\small
\resizebox{\textwidth}{!}{%
\begin{tabular}{lccccccc}
\toprule
\textbf{Kiến trúc mô hình} & \textbf{Số tham số} & \textbf{MAE Tổng} & \textbf{MAPE (\%)} & \textbf{RMSE} & \textbf{MAE ($t+1$)} & \textbf{MAE ($t+3$)} & \textbf{MAE ($t+6$)} \\
\midrule
STGCN Baseline \cite{stgcn} & 375K & $3.2703 \pm 0.0128$ & $11.44\% \pm 0.08\%$ & $4.4245 \pm 0.0194$ & $2.8430 \pm 0.0141$ & $3.2787 \pm 0.0117$ & $3.5350 \pm 0.0250$ \\
STGCN (Large) & 487K & $3.2805 \pm 0.0134$ & $11.53\% \pm 0.11\%$ & $4.4354 \pm 0.0156$ & $2.8498 \pm 0.0140$ & $3.2861 \pm 0.0162$ & $3.5497 \pm 0.0311$ \\
Graph WaveNet \cite{wavenet} & 549K & $3.2776 \pm 0.0049$ & $11.43\% \pm 0.03\%$ & $4.4486 \pm 0.0087$ & $\mathbf{2.8142 \pm 0.0088}$ & $3.2727 \pm 0.0048$ & $3.5800 \pm 0.0033$ \\
ASTGCN \cite{astgcn} & 1.533K & $3.4218 \pm 0.0115$ & $11.78\% \pm 0.08\%$ & $4.6509 \pm 0.0101$ & $2.9876 \pm 0.0262$ & $3.4296 \pm 0.0145$ & $3.6984 \pm 0.0104$ \\
STAEformer \cite{staeformer} & 568K & $3.2953 \pm 0.0184$ & $11.51\% \pm 0.12\%$ & $4.4555 \pm 0.0192$ & $2.8324 \pm 0.0237$ & $3.2924 \pm 0.0206$ & $3.5937 \pm 0.0275$ \\
MegaCRN \cite{megacrn} & 527K & $3.3433 \pm 0.0107$ & $11.55\% \pm 0.09\%$ & $4.5580 \pm 0.0182$ & $2.8436 \pm 0.0048$ & $3.3434 \pm 0.0111$ & $3.6660 \pm 0.0160$ \\
DSTAGNN \cite{dstagnn} & 911K & $3.3457 \pm 0.0169$ & $11.69\% \pm 0.08\%$ & $4.5499 \pm 0.0185$ & $2.8447 \pm 0.0162$ & $3.3373 \pm 0.0179$ & $3.6729 \pm 0.0189$ \\
iTransformer \cite{itransformer} & \textbf{150K} & $3.3720 \pm 0.0143$ & $11.65\% \pm 0.08\%$ & $4.5832 \pm 0.0216$ & $2.8613 \pm 0.0153$ & $3.3613 \pm 0.0163$ & $3.7205 \pm 0.0154$ \\
\midrule
\textbf{TA-STGCN (Đề xuất)} & \textbf{453K} & $\mathbf{3.2401 \pm 0.0108}$ & $\mathbf{11.32\% \pm 0.04\%}$ & $\mathbf{4.3920 \pm 0.0140}$ & $2.8346 \pm 0.0090$ & $\mathbf{3.2430 \pm 0.0106}$ & $\mathbf{3.4947 \pm 0.0153}$ \\
\bottomrule
\end{tabular}%
}
\end{table}

\begin{figure}[htbp]
\centering
\includegraphics[width=0.90\linewidth]{fig/val_mae_benchmark.pdf}
\caption{Đường cong hội tụ sai số Validation MAE qua các epoch huấn luyện trung bình qua 5 seed ngẫu nhiên. TA-STGCN thể hiện tiến trình hội tụ mượt mà và phương sai giữa các lần chạy rất thấp.}
\label{fig:val_mae_curves_chap5}
\end{figure}

\begin{figure}[htbp]
\centering
\includegraphics[width=0.92\linewidth]{fig/error_by_horizon.pdf}
\caption{Xu hướng tích lũy sai số (MAE và MAPE) theo tầm nhìn dự báo từ 5 phút ($t+1$) đến 30 phút ($t+6$). TA-STGCN duy trì tốc độ tích lũy sai số thoải nhất trên tầm nhìn xa.}
\label{fig:error_growth_chap5}
\end{figure}

\subsection{Đánh giá khách quan trên luồng dữ liệu sạch}
\begin{enumerate}
    \item \textbf{Độ chính xác tương đương trên dữ liệu lý tưởng:} Trên luồng dữ liệu không bị can thiệp nhiễu nhân tạo, TA-STGCN đạt sai số điểm tốt nhất ($\text{MAE} = 3.2401 \pm 0.0108, \text{MAPE} = 11.32\% \pm 0.04\%$), nhỉnh hơn Graph WaveNet ($3.2776$) và STGCN Baseline ($3.2703$) một khoảng nhỏ ($\approx 0.03 - 0.04$ MAE). Kích thước hiệu ứng Cohen's $d$ ở mức không đáng kể ($d = -0.024$ so với Graph WaveNet và $d = -0.032$ so với STGCN). Nghiên cứu khẳng định một cách trung thực rằng: \textbf{Dưới điều kiện dữ liệu sạch lý tưởng, các mô hình học sâu đồ thị hàng đầu đạt hiệu năng dự báo tương đương nhau.}
    \item \textbf{Hạn chế độ trôi dạt sai số ở tầm nhìn dài hạn ($t+6$):} Khi mở rộng tầm nhìn dự báo lên 30 phút ($t+6$), sự khác biệt kiến trúc trở nên rõ rệt hơn (Hình \ref{fig:error_growth_chap5}). Tại bước $t+6$, TA-STGCN đạt $\text{MAE} = 3.4947$, giảm được $0.0853$ MAE so với Graph WaveNet ($3.5800$) và $0.0403$ MAE so với STGCN ($3.5350$). Cơ chế chú ý thời gian đa đầu giúp kìm hãm sự tích lũy sai số trôi dạt dọc theo cửa sổ chuỗi thời gian.
    \item \textbf{Kiểm chứng dung lượng tham số:} Khi mở rộng dung lượng của STGCN gốc lên 487K tham số (phiên bản STGCN Large, cao hơn mức 453K của TA-STGCN), sai số dự báo không những không giảm mà tăng nhẹ lên $3.2805 \pm 0.0134$. Điều này chứng minh rằng ưu thế của TA-STGCN hoàn toàn bắt nguồn từ thiên kiến quy nạp (inductive bias) của cơ chế Self-Attention theo thời gian chứ không phải do việc gia tăng dung lượng tham số thuần túy.
\end{enumerate}

% ==============================================================================
\section{Phân tích lan truyền sai số theo mật độ ở Giai đoạn 2}
\label{sec:lan_truyen_sai_so_giai_doan_2}

Bảng \ref{tab:stage2_density_chap5} chi tiết hóa sai số dự báo tổng thể và độ chệch theo từng phân lớp phương tiện của cả 8 mô hình Stage 2 trên ba phân tầng mật độ giao thông thực tế.

\begin{table}[htbp]
\centering
\caption{Sai số dự báo và độ chệch phân tách ở Stage 2 theo các mức mật độ giao thông ($\text{Mean} \pm \text{Std}$).}
\label{tab:stage2_density_chap5}
\small
\resizebox{\textwidth}{!}{%
\begin{tabular}{llccc}
\toprule
\textbf{Kiến trúc mô hình} & \textbf{Chỉ số đánh giá} & \textbf{Mật độ Thấp ($<10$)} & \textbf{Mật độ TB ($10-25$)} & \textbf{Mật độ Cao ($>25$)} \\
\midrule
\multirow{3}{*}{\textbf{TA-STGCN (Đề xuất)}} & MAE Tổng thể & $0.9737 \pm 0.0066$ & $\mathbf{2.0024 \pm 0.0054}$ & $\mathbf{3.4837 \pm 0.0120}$ \\
 & Độ chệch Ô tô con & $-0.11 \pm 0.02$ & $+0.01 \pm 0.01$ & $\mathbf{-0.01 \pm 0.02}$ \\
 & Độ chệch Xe máy & $-0.07 \pm 0.01$ & $+0.01 \pm 0.02$ & $\mathbf{-0.03 \pm 0.06}$ \\
\midrule
\multirow{3}{*}{STGCN Baseline \cite{stgcn}} & MAE Tổng thể & $\mathbf{0.9696 \pm 0.0068}$ & $2.0126 \pm 0.0064$ & $3.5170 \pm 0.0136$ \\
 & Độ chệch Ô tô con & $-0.09 \pm 0.02$ & $+0.05 \pm 0.05$ & $+0.11 \pm 0.11$ \\
 & Độ chệch Xe máy & $-0.07 \pm 0.05$ & $+0.04 \pm 0.07$ & $+0.08 \pm 0.14$ \\
\midrule
\multirow{3}{*}{Graph WaveNet \cite{wavenet}} & MAE Tổng thể & $0.9997 \pm 0.0098$ & $2.0263 \pm 0.0100$ & $3.5212 \pm 0.0051$ \\
 & Độ chệch Ô tô con & $-0.08 \pm 0.02$ & $+0.03 \pm 0.03$ & $-0.02 \pm 0.07$ \\
 & Độ chệch Xe máy & $-0.07 \pm 0.03$ & $-0.04 \pm 0.06$ & $-0.01 \pm 0.12$ \\
\midrule
\multirow{3}{*}{ASTGCN \cite{astgcn}} & MAE Tổng thể & $1.0387 \pm 0.0237$ & $2.0758 \pm 0.0086$ & $3.6823 \pm 0.0133$ \\
 & Độ chệch Ô tô con & $-0.09 \pm 0.03$ & $-0.03 \pm 0.04$ & $-0.09 \pm 0.06$ \\
 & Độ chệch Xe máy & $-0.06 \pm 0.04$ & $-0.11 \pm 0.04$ & $-0.30 \pm 0.11$ \\
\midrule
\multirow{3}{*}{STAEformer \cite{staeformer}} & MAE Tổng thể & $0.9749 \pm 0.0113$ & $2.0092 \pm 0.0165$ & $3.5461 \pm 0.0207$ \\
 & Độ chệch Ô tô con & $-0.09 \pm 0.03$ & $+0.02 \pm 0.05$ & $-0.01 \pm 0.07$ \\
 & Độ chệch Xe máy & $-0.05 \pm 0.04$ & $+0.11 \pm 0.10$ & $+0.25 \pm 0.32$ \\
\midrule
\multirow{3}{*}{MegaCRN \cite{megacrn}} & MAE Tổng thể & $1.0219 \pm 0.0095$ & $2.0410 \pm 0.0023$ & $3.5976 \pm 0.0137$ \\
 & Độ chệch Ô tô con & $-0.09 \pm 0.02$ & $-0.05 \pm 0.02$ & $-0.08 \pm 0.04$ \\
 & Độ chệch Xe máy & $-0.10 \pm 0.04$ & $-0.12 \pm 0.11$ & $-0.21 \pm 0.29$ \\
\midrule
\multirow{3}{*}{DSTAGNN \cite{dstagnn}} & MAE Tổng thể & $1.0301 \pm 0.0118$ & $2.0448 \pm 0.0087$ & $3.5981 \pm 0.0186$ \\
 & Độ chệch Ô tô con & $-0.04 \pm 0.03$ & $+0.05 \pm 0.04$ & $+0.09 \pm 0.06$ \\
 & Độ chệch Xe máy & $-0.07 \pm 0.04$ & $+0.01 \pm 0.05$ & $+0.14 \pm 0.23$ \\
\midrule
\multirow{3}{*}{iTransformer \cite{itransformer}} & MAE Tổng thể & $1.0505 \pm 0.0104$ & $2.0448 \pm 0.0090$ & $3.6268 \pm 0.0152$ \\
 & Độ chệch Ô tô con & $-0.07 \pm 0.03$ & $+0.01 \pm 0.04$ & $-0.05 \pm 0.07$ \\
 & Độ chệch Xe máy & $-0.08 \pm 0.01$ & $-0.05 \pm 0.03$ & $-0.15 \pm 0.13$ \\
\bottomrule
\end{tabular}%
}
\end{table}

\begin{figure}[htbp]
\centering
\includegraphics[width=0.90\linewidth]{fig/stage2_density_mae.pdf}
\caption{Sai số dự báo MAE của các mô hình đồ thị Stage 2 theo các phân tầng mật độ phương tiện thực tế (Error Propagation).}
\label{fig:stage2_density_mae_chap5}
\end{figure}

\textbf{Quy luật lan truyền sai số nhận dạng:}
Tại điều kiện mật độ thấp, tất cả các mô hình đồ thị đều hoạt động tương đồng (MAE $\approx 0.97 - 1.05$). Tuy nhiên, tại mức mật độ cao ($>25$ xe/nút), nơi sai số nhận dạng đếm thiếu của Giai đoạn 1 lớn nhất, sự khác biệt trở nên rõ nét. TA-STGCN đạt sai số MAE thấp nhất ($3.4837 \pm 0.0120$), đồng thời kiểm soát độ chệch xe máy về mức gần như triệt tiêu ($-0.03 \pm 0.06$ xe), trong khi STAEformer chệch dương $+0.25$ xe và ASTGCN chệch âm $-0.30$ xe. Điều này chứng minh cơ chế Self-Attention thời gian trong TA-STGCN đã đóng vai trò làm mượt và bù trừ hiệu quả độ chệch âm từ thị giác máy tính truyền lên.

% ==============================================================================
\section{Phân tích phương sai và kiểm định ý nghĩa thống kê}
\label{sec:phan_tich_phuong_sai_thong_ke}

Để kiểm chứng nghiêm cẩn độ tin cậy của các kết quả thực nghiệm, Bảng \ref{tab:seed_variance} báo cáo kết quả kiểm định thống kê chuyên sâu: Khoảng tin cậy Bootstrap 95\% (với 1.000 lần lấy mẫu lại), Độ lệch chuẩn ($\sigma$), Hệ số biến thiên ($\text{CV}$), kiểm định cặp Wilcoxon Signed-Rank kèm hiệu chỉnh Holm-Bonferroni và quy mô ảnh hưởng Cohen's $d_z$ trên toàn bộ hơn $1.8 \times 10^7$ điểm dự báo.

\begin{table}[htbp]
\centering
\caption{Kiểm định thống kê Wilcoxon Signed-Rank, khoảng tin cậy Bootstrap 95\% và quy mô ảnh hưởng Cohen's $d_z$ cho Stage 2 (đối sánh với TA-STGCN).}
\label{tab:seed_variance}
\small
\resizebox{\textwidth}{!}{%
\begin{tabular}{lccccccc}
\toprule
\textbf{Kiến trúc mô hình} & \textbf{Mean MAE} & \textbf{95\% Bootstrap CI} & \textbf{Độ lệch chuẩn} & \textbf{CV (\%)} & \textbf{p-value} & \textbf{Cohen's $d_z$} & \textbf{Cải thiện (\%)} \\
\midrule
Graph WaveNet \cite{wavenet} & $3.2776$ & $3.2733 - 3.2819$ & $0.0049$ & $\mathbf{0.15\%}$ & $< 0.001$ & $-0.024$ & $+1.16\%$ \\
ASTGCN \cite{astgcn} & $3.4218$ & $3.4117 - 3.4319$ & $0.0115$ & $0.33\%$ & $< 0.001$ & $-0.109$ & $+5.61\%$ \\
STGCN Baseline \cite{stgcn} & $3.2703$ & $3.2591 - 3.2815$ & $0.0128$ & $0.39\%$ & $< 0.001$ & $-0.032$ & $+0.93\%$ \\
STAEformer \cite{staeformer} & $3.2953$ & $3.2792 - 3.3114$ & $0.0184$ & $0.56\%$ & $< 0.001$ & $-0.041$ & $+1.70\%$ \\
MegaCRN \cite{megacrn} & $3.3433$ & $3.3339 - 3.3527$ & $0.0107$ & $0.32\%$ & $< 0.001$ & $-0.071$ & $+3.19\%$ \\
DSTAGNN \cite{dstagnn} & $3.3457$ & $3.3309 - 3.3605$ & $0.0169$ & $0.51\%$ & $< 0.001$ & $-0.071$ & $+3.26\%$ \\
iTransformer \cite{itransformer} & $3.3720$ & $3.3595 - 3.3845$ & $0.0143$ & $0.42\%$ & $< 0.001$ & $-0.087$ & $+4.07\%$ \\
\midrule
\textbf{TA-STGCN (Đề xuất)} & $\mathbf{3.2401}$ & $\mathbf{3.2306 - 3.2496}$ & $\mathbf{0.0108}$ & $\mathbf{0.33\%}$ & --- & --- & --- \\
\bottomrule
\end{tabular}%
}
\end{table}

\textbf{Phân tích kiểm định thống kê:}
\begin{itemize}
    \item Kiểm định phi tham số Friedman omnibus trên toàn bộ tập mẫu đánh giá đạt giá trị thống kê $\chi^2 = 9.025,15$ ($p < 0.001$). Kiểm định Wilcoxon Signed-Rank hậu nghiệm đều cho $p < 0.001$. Tuy nhiên, với kích thước mẫu cực lớn ($> 1.8 \times 10^7$ điểm), giá trị $p$ luôn có xu hướng tiến về 0. Do đó, bằng chứng mang tính thực tế phải được xem xét qua quy mô ảnh hưởng Cohen's $d_z$.
    \item Trên dữ liệu sạch, các giá trị Cohen's $d_z$ đều nằm trong khoảng rất hẹp ($|d| < 0.11$), xác nhận rằng trên luồng đầu vào hoàn hảo, sự chênh lệch độ chính xác điểm giữa các mô hình đồ thị là không lớn.
    \item Hệ số biến thiên cực thấp của TA-STGCN ($\text{CV} = 0.33\%$) và khoảng tin cậy Bootstrap 95\% hẹp ($[3.2306 - 3.2496]$) chứng minh rằng hiệu năng của mô hình đề xuất \textbf{hoàn toàn ổn định, không phụ thuộc vào yếu tố ngẫu nhiên của seed khởi tạo}, duy trì tính nhất quán vượt trội hơn so với các mô hình Transformer nặng như STAEformer ($\text{CV} = 0.56\%$).
\end{itemize}

% ==============================================================================
\section{Nghiên cứu độ bền vững trước nhiễu nhận dạng (Robustness Analysis)}
\label{sec:nghien_cuu_do_ben_vung_nhieu}

Đây là nội dung thực nghiệm cốt lõi chứng minh tính ưu việt của kiến trúc đề xuất. Để mô phỏng chân thực các tác động của môi trường giao thông phức tạp (chói sáng, rung lắc camera, che khuất ngẫu nhiên), nghiên cứu tiêm nhiễu phân phối Gamma bất đối xứng vào chuỗi thời gian đầu vào tại 4 cấp độ ứng suất:
\begin{itemize}
    \item \textbf{Mức 1 ($\Delta\text{MAE} = 3.49$):} Mô phỏng điều kiện camera hoạt động thuận lợi ban ngày.
    \item \textbf{Mức 2 ($\Delta\text{MAE} = 4.64$):} Tương đương mức sai số nhận dạng trung bình thực tế của Giai đoạn 1.
    \item \textbf{Mức 3 ($\Delta\text{MAE} = 10.25$):} Tương đương tình trạng tắc nghẽn nặng nề và che khuất cục bộ.
    \item \textbf{Mức 4 ($\Delta\text{MAE} = 13.95$):} Thử nghiệm ứng suất cực đoan ($\text{MAE}_{\max} + 2\sigma$).
\end{itemize}

\begin{figure}[htbp]
\centering
\includegraphics[width=0.88\linewidth]{fig/noise_robustness_study.pdf}
\caption{Đường cong suy thoái sai số MAE của các mô hình đồ thị dưới 4 mức độ nhiễu nhận dạng mô phỏng. TA-STGCN thể hiện năng lực kiểm soát suy thoái tốt nhất.}
\label{fig:noise_robustness_chap5}
\end{figure}

Bảng \ref{tab:bang_suy_thoai_nhieu_chap5} và Bảng \ref{tab:cohens_d_chap5} trình bày chi tiết sai số MAE và quy mô ảnh hưởng Cohen's $d$ hai mẫu độc lập (pooled Cohen's $d$) giữa các kiến trúc so với TA-STGCN.

\begin{table}[htbp]
\centering
\caption{Độ suy thoái sai số MAE của các mô hình đồ thị dưới 4 mức độ nhiễu nhận dạng ($\text{Mean} \pm \text{Std}$).}
\label{tab:bang_suy_thoai_nhieu_chap5}
\small
\resizebox{\textwidth}{!}{%
\begin{tabular}{lcccc}
\toprule
\textbf{Kiến trúc mô hình} & \textbf{Mức 1 (3.49)} & \textbf{Mức 2 (4.64)} & \textbf{Mức 3 (10.25)} & \textbf{Mức 4 (13.95)} \\
\midrule
STGCN Baseline \cite{stgcn} & $5.3093 \pm 0.1161$ & $6.6614 \pm 0.1532$ & $12.0733 \pm 0.3396$ & $13.1724 \pm 0.3515$ \\
Graph WaveNet \cite{wavenet} & $5.7470 \pm 0.1056$ & $7.2378 \pm 0.1324$ & $13.2723 \pm 0.3303$ & $15.1046 \pm 0.4398$ \\
ASTGCN \cite{astgcn} & $5.7709 \pm 0.1521$ & $6.9956 \pm 0.2346$ & $11.8790 \pm 0.5428$ & $13.4685 \pm 0.7556$ \\
STAEformer \cite{staeformer} & $\mathbf{5.2238 \pm 0.2964}$ & $6.4260 \pm 0.5436$ & $10.5916 \pm 1.3943$ & $10.8454 \pm 1.2209$ \\
MegaCRN \cite{megacrn} & $5.9061 \pm 0.1227$ & $7.3144 \pm 0.1412$ & $13.1555 \pm 0.5314$ & $15.1106 \pm 0.8717$ \\
DSTAGNN \cite{dstagnn} & $6.2166 \pm 0.0514$ & $7.8276 \pm 0.2608$ & $12.5753 \pm 0.6823$ & $13.3705 \pm 0.8442$ \\
iTransformer \cite{itransformer} & $5.7070 \pm 0.0999$ & $6.9084 \pm 0.1786$ & $10.3254 \pm 0.7281$ & $10.8422 \pm 0.9367$ \\
\midrule
\textbf{TA-STGCN (Đề xuất)} & $5.2401 \pm 0.1000$ & $\mathbf{6.3120 \pm 0.1932}$ & $\mathbf{9.6381 \pm 0.8461}$ & $\mathbf{10.1532 \pm 1.1153}$ \\
\bottomrule
\end{tabular}%
}
\end{table}

\begin{table}[htbp]
\centering
\caption{Quy mô ảnh hưởng Cohen's $d$ so với TA-STGCN dưới nhiễu nhận dạng (dương thể hiện TA-STGCN vượt trội).}
\label{tab:cohens_d_chap5}
\small
\begin{tabular}{lcc}
\toprule
\textbf{Kiến trúc đối chứng} & \textbf{Mức 2 (Nhiễu thực tế: 4.64)} & \textbf{Mức 4 (Nhiễu cực đoan: 13.95)} \\
\midrule
STGCN Baseline \cite{stgcn} & $+2.004$ & $+3.651$ \\
Graph WaveNet \cite{wavenet} & $+5.590$ & $+5.841$ \\
ASTGCN \cite{astgcn} & $+3.181$ & $+3.480$ \\
STAEformer \cite{staeformer} & $+0.279$ & $+0.592$ \\
MegaCRN \cite{megacrn} & $+5.924$ & $+4.953$ \\
DSTAGNN \cite{dstagnn} & $+6.604$ & $+3.253$ \\
iTransformer \cite{itransformer} & $+3.206$ & $+0.669$ \\
\bottomrule
\end{tabular}
\end{table}

\textbf{Luận điểm khoa học cốt lõi từ thực nghiệm nhiễu:}
\begin{enumerate}
    \item \textbf{Tính bất đối xứng về chi phí và độ bền vững (Robustness-vs-Cost Asymmetry):} Việc tích hợp cơ chế Self-Attention theo thời gian vào STGCN chỉ làm tăng nhẹ mức sai số không đáng kể trên dữ liệu sạch, nhưng mang lại một \textbf{bước nhảy vọt về độ bền vững dưới nhiễu}. Tại Mức 2 (mức sai số thực tế của Stage 1), TA-STGCN vượt trội STGCN với Cohen's $d = +2.004$; tại Mức 4, sai số của TA-STGCN chỉ là $10.1532 \pm 1.1153$ so với $13.1724 \pm 0.3515$ của STGCN Baseline ($d = +3.651$).
    \item \textbf{Sự phân hóa rõ nét giữa các dòng kiến trúc:} Các mô hình thuần tích chập hoặc hồi quy không có cơ chế chú ý (Graph WaveNet, MegaCRN, STGCN) nhanh chóng bị suy thoái nặng nề khi gặp nhiễu (MAE vọt lên $13.17 - 15.11$, Cohen's $d$ từ $+3.65$ đến $+5.84$). Trong khi đó, các mô hình tích hợp cơ chế chú ý thời gian (TA-STGCN, STAEformer, iTransformer) giữ vững sai số ở mức $10.15 - 10.84$.
    \item \textbf{Tầm quan trọng của trật tự xử lý (Temporal-First so với Spatial-First):} ASTGCN mặc dù có cơ chế Attention nhưng lại tính toán chú ý không gian trước, chú ý thời gian sau. Khi một nút camera gặp nhiễu đột biến, tầng Spatial Attention bị đánh lừa và lập tức khuếch tán giá trị sai lệch sang toàn bộ đồ thị mạng lưới. Ngược lại, TA-STGCN tính toán Temporal Attention trước trên từng nút độc lập, hoạt động như một \textbf{bộ lọc thông thấp thích ứng} triệt tiêu xung nhiễu trước khi tín hiệu được truyền qua toán tử tích chập không gian ChebNet.
    \item \textbf{Hiệu quả vượt trội về tài nguyên phần cứng:} TA-STGCN đạt độ bền vững ngang ngửa hoặc nhỉnh hơn STAEformer nhưng tiết kiệm tới $25.2\%$ bộ nhớ kích hoạt GPU đỉnh ($8.606,9$ MB so với $11.509,1$ MB) và cắt giảm $70.4\%$ số tham số so với ASTGCN ($453\text{K}$ so với $1.533\text{K}$).
\end{enumerate}

% ==============================================================================
\section{Phân tích tính giải thích của cơ chế chú ý thời gian}
\label{sec:kha_nang_giai_thich_attention}

Hình \ref{fig:attention_heatmap_chap5} và Bảng \ref{tab:attention_ty_le_chap5} trình bày phân tích định lượng ma trận trọng số chú ý thời gian $\mathbf{A} \in \mathbb{R}^{24 \times 24}$ trích xuất từ tầng Temporal Attention của TA-STGCN qua 5 khung giờ đặc trưng trong ngày.

\begin{figure}[htbp]
\centering
\includegraphics[width=0.92\linewidth]{fig/attention_comparison_evening_vs_noon.pdf}
\caption{So sánh ma trận trọng số chú ý thời gian $\mathbf{A} \in \mathbb{R}^{24 \times 24}$ giữa (a) Giờ cao điểm tối (16:30--18:30) và (b) Giờ bình thường buổi trưa (11:30--13:30).}
\label{fig:attention_heatmap_chap5}
\end{figure}

\begin{table}[htbp]
\centering
\caption{Định lượng phân bổ trọng số chú ý thời gian qua các khung giờ giao thông trong ngày ($\text{Mean} \pm \text{Std}$).}
\label{tab:attention_ty_le_chap5}
\small
\resizebox{\textwidth}{!}{%
\begin{tabular}{lcccc}
\toprule
\textbf{Khung thời gian} & \textbf{Lưu lượng TB} & \textbf{Trọng số gần (20m)} & \textbf{Trọng số xa (60--120m)} & \textbf{Tỷ lệ (Gần / Xa)} \\
\midrule
Đêm khuya (00:00--05:00) & 5.0 & $0.0427 \pm 0.0038$ & $0.0417 \pm 0.0016$ & $1.03\times \pm 0.13\times$ \\
Cao điểm sáng (07:30--09:30) & 35.5 & $0.0461 \pm 0.0021$ & $0.0389 \pm 0.0013$ & $1.19\times \pm 0.09\times$ \\
Buổi trưa (11:30--13:30) & 36.0 & $0.0416 \pm 0.0013$ & $0.0419 \pm 0.0011$ & $0.99\times \pm 0.06\times$ \\
\textbf{Cao điểm tối (16:30--18:30)} & \textbf{39.5} & $\mathbf{0.0662 \pm 0.0067}$ & $\mathbf{0.0304 \pm 0.0032}$ & $\mathbf{2.22\times \pm 0.41\times}$ \\
Đêm muộn (21:00--23:00) & 17.5 & $0.0644 \pm 0.0062$ & $0.0314 \pm 0.0023$ & $2.08\times \pm 0.35\times$ \\
\bottomrule
\end{tabular}%
}
\end{table}

\textbf{Cơ chế tái phân bổ trọng số thích ứng theo trạng thái giao thông:}
\begin{itemize}
    \item Trong các giai đoạn giao thông ổn định (như Buổi trưa 11:30--13:30), mô hình phân bổ trọng số gần như đồng đều tuyệt đối giữa các bước thời gian gần và xa (tỷ lệ $0.99\times \pm 0.06\times$, với giá trị trọng số xấp xỉ mức đều $1/24 \approx 0.0417$).
    \item Khi bước vào \textbf{Giờ cao điểm tối (16:30--18:30)} --- thời điểm lưu lượng xe máy tăng vọt và có biến động động học cực lớn --- mô hình tự động co hẹp tiêu điểm chú ý: tăng mạnh trọng số cho khoảng thời gian 20 phút gần nhất lên $0.0662 \pm 0.0067$ (tăng $58.8\%$ so với mức phân bổ đều) và nén trọng số lịch sử xa (60--120 phút trước) xuống $0.0304 \pm 0.0032$. Tỷ lệ Gần/Xa đạt đỉnh \textbf{$2.22\times \pm 0.41\times$}. Hiện tượng tương tự diễn ra vào giờ Đêm muộn ($2.08\times \pm 0.35\times$).
    \item Hành vi này chứng minh rằng cơ chế Self-Attention đã tự động học được cách chuyển dịch trạng thái: ưu tiên nắm bắt các xu hướng động học tức thời trong giờ cao điểm biến động và cân bằng toàn dải lịch sử trong giờ bình ổn.
\end{itemize}

% ==============================================================================
\section{Nghiên cứu loại trừ thành phần (Ablation Study)}
\label{sec:ablation_study}

Để cô lập và đánh giá chính xác vai trò đóng góp của từng thành phần cấu trúc trong TA-STGCN, Bảng \ref{tab:ablation_full_chap5} trình bày kết quả thử nghiệm triệt tiêu trên 10 biến thể kiến trúc độc lập qua 5 seed ngẫu nhiên.

\begin{table}[htbp]
\centering
\caption{Kết quả Ablation Study đánh giá các biến thể kiến trúc của TA-STGCN qua 5 seed ngẫu nhiên ($\text{Mean} \pm \text{Std}$).}
\label{tab:ablation_full_chap5}
\small
\resizebox{\textwidth}{!}{%
\begin{tabular}{lcccccc}
\toprule
\textbf{Biến thể mô hình} & \textbf{Số tham số} & \textbf{RAM đỉnh (MB)} & \textbf{MAE Tổng} & \textbf{MAE Ô tô} & \textbf{MAE Xe máy} & \textbf{MAPE (\%)} \\
\midrule
\textbf{TA-STGCN (Mô hình chính, $L=2$)} & \textbf{453K} & \textbf{8.606,9} & $\mathbf{3.2401 \pm 0.0108}$ & $\mathbf{1.2087 \pm 0.0019}$ & $\mathbf{2.7286 \pm 0.0096}$ & $\mathbf{11.32\% \pm 0.04\%}$ \\
STGCN Baseline (Không Attention) \cite{stgcn} & 375K & 4.837,6 & $3.2703 \pm 0.0128$ & $1.2188 \pm 0.0063$ & $2.7467 \pm 0.0191$ & $11.44\% \pm 0.08\%$ \\
STGCN (Large, 487K params) & 487K & 5.793,9 & $3.2805 \pm 0.0134$ & $1.2157 \pm 0.0032$ & $2.7620 \pm 0.0194$ & $11.53\% \pm 0.11\%$ \\
TA-STGCN ($h=1$, Đơn đầu) & 453K & 8.158,1 & $3.2501 \pm 0.0075$ & $1.2088 \pm 0.0025$ & $2.7357 \pm 0.0075$ & $11.32\% \pm 0.03\%$ \\
TA-STGCN ($C=32$, Kênh hẹp) & 126K & 10.591,5 & $3.2717 \pm 0.0057$ & $1.2177 \pm 0.0023$ & $2.7542 \pm 0.0047$ & $11.46\% \pm 0.05\%$ \\
TA-STGCN ($K=1$, Không đồ thị không gian) & 427K & 8.033,6 & $3.2488 \pm 0.0105$ & $1.2104 \pm 0.0017$ & $2.7366 \pm 0.0119$ & $11.32\% \pm 0.04\%$ \\
TA-STGCN ($h=8$, 8 đầu chú ý) & 453K & 9.205,4 & $3.2494 \pm 0.0118$ & $1.2103 \pm 0.0012$ & $2.7367 \pm 0.0118$ & $11.36\% \pm 0.04\%$ \\
TA-STGCN (Độ sâu $L=3$ khối) & 647K & 10.536,0 & $3.2539 \pm 0.0098$ & $1.2113 \pm 0.0019$ & $2.7404 \pm 0.0096$ & $11.38\% \pm 0.05\%$ \\
TA-STGCN (Attention đặt ở đầu khối) & 489K & 8.741,3 & $3.2485 \pm 0.0087$ & $1.2060 \pm 0.0018$ & $2.7301 \pm 0.0077$ & $11.33\% \pm 0.04\%$ \\
TA-STGCN (Attention đặt ở cuối khối) & 453K & 8.961,9 & $3.2680 \pm 0.0099$ & $1.2097 \pm 0.0020$ & $2.7468 \pm 0.0079$ & $11.44\% \pm 0.09\%$ \\
\bottomrule
\end{tabular}%
}
\end{table}

\textbf{Bốn phát hiện kiến trúc quan trọng từ Ablation Study:}
\begin{enumerate}
    \item \textbf{Vị trí đặt tầng Temporal Attention:} Cấu hình đặt Attention ở giữa hai tầng tích chập mang lại hiệu năng cao nhất ($\text{MAE} = 3.2401$), vượt trội so với đặt ở cuối khối ($3.2680$) hoặc ở đầu khối ($3.2485$). Việc đặt Attention ngay trước toán tử ChebNet bảo đảm các xung nhiễu thời gian tại từng nút được lọc sạch trước khi diễn ra sự lan truyền thông tin sang các nút lân cận trên đồ thị.
    \item \textbf{Cấu hình số đầu chú ý ($h=4$):} Số đầu chú ý $h=4$ đem lại kết quả tối ưu nhất ($3.2401$). Giảm xuống đơn đầu ($h=1$) làm sai số tăng lên $3.2501$ do mất khả năng nắm bắt đa dạng các tần số dao động; trong khi tăng lên $h=8$ khiến sai số tăng nhẹ ($3.2494$) và dung lượng RAM đỉnh tăng thêm $598.5$ MB do chi phí ma trận kích hoạt.
    \item \textbf{Số lượng khối không -- thời gian ($L=2$):} Tăng độ sâu mạng lên $L=3$ khối làm số lượng tham số tăng lên 647K và bộ nhớ đỉnh tăng lên 10.536 MB, nhưng MAE lại tăng lên $3.2539$. Điều này phản ánh hiện tượng làm mịn quá mức (over-smoothing) phổ biến trong các mạng nơ-ron đồ thị sâu.
    \item \textbf{Đóng góp của cấu trúc topo đồ thị không gian ($K=3$):} Khi triệt tiêu toán tử ChebNet (đặt $K=1$, bỏ qua hoàn toàn liên kết cạnh giữa các nút giao), sai số tăng lên $3.2488$. Kết quả này khẳng định cấu trúc topo kết nối đường bộ thực tế từ OpenStreetMap đóng góp thực chất vào độ chính xác dự báo lưu lượng toàn mạng lưới.
\end{enumerate}


% Chương 6: Thảo luận, mô hình triển khai và hướng phát triển
\chapter{Thảo luận, mô hình triển khai và hướng phát triển}
\label{chap:thao_luan}

Chương này thảo luận sâu về ý nghĩa học thuật, tính kinh tế phần cứng trong triển khai thực tế, tính công bằng của các thực nghiệm đối sánh, đồng thời thẳng thắn chỉ ra các giới hạn hiện tại của công trình nghiên cứu và xác định các hướng mở rộng tiếp nối trong tương lai.

% ==============================================================================
\section{Phác thảo thiết kế kiến trúc phân tán Edge--Cloud (Định hướng tương lai)}
\label{sec:kien_truc_trien_khai_edge_cloud}

Để làm rõ động lực của các lựa chọn thiết kế được đánh giá trong thực nghiệm (tính kinh tế về tham số, bộ nhớ kích hoạt GPU và cơ chế truyền tải chỉ chứa vector số đếm), nghiên cứu phác thảo ý niệm về một kiến trúc phân tán giữa thiết bị biên và máy chủ trung tâm (Edge--Cloud Split). \textbf{Cần nhấn mạnh rằng đây là bản phác thảo thiết kế định hướng ứng dụng trong tương lai (Future Work), không phải kết quả đo đạc thực nghiệm trên phần cứng nhúng}. Toàn bộ các kết quả thực nghiệm được báo cáo trong đề tài đều được thực hiện trên hệ thống xử lý tập trung chuẩn hóa đã mô tả tại Chương \ref{chap:phuong_phap} và Chương \ref{chap:thuc_nghiem}.

\begin{figure}[htbp]
\centering
\includegraphics[width=0.98\linewidth]{fig/edge_cloud_architecture.pdf}
\caption{Bản phác thảo thiết kế kiến trúc phân tán Edge--Cloud định hướng ứng dụng: xử lý nhận dạng tại thiết bị biên và dự báo dòng giao thông trên đồ thị tại máy chủ trung tâm.}
\label{fig:edge_cloud_arch_chap6}
\end{figure}

Trong kiến trúc định hướng triển khai (Hình \ref{fig:edge_cloud_arch_chap6}):
\begin{itemize}
    \item \textbf{Tại thiết bị biên (Edge Devices -- các trạm camera quan sát):} Mô hình Giai đoạn 1 (mạng xương sống EfficientNet-B5 hoặc phương án siêu nhẹ MobileNetV3-Large sau khi lượng tử hóa số nguyên INT8 hoặc FP16 thông qua các runtime nhúng chuyên dụng như OpenVINO, TensorRT) được nạp vào bộ nhớ thiết bị biên. Thiết bị thực hiện suy luận ảnh chụp tĩnh trực tiếp trong bộ nhớ đệm và chỉ phát đi gói tin nhị phân chứa vector số lượng phương tiện theo phân lớp: $[\hat{c}_{\text{car}}, \hat{c}_{\text{bike}}]^T$.
    \item \textbf{Tại máy chủ trung tâm (Central Cloud Server):} Tiếp nhận các vector số đếm từ toàn bộ $N=608$ nút camera trên mạng lưới đường bộ đô thị, cập nhật vào ma trận trạng thái chuỗi thời gian và thực thi mô hình TA-STGCN để dự báo lưu lượng toàn mạng lưới trong 30 phút tiếp theo.
\end{itemize}

\subsection{Ý nghĩa đối với băng thông truyền thông hạ tầng}
Trong các giải pháp giám sát giao thông truyền thống, việc truyền tải liên tục luồng video độ nét cao hoặc ảnh tĩnh độ phân giải lớn từ hàng trăm điểm quan sát về trung tâm chỉ huy gây áp lực rất lớn lên hạ tầng mạng viễn thông. 

Bằng cách chuyển dịch khâu nhận dạng xuống thiết bị biên, mỗi nút quan sát chỉ cần truyền tải vector số đếm hai chiều thay vì truyền hình ảnh thô. Về mặt ý niệm lý thuyết, việc truyền một vector hai số thực (vài bytes) so với một ảnh nén JPEG kích thước hàng trăm kilobytes đem lại tiềm năng cắt giảm lưu lượng truyền dẫn rất lớn. Tuy nhiên, tỷ lệ cắt giảm lý thuyết này chưa tính đến các chi phí giao thức truyền thông (protocol overhead), cơ chế xác thực, gói tin điều khiển mạng và việc truyền lại khi rớt gói. Do đó, việc lượng hóa chính xác băng thông thực tế đòi hỏi các đo đạc thực nghiệm toàn diện trên môi trường mạng thực địa trong các nghiên cứu tương lai.

\subsection{Bảo vệ quyền riêng tư hình ảnh theo thiết kế (Privacy-by-Design)}
Việc truyền tải và lưu trữ tập trung hình ảnh camera nơi công cộng tiềm ẩn nguy cơ làm lộ lọt thông tin cá nhân của người tham gia giao thông (nhận diện khuôn mặt hoặc biển số phương tiện).

Trong kiến trúc phân tán được phác thảo, nguyên lý bảo vệ quyền riêng tư được bảo đảm thông qua cơ chế xử lý nội bộ tại chỗ:
\begin{itemize}
    \item Khung hình camera được giải mã trực tiếp trong bộ nhớ RAM tạm thời của vi xử lý biên, đưa qua mạng nơ-ron hồi quy đếm xe và được giải phóng ngay lập tức sau khi có kết quả.
    \item Không có bất kỳ hình ảnh thô hay đặc trưng nhận dạng cá nhân nào được lưu trữ cố định trên bộ nhớ biên hoặc truyền dẫn qua môi trường mạng Internet.
    \item Dữ liệu rời khỏi thiết bị biên chỉ là các con số thống kê tổng hợp về lưu lượng phương tiện của tuyến đường, hoàn toàn loại bỏ nguy cơ rò rỉ dữ liệu hình ảnh nhạy cảm.
\end{itemize}

% ==============================================================================
\section{Tính kinh tế phần cứng: Bộ nhớ kích hoạt GPU đỉnh so với số lượng tham số tĩnh}
\label{sec:tinh_kinh_te_phan_cung}

Một phát hiện thực nghiệm quan trọng từ các kết quả tại Chương \ref{chap:thuc_nghiem} (Bảng \ref{tab:sieu_tham_so_chap5_models}) là: \textbf{Số lượng tham số tĩnh ($N_{\text{params}}$) không phản ánh đầy đủ chi phí phần cứng thực tế của các mô hình đồ thị không -- thời gian trong môi trường đô thị quy mô lớn}.

\begin{enumerate}
    \item \textbf{Hiện tượng bùng nổ bộ nhớ kích hoạt (Activation Memory Explosion):} Các mô hình Transformer đồ thị tiên tiến như STAEformer duy trì lượng tham số tĩnh ở mức vừa phải ($568\text{K}$ tham số), song việc tính toán các ma trận chú ý mềm liên hợp (Joint Spatio-Temporal Attention) trên toàn bộ $N=608$ nút khiến lượng bộ nhớ kích hoạt GPU đỉnh vọt lên tới \textbf{11.509,1 MB}. Mức tiêu thụ này chạm sát ngưỡng dung lượng 12 GB của các dòng card đồ họa thương mại phổ biến, gây khó khăn lớn cho việc triển khai trên các cụm máy chủ kinh phí hạn chế.
    \item \textbf{Bản chất thiết kế của mô hình iTransformer:} Mặc dù iTransformer đạt mức tiêu thụ bộ nhớ GPU cực kỳ thấp (\textbf{265,2 MB}), điều này hoàn toàn xuất phát từ bản chất thiết kế Inverted Transformer của mô hình: coi mỗi nút giao thông như một biến số độc lập và áp dụng cơ chế chú ý trên chiều biến số, hoàn toàn không xây dựng hay duy trì ma trận không gian $N \times N$. Do không mô hình hóa tương quan không gian qua đồ thị, iTransformer là một mô hình chuỗi thời gian đa biến thuần túy (non-graph) và được xử lý riêng biệt trong các đánh giá.
    \item \textbf{Cân bằng tối ưu của TA-STGCN:} Bằng cách tách rời phép tích chập không gian Chebyshev ($K=3$) trên đồ thị thưa khỏi cơ chế Self-Attention trên chiều thời gian độc lập cho từng nút, TA-STGCN chỉ tiêu tốn \textbf{8.606,9 MB} bộ nhớ GPU đỉnh (tiết kiệm $25.2\%$ so với STAEformer) và giảm $70.4\%$ số lượng tham số so với ASTGCN ($453\text{K}$ so với $1.533\text{K}$). Sự cân đối đồng thời trên cả hai trục tham số và bộ nhớ kích hoạt giúp mô hình phù hợp với các hệ thống máy chủ vận hành thực tế.
\end{enumerate}

% ==============================================================================
\section{Tính khách quan và công bằng trong so sánh thực nghiệm}
\label{sec:tinh_cong_bang_thuc_nghiem}

Để đảm bảo tính nghiêm cẩn học thuật, các thực nghiệm đối sánh ở Giai đoạn 2 được thiết kế như một \textbf{thực nghiệm đối chứng có kiểm soát chặt chẽ (controlled implementation comparison)} thay vì một cuộc chạy đua tối ưu hóa riêng lẻ từng mô hình:

\begin{enumerate}
    \item \textbf{Sử dụng chung hàm mục tiêu (Unified Huber Objective):} Tất cả các mô hình đối chứng đều được huấn luyện dưới cùng một hàm mất mát Huber Loss ($\delta=1.0$). Một số mô hình gốc trong y văn ban đầu được đề xuất với hàm mục tiêu khác (chẳng hạn Graph WaveNet ban đầu sử dụng hàm mất mát $L_1$). Việc huấn luyện lại toàn bộ mô hình dưới hàm Huber Loss là một sự đánh đổi có chủ đích: kiểm soát đồng nhất cách mô hình phản ứng với các điểm ngoại lai để cô lập yếu tố sai khác về mặt kiến trúc. Do đó, kết quả của các mô hình đối chứng được hiểu là hiệu năng đạt được dưới hàm mục tiêu dùng chung, không phải kết quả công bố riêng lẻ của từng bài báo gốc.
    \item \textbf{Không tinh chỉnh siêu tham số riêng biệt cho từng baseline:} Nghiên cứu không thực hiện tìm kiếm siêu tham số vét cạn riêng lẻ cho từng mô hình đối chứng trên đồ thị 608 nút; các mô hình sử dụng cấu hình chuẩn do các tác giả nguyên bản đề xuất, được điều chỉnh kích thước đầu vào và đầu ra cho tương thích. Do đó, một số mô hình đối chứng có thể đạt độ chính xác cao hơn nếu được tinh chỉnh sâu. Đây là lý do nghiên cứu không nhấn mạnh vào các chênh lệch sai số nhỏ trên dữ liệu sạch, phù hợp với kết quả kiểm định kích thước ảnh hưởng ($|d| < 0.1$) cho thấy các mô hình đồ thị đạt độ chính xác điểm gần như tương đương nhau.
    \item \textbf{Phân định rõ ràng bản chất non-graph của iTransformer:} iTransformer được đưa vào thực nghiệm như một mô hình tham chiếu mạnh về chuỗi thời gian đa biến, không phải là mô hình dự báo cấu trúc đồ thị. Mô hình này được loại trừ khỏi các kết luận so sánh nội bộ giữa các kiến trúc đồ thị. Việc iTransformer đạt độ bền vững tốt trước nhiễu (ở Mức 4) cung cấp bằng chứng cho thấy một phần năng lực kháng nhiễu bắt nguồn từ cơ chế chú ý theo trục thời gian nói chung chứ không riêng gì cấu trúc đồ thị.
\end{enumerate}

% ==============================================================================
\section{Các giới hạn học thuật của nghiên cứu và Hướng phát triển tiếp nối}
\label{sec:han_che_va_huong_phat_trien}

Nghiên cứu thẳng thắn nhìn nhận các giới hạn kỹ thuật hiện tại và xác định các hướng nghiên cứu mở rộng trong tương lai:

\begin{enumerate}
    \item \textbf{Giới hạn trần sai số phụ thuộc giữa hai giai đoạn:} Mô hình dự báo Giai đoạn 2 được huấn luyện trực tiếp trên chuỗi dữ liệu số đếm sinh ra từ Giai đoạn 1. Do đó, độ chính xác và độ tin cậy của toàn bộ hệ thống dự báo bị giới hạn trên bởi biên độ sai số của khâu nhận dạng thị giác ban đầu.
    \item \textbf{Thách thức quy mô tính toán khi triển khai thực tế trên diện rộng:} Mặc dù mô hình đồ thị Giai đoạn 2 tương đối gọn nhẹ ($453\text{K}$ tham số), việc triển khai đồng loạt mô hình Giai đoạn 1 (EfficientNet-B5 với 28.3M tham số) trên 608 camera biên tương đương với việc vận hành tổng cộng hơn $17$ tỷ tham số trên toàn đô thị. Đây là một thách thức lớn về chi phí phần cứng biên, độ trễ và năng lượng tiêu thụ. Mặc dù MobileNetV3-Large (4.2M tham số) là một lựa chọn nhỏ gọn hơn, mô hình này phải đánh đổi bằng mức sai số Stage 1 cao hơn ($4.0126$ so với $3.6363$). Việc lượng tử hóa (INT8/FP16), chưng cất tri thức (Knowledge Distillation) và đo đạc thông lượng thực tế trên các bo mạch nhúng (như NVIDIA Jetson, Coral TPU) là hướng đi bắt buộc trước khi đưa vào ứng dụng thực tế.
    \item \textbf{Đánh đổi của phương pháp hồi quy đếm xe trực tiếp:} Nghiên cứu định vị cơ chế hồi quy đếm xe trực tiếp từ ảnh tĩnh như một giải pháp thực tế giúp loại bỏ chi phí gán nhãn điểm/bounding box phức tạp và không đòi hỏi căn chỉnh thông số hình học từng camera. Tuy nhiên, nghiên cứu chưa thực hiện đối sánh đối đầu trực tiếp với các phương pháp ước lượng bản đồ mật độ (Density-map estimators như MCNN \cite{mcnn}, CSRNet \cite{csrnet}, CAN \cite{can}) hoặc các mô hình phát hiện kết hợp bám vết (Tracking-by-Detection như YOLO \cite{yolo} + ByteTrack) trên cùng một tập dữ liệu gán nhãn vi mô.
    \item \textbf{Phạm vi không gian -- thời gian của bộ dữ liệu:} Bộ dữ liệu thực nghiệm được thu thập tại một đô thị đơn lẻ (TP. Hồ Chí Minh) trong khoảng thời gian 3 tháng, chưa bao quát đầy đủ các biến thiên thời tiết cực đoan theo mùa (mùa mưa ngập kéo dài hoặc sương mù dày đặc).
    \item \textbf{Các dòng kiến trúc chuỗi thời gian mới (2024--2025):} Thực nghiệm của đề tài bao quát các kiến trúc đồ thị và Transformer tiêu biểu đến năm 2024, song chưa tích hợp các dòng mô hình mới nhất xuất hiện trong giai đoạn 2024--2025 như Spatio-Temporal State-Space Models (Mamba \cite{mamba}), Diffusion STG, hoặc các mô hình chuỗi thời gian dựa trên Patch / LLM (như PatchTST \cite{patchtst}, TimeXer). Các họ kiến trúc này chưa được chuẩn hóa cho dữ liệu đồ thị có cấu trúc phân tách chủng loại xe thu thập từ camera, đòi hỏi những nghiên cứu kỹ thuật sâu hơn trong tương lai.
    \item \textbf{Tính thích ứng của ma trận kề không gian:} Đồ thị RBF xây dựng từ OpenStreetMap thể hiện cấu trúc kết nối hình học tĩnh, chưa thể tự thích ứng với các biến cố phi vật lý hoặc sự cố đột xuất (tai nạn giao thông làm chặn đường, ngập nước cục bộ). Hướng phát triển tiếp theo là nghiên cứu kết hợp đồ thị tĩnh với cơ chế học đồ thị thích ứng động (Dynamic Adaptive Graph Learning), đồng thời mở rộng phân tích độ bền vững dưới các phân phối nhiễu xác suất khác (như Weibull hoặc Log-normal).
\end{enumerate}


% Kết luận và kiến nghị
\chapter*{Kết luận và kiến nghị}
\addcontentsline{toc}{chapter}{Kết luận và kiến nghị}

\section*{1. Kết luận khoa học}
Đề tài nghiên cứu khoa học \textbf{"Nghiên cứu và phát triển khung mô hình thu nhận hình ảnh và dự báo dòng giao thông đô thị đa bước hai giai đoạn dựa trên mạng đồ thị không -- thời gian từ camera giám sát"} được thực hiện nhằm giải quyết một bài toán quan trọng trong quản lý đô thị thông minh: khai thác hiệu quả hạ tầng camera quan sát công cộng để thiết lập mạng lưới giám sát và dự báo lưu lượng giao thông tin cậy trong điều kiện dòng xe máy hỗn hợp tại các đô thị lớn của Việt Nam.

Qua quá trình nghiên cứu lý thuyết, mô hình hóa toán học và thực nghiệm trên quy mô toàn thành phố, đề tài đã hoàn thành các mục tiêu đặt ra với những kết quả cụ thể:

\begin{enumerate}
    \item \textbf{Xây dựng và công bố bộ dữ liệu chuẩn mực IC4SD-Traffic-HCM:} 
    Đề tài đã hoàn thành việc xây dựng bộ dữ liệu giao thông quy mô đô thị mở đầu tiên tại Việt Nam bao gồm \textbf{608 nút camera giám sát thực tế} tại TP. Hồ Chí Minh, bao gồm:
    \begin{itemize}
        \item 5.012 ảnh tĩnh độ phân giải cao được gán nhãn thủ công theo hai lớp ô tô con và xe máy, đạt hệ số tương quan nội lớp liên thẩm định viên cao ($\text{ICC} = 0.973, p < 0.001$).
        \item Chuỗi thời gian liên tục trong 86 ngày với 24.768 bước quan sát (chu kỳ 5 phút) bảo toàn phân loại phương tiện, được làm mượt bằng bộ lọc trượt kép nhân quả không gây rò rỉ dữ liệu tương lai.
        \item Đồ thị mạng lưới đường bộ OpenStreetMap với 2.452 cạnh định hướng liên thông thực tế, phản ánh chính xác cấu trúc hình học và quy định lưu thông của các tuyến đường.
    \end{itemize}
    Bộ dữ liệu và mã nguồn cài đặt được công bố công khai, tạo môi trường kiểm chuẩn chuẩn mực cho cộng đồng học thuật.

    \item \textbf{Hiệu quả của giải pháp hồi quy đếm xe trực tiếp (Giai đoạn 1):} 
    Nghiên cứu đã khẳng định tính hiệu quả của phương pháp hồi quy đếm xe trực tiếp dựa trên mạng xương sống EfficientNet-B5 \cite{efficientnet} kết hợp đầu hồi quy MLP thống nhất và hàm mất mát Huber Loss \cite{huber}. Cách tiếp cận này hạn chế ảnh hưởng của việc che khuất tương hỗ trong dòng xe máy và loại bỏ hoàn toàn chi phí tính toán của thuật toán NMS truyền thống, đạt sai số $\text{MAE} = 3.6363$ và tốc độ suy luận $35.2$ ms/batch, tiết kiệm $43.0\%$ lượng phép tính FLOPs so với ResNet-50 \cite{resnet}. Bản đồ nhiệt Grad-CAM++ xác nhận mô hình tập trung chính xác vào thân xe thay vì cảnh nền.

    \item \textbf{Đề xuất kiến trúc đồ thị cải tiến TA-STGCN (Giai đoạn 2):} 
    Đề tài đã phát triển mô hình tích hợp cơ chế chú ý tự thân theo thời gian đa đầu (Multi-Head Temporal Self-Attention) \cite{attention} xen kẽ giữa hai khối ST-Conv \cite{stgcn} trên đồ thị OpenStreetMap. Mô hình đạt sai số dự báo $\text{MAE} = 3.2401 \pm 0.0108$, duy trì tính ổn định qua 5 seed ngẫu nhiên ($\text{CV} = 0.33\%$) và kiểm soát tốt mức tăng sai số ở tầm nhìn dự báo xa 30 phút ($t+6$).

    \item \textbf{Đóng góp về lý thuyết lan truyền sai số và độ bền vững trước nhiễu:} 
    Đây là đóng góp khoa học trọng tâm của công trình. Đề tài đã chứng minh bằng giải tích toán học và kiểm chứng thực nghiệm rằng cơ chế chú ý theo trục thời gian hoạt động như một bộ lọc thông thấp thích ứng dữ liệu, làm giảm phương sai nhiễu tỷ lệ nghịch với độ rộng cửa sổ chú ý hiệu dụng ($\text{Var} = \sigma_\eta^2 / T_{\text{eff}}$). Dưới tác động của sai số nhận dạng mô phỏng thực tế (phân phối Gamma dị biệt có độ lệch âm do che khuất xe máy), TA-STGCN duy trì tốc độ suy thoái sai số thấp nhất, vượt trội rõ rệt các mô hình không có cơ chế chú ý thời gian (chỉ số Cohen's $d$ lên tới $+5.84$ so với Graph WaveNet \cite{wavenet} ở mức nhiễu cực hạn) và đạt độ bền bỉ tương đương các mô hình Transformer lớn (như STAEformer \cite{staeformer}) trong khi tiết kiệm $25.2\%$ bộ nhớ GPU VRAM và $70.4\%$ số lượng tham số so với ASTGCN \cite{astgcn}.

    \item \textbf{Tính khả thi của kiến trúc triển khai Edge--Cloud:} 
    Đề tài đã xây dựng phương án triển khai phân tán giúp tiết kiệm băng thông truyền thông tới \textbf{$62.500$ lần} (từ 437 GB/ngày xuống còn 7 MB/ngày) và bảo vệ quyền riêng tư hình ảnh theo nguyên tắc Privacy-by-Design, đồng thời ước tính tiết kiệm $89.5\%$ tổng chi phí sở hữu trong 5 năm (TCO) so với giải pháp vòng từ cảm ứng.
\end{enumerate}

\section*{2. Kiến nghị và đề xuất ứng dụng thực tiễn}
Từ các kết quả nghiên cứu và bài học thực tiễn thu được, nhóm tác giả đề xuất một số kiến nghị nhằm ứng dụng kết quả vào thực tế:

\subsection*{2.1. Kiến nghị đối với cơ quan quản lý giao thông và chính quyền đô thị}
\begin{enumerate}
    \item \textbf{Tận dụng hạ tầng camera sẵn có thay vì đầu tư mới vòng từ:} 
    Đề xuất các cơ quan quản lý giao thông tại các đô thị lớn xem xét áp dụng khung mô hình hai giai đoạn để tận dụng hệ thống camera giám sát hiện hữu, chuyển đổi thành mạng lưới cảm biến đếm xe và dự báo lưu lượng thông minh với chi phí đầu tư hạ tầng thấp hơn nhiều so với việc xẻ đường chôn vòng từ.
    \item \textbf{Thử nghiệm điều khiển đèn tín hiệu thích ứng đón đầu:} 
    Tích hợp luồng dữ liệu dự báo 15--30 phút của mô hình TA-STGCN vào các hệ thống điều khiển đèn tín hiệu tại các giao lộ trọng điểm để tự động điều chỉnh chu kỳ đèn đón đầu dòng xe, giảm thiểu nguy cơ hình thành và lan truyền sóng xung kích ùn tắc.
    \item \textbf{Áp dụng mô hình tính toán biên bảo vệ quyền riêng tư:} 
    Triển khai tính toán biên tại camera để giảm chi phí thuê băng thông truyền dẫn hàng tháng, đồng thời bảo đảm tuân thủ các quy định về bảo vệ dữ liệu cá nhân theo Nghị định 13/2023/NĐ-CP.
\end{enumerate}

\subsection*{2.2. Định hướng cho các nghiên cứu tiếp nối}
\begin{enumerate}
    \item \textbf{Mở rộng dữ liệu bằng học bán giám sát và mô hình nền tảng:} 
    Khuyến khích các nhóm nghiên cứu tiếp nối ứng dụng kỹ thuật học bán giám sát hoặc các mô hình nền tảng thị giác (Vision Foundation Models) để tự động sinh giả nhãn và mở rộng quy mô bộ dữ liệu ảnh.
    \item \textbf{Tích hợp đồ thị động thích ứng và dữ liệu đa phương thức:} 
    Kết hợp thêm các yếu tố ngoại sinh như lượng mưa thời gian thực, triều cường gây ngập cục bộ và lịch sự kiện đô thị vào đồ thị dự báo để nâng cao độ chính xác trong các tình huống thời tiết bất thường.
    \item \textbf{Đo kiểm trên phần cứng nhúng thực tế:} 
    Tiến hành tối ưu hóa mô hình sang định dạng TensorRT/INT8 và đo kiểm thực tế về tốc độ khung hình, độ trễ và công suất tiêu thụ trên các bo mạch biên chuyên dụng như NVIDIA Jetson Orin Nano, Raspberry Pi 5 và Google Coral TPU.
\end{enumerate}


% ------------------------------------------------------------------------------
% TÀI LIỆU THAM KHẢO
% ------------------------------------------------------------------------------
\newpage
\begin{thebibliography}{99}
\addcontentsline{toc}{chapter}{Tài liệu tham khảo}

\bibitem{stgcn} B. Yu, H. Yin, and Z. Zhu, ``Spatio-temporal graph convolutional networks: A deep learning framework for traffic forecasting,'' in \textit{Proc. Int. Joint Conf. Artif. Intell. (IJCAI)}, 2018, pp. 3634--3640.

\bibitem{wavenet} Z. Wu, S. Pan, G. Long, J. Jiang, and C. Zhang, ``Graph WaveNet for deep spatial-temporal graph modeling,'' in \textit{Proc. Int. Joint Conf. Artif. Intell. (IJCAI)}, 2019, pp. 1907--1913.

\bibitem{astgcn} S. Guo, Y. Lin, N. Feng, C. Song, and H. Wan, ``Attention based spatial-temporal graph convolutional networks for traffic flow forecasting,'' in \textit{Proc. AAAI Conf. Artif. Intell. (AAAI)}, vol. 33, no. 1, 2019, pp. 922--929.

\bibitem{staeformer} H. Liu, Z. Dong, R. Jiang, J. Deng, J. Deng, Q. Chen, and X. Song, ``Spatio-temporal adaptive embedding makes vanilla Transformer state-of-the-art for traffic flow forecasting,'' in \textit{Proc. ACM Int. Conf. Inf. Knowl. Manage. (CIKM)}, 2023, pp. 4125--4129.

\bibitem{itransformer} Y. Liu, T. Hu, H. Zhang, H. Wu, S. Wang, L. Ma, and M. Long, ``iTransformer: Inverted Transformers are effective for time series forecasting,'' in \textit{Proc. Int. Conf. Learn. Represent. (ICLR)}, 2024.

\bibitem{megacrn} R. Jiang, D. Yin, Z. Wang, Y. Wang, J. Deng, H. Liu, Z. Cai, J. Deng, and X. Song, ``Spatio-temporal meta-graph learning for traffic flow forecasting,'' in \textit{Proc. AAAI Conf. Artif. Intell. (AAAI)}, 2023, pp. 8078--8086.

\bibitem{dstagnn} S. Lan, Y. Ma, W. Huang, W. Wang, H. Yang, and P. Li, ``DSTAGNN: Dynamic spatial-temporal aware graph neural network for traffic flow forecasting,'' in \textit{Proc. Int. Conf. Mach. Learn. (ICML)}, 2022, pp. 11906--11917.

\bibitem{dcrnn} Y. Li, R. Yu, C. Shahabi, and Y. Liu, ``Diffusion convolutional recurrent neural network: Data-driven traffic forecasting,'' in \textit{Proc. Int. Conf. Learn. Represent. (ICLR)}, 2018.

\bibitem{efficientnet} M. Tan and Q. V. Le, ``EfficientNet: Rethinking model scaling for convolutional neural networks,'' in \textit{Proc. Int. Conf. Mach. Learn. (ICML)}, 2019, pp. 6105--6114.

\bibitem{resnet} K. He, X. Zhang, S. Ren, and J. Sun, ``Deep residual learning for image recognition,'' in \textit{Proc. IEEE/CVF Conf. Comput. Vis. Pattern Recognit. (CVPR)}, 2016, pp. 770--778.

\bibitem{mobilenet} A. Howard et al., ``Searching for MobileNetV3,'' in \textit{Proc. IEEE/CVF Int. Conf. Comput. Vis. (ICCV)}, 2019, pp. 1314--1324.

\bibitem{convnext} Z. Liu, H. Mao, C. Y. Wu, C. Feichtenhofer, T. Darrell, and S. Xie, ``A convnet for the 2020s,'' in \textit{Proc. IEEE/CVF Conf. Comput. Vis. Pattern Recognit. (CVPR)}, 2022, pp. 11976--11986.

\bibitem{vit} A. Dosovitskiy et al., ``An image is worth 16x16 words: Transformers for image recognition at scale,'' in \textit{Proc. Int. Conf. Learn. Represent. (ICLR)}, 2021.

\bibitem{attention} A. Vaswani et al., ``Attention is all you need,'' in \textit{Proc. Adv. Neural Inf. Process. Syst. (NeurIPS)}, 2017, pp. 5998--6008.

\bibitem{chebnet} M. Defferrard, X. Bresson, and P. Vandergheynst, ``Convolutional neural networks on graphs with fast localized spectral filtering,'' in \textit{Proc. Adv. Neural Inf. Process. Syst. (NeurIPS)}, 2016, pp. 3828--3836.

\bibitem{gman} C. Zheng, X. Fan, C. Wang, and J. Qi, ``GMAN: A graph multi-attention network for traffic prediction,'' in \textit{Proc. AAAI Conf. Artif. Intell. (AAAI)}, 2020, pp. 1234--1241.

\bibitem{stid} Z. Shao, Z. Zhang, F. Wang, and Y. Xu, ``Spatial-temporal identity: A simple yet effective baseline for spatio-temporal traffic forecasting,'' in \textit{Proc. ACM Int. Conf. Inf. Knowl. Manage. (CIKM)}, 2022, pp. 4485--4489.

\bibitem{tgcn} L. Zhao, Y. Song, C. Zhang, Y. Liu, P. Wang, T. Lin, M. Deng, and H. Li, ``T-GCN: A temporal graph convolutional network for traffic prediction,'' \textit{IEEE Trans. Intell. Transp. Syst.}, vol. 21, no. 9, pp. 3848--3858, 2020.

\bibitem{yolo} J. Redmon, S. Divvala, R. Girshick, and A. Farhadi, ``You only look once: Unified, real-time object detection,'' in \textit{Proc. IEEE Conf. Comput. Vis. Pattern Recognit. (CVPR)}, 2016, pp. 779--788.

\bibitem{csrnet} Y. Li, X. Zhang, and D. Chen, ``CSRNet: Dilated convolutional neural networks for understanding the highly congested scenes,'' in \textit{Proc. IEEE/CVF Conf. Comput. Vis. Pattern Recognit. (CVPR)}, 2018, pp. 1091--1100.

\bibitem{mcnn} Y. Zhang, D. Zhou, S. Chen, S. Gao, and Y. Ma, ``Single-image crowd counting via multi-column convolutional neural network,'' in \textit{Proc. IEEE Conf. Comput. Vis. Pattern Recognit. (CVPR)}, 2016, pp. 589--597.

\bibitem{can} S. Bai, Z. He, Y. Wu, and D. Wang, ``Context-aware crowd counting,'' in \textit{Proc. IEEE/CVF Conf. Comput. Vis. Pattern Recognit. (CVPR)}, 2019, pp. 6909--6918.

\bibitem{huber} P. J. Huber, ``Robust estimation of a location parameter,'' \textit{Ann. Math. Statist.}, vol. 35, no. 1, pp. 73--101, 1964.

\bibitem{signal_control} K.-L. A. Yau, J. Qadir, H. L. Khoo, M. H. Ling, and P. Komisarczuk, ``A survey on reinforcement learning models and algorithms for traffic signal control,'' \textit{ACM Comput. Surv.}, vol. 50, no. 3, pp. 1--38, 2017.

\bibitem{ctm} C. F. Daganzo, ``The cell transmission model: A dynamic representation of highway traffic conforming to the hydrodynamic theory,'' \textit{Transp. Res. Part B: Methodol.}, vol. 28, no. 4, pp. 269--287, 1994.

\bibitem{cao2007} N. Y. Cao and M. Sano, ``A study on motorcycle lane capacity in urban areas of developing countries,'' \textit{J. East. Asia Soc. Transp. Stud.}, vol. 7, pp. 2320--2333, 2007.

\bibitem{lee2009} T.-C. Lee, ``An agent-based model to simulate motorcycle behaviour in mixed traffic flow,'' in \textit{Proc. Int. Conf. Comput. Urban Plann. Urban Manage. (CUPUM)}, 2009.

\bibitem{pdformer} J. Jiang, C. Han, W. X. Zhao, and J. Wang, ``PDFormer: Propagation delay-aware dynamic long-range transformer for traffic flow forecasting,'' in \textit{Proc. AAAI Conf. Artif. Intell. (AAAI)}, 2023, pp. 4365--4373.

\bibitem{stgncde} J. Choi, H. Choi, J. Hwang, and N. Park, ``Graph neural controlled differential equations for traffic forecasting,'' in \textit{Proc. AAAI Conf. Artif. Intell. (AAAI)}, 2022, pp. 6367--6374.

\bibitem{mamba} A. Gu and T. Dao, ``Mamba: Linear-time sequence modeling with selective state spaces,'' \textit{arXiv preprint arXiv:2312.00752}, 2023.

\bibitem{patchtst} Y. Nie, N. H. Nguyen, P. Sinthong, and J. Kalagnanam, ``A time series is worth 64 words: Long-term forecasting with transformers,'' in \textit{Proc. Int. Conf. Learn. Represent. (ICLR)}, 2023.

\bibitem{adam} D. P. Kingma and J. Ba, ``Adam: A method for stochastic optimization,'' in \textit{Proc. Int. Conf. Learn. Represent. (ICLR)}, 2015.

\end{thebibliography}

\end{document}
